{.977}}
\put(22.86,90.75){\line(-1,0){.977}}
\put(20.91,90.93){\line(-1,0){.977}}
%\end
%\dashline{1}(24.75,103.59)(26.25,107.59)
\multiput(24.68,103.52)(.03125,.08333){8}{\line(0,1){.08333}}
\multiput(25.18,104.86)(.03125,.08333){8}{\line(0,1){.08333}}
\multiput(25.68,106.19)(.03125,.08333){8}{\line(0,1){.08333}}
%\end
%\dashline{1}(25.75,106.34)(27.25,110.34)
\multiput(25.68,106.27)(.03125,.08333){8}{\line(0,1){.08333}}
\multiput(26.18,107.61)(.03125,.08333){8}{\line(0,1){.08333}}
\multiput(26.68,108.94)(.03125,.08333){8}{\line(0,1){.08333}}
%\end
\put(69,104.34){\line(1,0){23.5}}
%\emline(92.5,104.34)(82.75,79.34)
\multiput(92.5,104.34)(-.033737024,-.08650519){289}{\line(0,-1){.08650519}}
%\end
%\emline(71.25,82.34)(77.5,95.84)
\multiput(71.25,82.34)(.03360215,.07258065){186}{\line(0,1){.07258065}}
%\end
\put(77.5,95.84){\line(-1,0){9}}
%\emline(77.25,95.34)(79.5,100.59)
\multiput(77.25,95.34)(.0335821,.0783582){67}{\line(0,1){.0783582}}
%\end
%\emline(79.5,100.59)(92.25,104.09)
\multiput(79.5,100.59)(.12259615,.03365385){104}{\line(1,0){.12259615}}
%\end
%\emline(79.75,100.59)(88.75,95.59)
\multiput(79.75,100.59)(.06040268,-.03355705){149}{\line(1,0){.06040268}}
%\end
%\dashline{1}(79.5,100.84)(81,105.09)
\multiput(79.43,100.77)(.03125,.08854){8}{\line(0,1){.08854}}
\multiput(79.93,102.19)(.03125,.08854){8}{\line(0,1){.08854}}
\multiput(80.43,103.61)(.03125,.08854){8}{\line(0,1){.08854}}
%\end
%\dashline{1}(77.25,95.59)(92.25,103.84)
\multiput(77.18,95.52)(.060729,.033401){13}{\line(1,0){.060729}}
\multiput(78.76,96.39)(.060729,.033401){13}{\line(1,0){.060729}}
\multiput(80.34,97.26)(.060729,.033401){13}{\line(1,0){.060729}}
\multiput(81.92,98.13)(.060729,.033401){13}{\line(1,0){.060729}}
\multiput(83.5,99)(.060729,.033401){13}{\line(1,0){.060729}}
\multiput(85.07,99.87)(.060729,.033401){13}{\line(1,0){.060729}}
\multiput(86.65,100.74)(.060729,.033401){13}{\line(1,0){.060729}}
\multiput(88.23,101.6)(.060729,.033401){13}{\line(1,0){.060729}}
\multiput(89.81,102.47)(.060729,.033401){13}{\line(1,0){.060729}}
\multiput(91.39,103.34)(.060729,.033401){13}{\line(1,0){.060729}}
%\end
%\dashline{1}(73,99.09)(86.75,89.59)
\multiput(72.93,99.02)(.047743,-.032986){16}{\line(1,0){.047743}}
\multiput(74.46,97.97)(.047743,-.032986){16}{\line(1,0){.047743}}
\multiput(75.99,96.91)(.047743,-.032986){16}{\line(1,0){.047743}}
\multiput(77.51,95.86)(.047743,-.032986){16}{\line(1,0){.047743}}
\multiput(79.04,94.8)(.047743,-.032986){16}{\line(1,0){.047743}}
\multiput(80.57,93.75)(.047743,-.032986){16}{\line(1,0){.047743}}
\multiput(82.1,92.69)(.047743,-.032986){16}{\line(1,0){.047743}}
\multiput(83.62,91.64)(.047743,-.032986){16}{\line(1,0){.047743}}
\multiput(85.15,90.58)(.047743,-.032986){16}{\line(1,0){.047743}}
%\end
%\dashline{1}(77.5,95.59)(84,82.34)
\multiput(77.43,95.52)(.03125,-.063702){13}{\line(0,-1){.063702}}
\multiput(78.24,93.87)(.03125,-.063702){13}{\line(0,-1){.063702}}
\multiput(79.05,92.21)(.03125,-.063702){13}{\line(0,-1){.063702}}
\multiput(79.87,90.56)(.03125,-.063702){13}{\line(0,-1){.063702}}
\multiput(80.68,88.9)(.03125,-.063702){13}{\line(0,-1){.063702}}
\multiput(81.49,87.24)(.03125,-.063702){13}{\line(0,-1){.063702}}
\multiput(82.3,85.59)(.03125,-.063702){13}{\line(0,-1){.063702}}
\multiput(83.12,83.93)(.03125,-.063702){13}{\line(0,-1){.063702}}
%\end
%\dashline{1}(84,82.34)(71.5,82.34)
\put(83.93,82.27){\line(-1,0){.962}}
\put(82.01,82.27){\line(-1,0){.962}}
\put(80.08,82.27){\line(-1,0){.962}}
\put(78.16,82.27){\line(-1,0){.962}}
\put(76.24,82.27){\line(-1,0){.962}}
\put(74.31,82.27){\line(-1,0){.962}}
\put(72.39,82.27){\line(-1,0){.962}}
%\end
\thicklines
\put(28.68,64.09){\line(0,-1){10.5}}
\put(56.18,20.59){\line(5,-1){13.75}}
\put(69.93,17.84){\line(0,1){10.25}}
%\emline(55.93,20.84)(51.68,21.84)
\multiput(55.93,20.84)(-.141667,.033333){30}{\line(-1,0){.141667}}
%\end
\thinlines
%\emline(69.93,17.84)(65.68,26.59)
\multiput(69.93,17.84)(-.03373016,.06944444){126}{\line(0,1){.06944444}}
%\end
%\emline(48.63,11.99)(51.63,22.24)
\multiput(48.63,11.99)(.0337079,.1151685){89}{\line(0,1){.1151685}}
%\end
%\emline(32.28,55.09)(54.78,47.84)
\multiput(32.28,55.09)(.10465116,-.03372093){215}{\line(1,0){.10465116}}
%\end
%\emline(71.43,37.59)(48.68,30.84)
\multiput(71.43,37.59)(-.11318408,-.03358209){201}{\line(-1,0){.11318408}}
%\end
%\dashline{1}(55.43,60.34)(32.43,55.09)
\multiput(55.36,60.27)(-.13143,-.03){7}{\line(-1,0){.13143}}
\multiput(53.52,59.85)(-.13143,-.03){7}{\line(-1,0){.13143}}
\multiput(51.68,59.43)(-.13143,-.03){7}{\line(-1,0){.13143}}
\multiput(49.84,59.01)(-.13143,-.03){7}{\line(-1,0){.13143}}
\multiput(48,58.59)(-.13143,-.03){7}{\line(-1,0){.13143}}
\multiput(46.16,58.17)(-.13143,-.03){7}{\line(-1,0){.13143}}
\multiput(44.32,57.75)(-.13143,-.03){7}{\line(-1,0){.13143}}
\multiput(42.48,57.33)(-.13143,-.03){7}{\line(-1,0){.13143}}
\multiput(40.64,56.91)(-.13143,-.03){7}{\line(-1,0){.13143}}
\multiput(38.8,56.49)(-.13143,-.03){7}{\line(-1,0){.13143}}
\multiput(36.96,56.07)(-.13143,-.03){7}{\line(-1,0){.13143}}
\multiput(35.12,55.65)(-.13143,-.03){7}{\line(-1,0){.13143}}
\multiput(33.28,55.23)(-.13143,-.03){7}{\line(-1,0){.13143}}
%\end
%\dashline{1}(40.43,60.34)(25.68,43.84)
\multiput(40.36,60.27)(-.032065,-.03587){20}{\line(0,-1){.03587}}
\multiput(39.08,58.84)(-.032065,-.03587){20}{\line(0,-1){.03587}}
\multiput(37.8,57.4)(-.032065,-.03587){20}{\line(0,-1){.03587}}
\multiput(36.51,55.97)(-.032065,-.03587){20}{\line(0,-1){.03587}}
\multiput(35.23,54.54)(-.032065,-.03587){20}{\line(0,-1){.03587}}
\multiput(33.95,53.1)(-.032065,-.03587){20}{\line(0,-1){.03587}}
\multiput(32.67,51.67)(-.032065,-.03587){20}{\line(0,-1){.03587}}
\multiput(31.38,50.23)(-.032065,-.03587){20}{\line(0,-1){.03587}}
\multiput(30.1,48.8)(-.032065,-.03587){20}{\line(0,-1){.03587}}
\multiput(28.82,47.36)(-.032065,-.03587){20}{\line(0,-1){.03587}}
\multiput(27.53,45.93)(-.032065,-.03587){20}{\line(0,-1){.03587}}
\multiput(26.25,44.49)(-.032065,-.03587){20}{\line(0,-1){.03587}}
%\end
%\dashline{1}(23.43,40.84)(38.18,24.34)
\multiput(23.36,40.77)(.032065,-.03587){20}{\line(0,-1){.03587}}
\multiput(24.64,39.34)(.032065,-.03587){20}{\line(0,-1){.03587}}
\multiput(25.93,37.9)(.032065,-.03587){20}{\line(0,-1){.03587}}
\multiput(27.21,36.47)(.032065,-.03587){20}{\line(0,-1){.03587}}
\multiput(28.49,35.04)(.032065,-.03587){20}{\line(0,-1){.03587}}
\multiput(29.77,33.6)(.032065,-.03587){20}{\line(0,-1){.03587}}
\multiput(31.06,32.17)(.032065,-.03587){20}{\line(0,-1){.03587}}
\multiput(32.34,30.73)(.032065,-.03587){20}{\line(0,-1){.03587}}
\multiput(33.62,29.3)(.032065,-.03587){20}{\line(0,-1){.03587}}
\multiput(34.9,27.86)(.032065,-.03587){20}{\line(0,-1){.03587}}
\multiput(36.19,26.43)(.032065,-.03587){20}{\line(0,-1){.03587}}
\multiput(37.47,24.99)(.032065,-.03587){20}{\line(0,-1){.03587}}
%\end
%\dashline{1}(29.18,28.34)(52.43,21.59)
\multiput(29.11,28.27)(.10333,-.03){9}{\line(1,0){.10333}}
\multiput(30.97,27.73)(.10333,-.03){9}{\line(1,0){.10333}}
\multiput(32.83,27.19)(.10333,-.03){9}{\line(1,0){.10333}}
\multiput(34.69,26.65)(.10333,-.03){9}{\line(1,0){.10333}}
\multiput(36.55,26.11)(.10333,-.03){9}{\line(1,0){.10333}}
\multiput(38.41,25.57)(.10333,-.03){9}{\line(1,0){.10333}}
\multiput(40.27,25.03)(.10333,-.03){9}{\line(1,0){.10333}}
\multiput(42.13,24.49)(.10333,-.03){9}{\line(1,0){.10333}}
\multiput(43.99,23.95)(.10333,-.03){9}{\line(1,0){.10333}}
\multiput(45.85,23.41)(.10333,-.03){9}{\line(1,0){.10333}}
\multiput(47.71,22.87)(.10333,-.03){9}{\line(1,0){.10333}}
\multiput(49.57,22.33)(.10333,-.03){9}{\line(1,0){.10333}}
\multiput(51.43,21.79)(.10333,-.03){9}{\line(1,0){.10333}}
%\end
%\dashline{1}(56.43,20.59)(71.43,37.59)
\multiput(56.36,20.52)(.032895,.037281){19}{\line(0,1){.037281}}
\multiput(57.61,21.94)(.032895,.037281){19}{\line(0,1){.037281}}
\multiput(58.86,23.36)(.032895,.037281){19}{\line(0,1){.037281}}
\multiput(60.11,24.77)(.032895,.037281){19}{\line(0,1){.037281}}
\multiput(61.36,26.19)(.032895,.037281){19}{\line(0,1){.037281}}
\multiput(62.61,27.61)(.032895,.037281){19}{\line(0,1){.037281}}
\multiput(63.86,29.02)(.032895,.037281){19}{\line(0,1){.037281}}
\multiput(65.11,30.44)(.032895,.037281){19}{\line(0,1){.037281}}
\multiput(66.36,31.86)(.032895,.037281){19}{\line(0,1){.037281}}
\multiput(67.61,33.27)(.032895,.037281){19}{\line(0,1){.037281}}
\multiput(68.86,34.69)(.032895,.037281){19}{\line(0,1){.037281}}
\multiput(70.11,36.11)(.032895,.037281){19}{\line(0,1){.037281}}
%\end
%\dashline{1}(68.93,52.84)(44.43,59.59)
\multiput(68.86,52.77)(-.11779,.03245){8}{\line(-1,0){.11779}}
\multiput(66.98,53.29)(-.11779,.03245){8}{\line(-1,0){.11779}}
\multiput(65.09,53.81)(-.11779,.03245){8}{\line(-1,0){.11779}}
\multiput(63.21,54.33)(-.11779,.03245){8}{\line(-1,0){.11779}}
\multiput(61.32,54.85)(-.11779,.03245){8}{\line(-1,0){.11779}}
\multiput(59.44,55.37)(-.11779,.03245){8}{\line(-1,0){.11779}}
\multiput(57.55,55.89)(-.11779,.03245){8}{\line(-1,0){.11779}}
\multiput(55.67,56.41)(-.11779,.03245){8}{\line(-1,0){.11779}}
\multiput(53.78,56.93)(-.11779,.03245){8}{\line(-1,0){.11779}}
\multiput(51.9,57.45)(-.11779,.03245){8}{\line(-1,0){.11779}}
\multiput(50.01,57.97)(-.11779,.03245){8}{\line(-1,0){.11779}}
\multiput(48.13,58.49)(-.11779,.03245){8}{\line(-1,0){.11779}}
\multiput(46.25,59.01)(-.11779,.03245){8}{\line(-1,0){.11779}}
%\end
%\qbezvec[both](25,128.25)(27.13,129.13)(26.75,131.5)
\put(26.75,131.5){\vector(0,1){.07}}\put(25,128.25){\vector(-3,-2){.07}}\qbezier(25,128.25)(27.13,129.13)(26.75,131.5)
%\end
\put(28.75,16.75){\line(0,1){15.75}}
%\dashline{1}(28.75,53.5)(28.75,32.75)
\put(28.68,53.43){\line(0,-1){.988}}
\put(28.68,51.45){\line(0,-1){.988}}
\put(28.68,49.48){\line(0,-1){.988}}
\put(28.68,47.5){\line(0,-1){.988}}
\put(28.68,45.52){\line(0,-1){.988}}
\put(28.68,43.55){\line(0,-1){.988}}
\put(28.68,41.57){\line(0,-1){.988}}
\put(28.68,39.6){\line(0,-1){.988}}
\put(28.68,37.62){\line(0,-1){.988}}
\put(28.68,35.64){\line(0,-1){.988}}
\put(28.68,33.67){\line(0,-1){.988}}
%\end
%\dashline{1}(28,32.75)(28,32.75)
%\end
%\dashline{1}(69.75,28)(69.75,48.75)
\put(69.68,27.93){\line(0,1){.988}}
\put(69.68,29.91){\line(0,1){.988}}
\put(69.68,31.88){\line(0,1){.988}}
\put(69.68,33.86){\line(0,1){.988}}
\put(69.68,35.83){\line(0,1){.988}}
\put(69.68,37.81){\line(0,1){.988}}
\put(69.68,39.79){\line(0,1){.988}}
\put(69.68,41.76){\line(0,1){.988}}
\put(69.68,43.74){\line(0,1){.988}}
\put(69.68,45.72){\line(0,1){.988}}
\put(69.68,47.69){\line(0,1){.988}}
%\end
%\dashline{1}(69.75,48.75)(69.75,48.75)
%\end
%\dashline{1}(74.25,40.25)(59.75,58.25)
\multiput(74.18,40.18)(-.032222,.04){18}{\line(0,1){.04}}
\multiput(73.02,41.62)(-.032222,.04){18}{\line(0,1){.04}}
\multiput(71.86,43.06)(-.032222,.04){18}{\line(0,1){.04}}
\multiput(70.7,44.5)(-.032222,.04){18}{\line(0,1){.04}}
\multiput(69.54,45.94)(-.032222,.04){18}{\line(0,1){.04}}
\multiput(68.38,47.38)(-.032222,.04){18}{\line(0,1){.04}}
\multiput(67.22,48.82)(-.032222,.04){18}{\line(0,1){.04}}
\multiput(66.06,50.26)(-.032222,.04){18}{\line(0,1){.04}}
\multiput(64.9,51.7)(-.032222,.04){18}{\line(0,1){.04}}
\multiput(63.74,53.14)(-.032222,.04){18}{\line(0,1){.04}}
\multiput(62.58,54.58)(-.032222,.04){18}{\line(0,1){.04}}
\multiput(61.42,56.02)(-.032222,.04){18}{\line(0,1){.04}}
\multiput(60.26,57.46)(-.032222,.04){18}{\line(0,1){.04}}
%\end
\put(69.75,48.5){\line(0,1){16.25}}
%\emline(56,60.25)(69.75,64.5)
\multiput(56,60.25)(.10912698,.03373016){126}{\line(1,0){.10912698}}
%\end
\put(38.25,42.5){\line(0,-1){18.5}}
%\emline(38,24)(28.75,26.25)
\multiput(38,24)(-.1380597,.0335821){67}{\line(-1,0){.1380597}}
%\end
%\emline(25.75,44)(15.75,33)
\multiput(25.75,44)(-.033670034,-.037037037){297}{\line(0,-1){.037037037}}
%\end
%\emline(56.25,20.5)(48.75,11.75)
\multiput(56.25,20.5)(-.03363229,-.03923767){223}{\line(0,-1){.03923767}}
%\end
%\emline(60,58.25)(49.25,70.75)
\multiput(60,58.25)(-.03369906,.039184953){319}{\line(0,1){.039184953}}
%\end
%\emline(41,60.75)(49.25,70.75)
\multiput(41,60.75)(.033673469,.040816327){245}{\line(0,1){.040816327}}
%\end
%\emline(49.25,70.75)(44.5,59.25)
\multiput(49.25,70.75)(-.03368794,-.08156028){141}{\line(0,-1){.08156028}}
%\end
%\emline(44.71,59.7)(56.27,62.65)
\multiput(44.71,59.7)(.1313901,.0334447){88}{\line(1,0){.1313901}}
%\end
\put(59.46,58.8){\line(0,-1){20.81}}
%\emline(54.8,32.58)(69.73,49.82)
\multiput(54.8,32.58)(.033692802,.038912814){443}{\line(0,1){.038912814}}
%\end
%\emline(69.73,52.56)(81.92,49.61)
\multiput(69.73,52.56)(.1385568,-.0334447){88}{\line(1,0){.1385568}}
%\end
\put(81.92,49.61){\line(-1,0){12.19}}
%\emline(82.13,49.4)(71.62,37.63)
\multiput(82.13,49.4)(-.03368976,-.037732531){312}{\line(0,-1){.037732531}}
%\end
%\emline(69.73,49.61)(75.82,41.83)
\multiput(69.73,49.61)(.03368231,-.04297399){181}{\line(0,-1){.04297399}}
%\end
%\emline(59.65,58.69)(69.73,64.35)
\multiput(59.65,58.69)(.05997781,.03367175){168}{\line(1,0){.05997781}}
%\end
%\emline(59.65,58.69)(69.73,55.68)
\multiput(59.65,58.69)(.1119586,-.0333912){90}{\line(1,0){.1119586}}
%\end
%\emline(44.62,59.4)(59.3,42.25)
\multiput(44.62,59.4)(.033729806,-.039419171){435}{\line(0,-1){.039419171}}
%\end
%\emline(15.81,32.88)(28.71,28.46)
\multiput(15.81,32.88)(.09850915,-.03373601){131}{\line(1,0){.09850915}}
%\end
\put(15.63,32.7){\line(1,0){12.9}}
%\emline(28.54,32.7)(22,39.77)
\multiput(28.54,32.7)(-.03371514,.0364488){194}{\line(0,1){.0364488}}
%\end
%\emline(51.87,21.92)(38.26,37.3)
\multiput(51.87,21.92)(-.033692588,.038068249){404}{\line(0,1){.038068249}}
%\end
%\emline(32.25,54.98)(15.45,50.91)
\multiput(32.25,54.98)(-.13879162,-.03360218){121}{\line(-1,0){.13879162}}
%\end
%\emline(15.45,50.73)(23.23,41.19)
\multiput(15.45,50.73)(.03367175,-.04132442){231}{\line(0,-1){.04132442}}
%\end
%\emline(15.63,50.73)(25.88,44.02)
\multiput(15.63,50.73)(.05126524,-.03358757){200}{\line(1,0){.05126524}}
%\end
%\emline(25.88,44.02)(49.04,49.67)
\multiput(25.88,44.02)(.13784373,.03367175){168}{\line(1,0){.13784373}}
%\end
%\emline(28.71,32.53)(43.56,48.26)
\multiput(28.71,32.53)(.033671751,.035676022){441}{\line(0,1){.035676022}}
%\end
%\emline(32.6,54.98)(38.08,61.34)
\multiput(32.6,54.98)(.03362011,.03904271){163}{\line(0,1){.03904271}}
%\end
%\emline(32.45,54.95)(28.56,64.14)
\multiput(32.45,54.95)(-.03352661,.07924473){116}{\line(0,1){.07924473}}
%\end
%\emline(28.71,64.17)(44.62,59.57)
\multiput(28.71,64.17)(.11613068,-.03354886){137}{\line(1,0){.11613068}}
%\end
\put(25.88,44.19){\line(0,1){9.19}}
%\emline(38.41,24.01)(48.84,11.82)
\multiput(38.41,24.01)(.033644597,-.039347071){310}{\line(0,-1){.039347071}}
%\end
%\emline(65.36,26.49)(59.53,19.95)
\multiput(65.36,26.49)(-.03372041,-.03780773){173}{\line(0,-1){.03780773}}
%\end
%\dashline{1}(40.73,20.86)(65.48,26.52)
\multiput(40.66,20.79)(.13598,.03108){7}{\line(1,0){.13598}}
\multiput(42.57,21.22)(.13598,.03108){7}{\line(1,0){.13598}}
\multiput(44.47,21.66)(.13598,.03108){7}{\line(1,0){.13598}}
\multiput(46.37,22.09)(.13598,.03108){7}{\line(1,0){.13598}}
\multiput(48.28,22.53)(.13598,.03108){7}{\line(1,0){.13598}}
\multiput(50.18,22.97)(.13598,.03108){7}{\line(1,0){.13598}}
\multiput(52.09,23.4)(.13598,.03108){7}{\line(1,0){.13598}}
\multiput(53.99,23.84)(.13598,.03108){7}{\line(1,0){.13598}}
\multiput(55.89,24.27)(.13598,.03108){7}{\line(1,0){.13598}}
\multiput(57.8,24.71)(.13598,.03108){7}{\line(1,0){.13598}}
\multiput(59.7,25.14)(.13598,.03108){7}{\line(1,0){.13598}}
\multiput(61.6,25.58)(.13598,.03108){7}{\line(1,0){.13598}}
\multiput(63.51,26.01)(.13598,.03108){7}{\line(1,0){.13598}}
%\end
%\dashline{1}(65.48,26.52)(65.48,26.52)
%\end
%\emline(28.71,16.62)(38.26,23.86)
\multiput(28.71,16.62)(.04439973,.0337109){215}{\line(1,0){.04439973}}
%\end
%\emline(28.71,16.79)(40.91,20.68)
\multiput(28.71,16.79)(.10515165,.03352661){116}{\line(1,0){.10515165}}
%\end
\put(51.87,21.74){\line(-3,-1){9.02}}
%\emline(65.48,26.52)(81.75,31.11)
\multiput(65.48,26.52)(.11871136,.03354886){137}{\line(1,0){.11871136}}
%\end
%\emline(81.75,31.11)(74.14,40.31)
\multiput(81.75,31.11)(-.0336345,.04067428){226}{\line(0,1){.04067428}}
%\end
%\emline(71.57,37.61)(81.83,30.92)
\multiput(71.57,37.61)(.05154227,-.03361452){199}{\line(1,0){.05154227}}
%\end
\put(71.72,37.61){\line(0,-1){9.22}}
%\emline(65.62,26.46)(42.43,32.55)
\multiput(65.62,26.46)(-.128119,.0336723){181}{\line(-1,0){.128119}}
%\end
\put(22.7,106.5){\makebox(0,0)[cc]{$A_{4k}$}}
\put(94,105.59){\makebox(0,0)[cc]{$A_{4k}$}}
\put(21.5,134.34){\makebox(0,0)[cc]{$A_{4k}$}}
\put(36,89.84){\makebox(0,0)[cc]{$A_{4k+2}$}}
\put(23.25,119.84){\makebox(0,0)[cc]{$A_{4k+1}$}}
\put(39,119.84){\makebox(0,0)[cc]{$A_{4k+2}$}}
\put(84.25,94.84){\makebox(0,0)[cc]{$A_{4k+2}$}}
\put(36,96.09){\makebox(0,0)[cc]{$B$}}
\put(77,101.34){\makebox(0,0)[cc]{$B$}}
\put(32.5,120.84){\makebox(0,0)[cc]{$B$}}
\put(31.25,129.09){\makebox(0,0)[cc]{$\frac{2\pi}{10}$}}
\put(57.5,129.09){\makebox(0,0)[cc]{$\frac{\pi}{5}$}}
\put(54.5,126.34){\makebox(0,0)[cc]{$\frac{\pi}{5}$}}
\put(50.75,112.34){\makebox(0,0)[cc]{(b)}}
\put(15.75,75.84){\makebox(0,0)[cc]{(c)}}
\put(77.25,75.59){\makebox(0,0)[cc]{(d)}}
\put(34.75,114.84){\makebox(0,0)[cc]{Fold Line}}
\put(33.5,107){\makebox(0,0)[cc]{Fold Line}}
\put(48.43,3.34){\makebox(0,0)[cc]{Figure 5.3}}
\end{picture}
\end{center}

\begin{center}
%TeXCAD Picture [fig5.4.pic]. Options:
%\grade{\on}
%\emlines{\off}
%\epic{\off}
%\beziermacro{\on}
%\reduce{\on}
%\snapping{\off}
%\quality{8.00}
%\graddiff{0.01}
%\snapasp{1}
%\zoom{4.0000}
\unitlength 1mm % = 2.85pt
\linethickness{0.4pt}
\ifx\plotpoint\undefined\newsavebox{\plotpoint}\fi % GNUPLOT compatibility
\begin{picture}(62.25,87)(0,0)
\put(6,87){\line(1,0){52.25}}
\put(58.25,87){\line(0,1){0}}
\put(57.75,81.75){\line(-1,0){51.5}}
%\dottedline(8.75,81.75)(6,85.5)
\multiput(8.68,81.68)(-.55,.75){6}{{\rule{.4pt}{.4pt}}}
%\end
%\dottedline(6,85.5)(15.5,81.75)
\multiput(5.93,85.43)(.8636,-.3409){12}{{\rule{.4pt}{.4pt}}}
%\end
%\dottedline(15.5,81.75)(18,86.75)
\multiput(15.43,81.68)(.3571,.7143){8}{{\rule{.4pt}{.4pt}}}
%\end
%\dottedline(18,86.75)(22.75,82)
\multiput(17.93,86.68)(.5938,-.5938){9}{{\rule{.4pt}{.4pt}}}
%\end
%\dottedline(22.75,82)(31.75,86.75)
\multiput(22.68,81.93)(.8182,.4318){12}{{\rule{.4pt}{.4pt}}}
%\end
%\dottedline(31.75,86.75)(35.75,81.75)
\multiput(31.68,86.68)(.5,-.625){9}{{\rule{.4pt}{.4pt}}}
%\end
%\dottedline(35.75,81.75)(40,86.75)
\multiput(35.68,81.68)(.5313,.625){9}{{\rule{.4pt}{.4pt}}}
%\end
%\dottedline(40,86.75)(43.75,82)
\multiput(39.93,86.68)(.5357,-.6786){8}{{\rule{.4pt}{.4pt}}}
%\end
%\dottedline(40.5,86.75)(51.25,82)
\multiput(40.43,86.68)(.8958,-.3958){13}{{\rule{.4pt}{.4pt}}}
%\end
%\dottedline(51.25,82)(52.25,86.5)
\multiput(51.18,81.93)(.1667,.75){7}{{\rule{.4pt}{.4pt}}}
%\end
%\dottedline(52.25,86.5)(58.25,81.75)
\multiput(52.18,86.43)(.6667,-.5278){10}{{\rule{.4pt}{.4pt}}}
%\end
%\dottedline(23,82)(25.5,87)
\multiput(22.93,81.93)(.3571,.7143){8}{{\rule{.4pt}{.4pt}}}
%\end
\thicklines
%\emline(42.25,11.25)(43,22.75)
\multiput(42.25,11.25)(.032609,.5){23}{\line(0,1){.5}}
%\end
%\emline(61.5,52.25)(50,53)
\multiput(61.5,52.25)(-.5,.032609){23}{\line(-1,0){.5}}
%\end
\thinlines
%\emline(29.25,34.75)(49.5,26.75)
\multiput(29.25,34.75)(.08508403,-.03361345){238}{\line(1,0){.08508403}}
%\end
%\emline(38,39.25)(46,59.5)
\multiput(38,39.25)(.03361345,.08508403){238}{\line(0,1){.08508403}}
%\end
%\emline(29.5,34.5)(28.25,35)
\multiput(29.5,34.5)(-.083333,.033333){15}{\line(-1,0){.083333}}
%\end
%\emline(38.25,39.5)(37.75,38.25)
\multiput(38.25,39.5)(-.033333,-.083333){15}{\line(0,-1){.083333}}
%\end
\put(19,78){\makebox(0,0)[cc]{$D^{1} U^{1} D^{2}$ tape}}
%\emline(38,25)(52.5,19)
\multiput(38,25)(.08146067,-.03370787){178}{\line(1,0){.08146067}}
%\end
%\emline(47.75,48)(53.75,62.5)
\multiput(47.75,48)(.03370787,.08146067){178}{\line(0,1){.08146067}}
%\end
%\emline(52.5,19)(48,29.5)
\multiput(52.5,19)(-.03358209,.07835821){134}{\line(0,1){.07835821}}
%\end
%\emline(53.75,62.5)(43.25,58)
\multiput(53.75,62.5)(-.07835821,-.03358209){134}{\line(-1,0){.07835821}}
%\end
\put(49.75,25.5){\line(-1,-1){3.75}}
\put(47.25,59.75){\line(1,-1){3.75}}
\put(8.5,55.75){\line(1,1){3.75}}
\put(16,20.25){\line(-1,1){3.75}}
%\emline(59.5,28)(41,30.25)
\multiput(59.5,28)(-.2761194,.0335821){67}{\line(-1,0){.2761194}}
%\end
%\emline(44.75,69.5)(42.5,51)
\multiput(44.75,69.5)(-.0335821,-.2761194){67}{\line(0,-1){.2761194}}
%\end
%\emline(28.25,35.25)(55,33.75)
\multiput(28.25,35.25)(.5944444,-.0333333){45}{\line(1,0){.5944444}}
%\end
%\emline(59.5,28)(52.25,37)
\multiput(59.5,28)(-.03372093,.04186047){215}{\line(0,1){.04186047}}
%\end
%\emline(44.75,69.5)(35.75,62.25)
\multiput(44.75,69.5)(-.04186047,-.03372093){215}{\line(-1,0){.04186047}}
%\end
%\emline(56,32.75)(53.25,28.75)
\multiput(56,32.75)(-.0335366,-.0487805){82}{\line(0,-1){.0487805}}
%\end
%\emline(53.75,52.75)(34.25,35.25)
\multiput(53.75,52.75)(-.037572254,-.03371869){519}{\line(-1,0){.037572254}}
%\end
%\emline(24.5,27.25)(30.5,11.75)
\multiput(24.5,27.25)(.03370787,-.08707865){178}{\line(0,-1){.08707865}}
%\end
%\emline(46,34.75)(61.5,40.75)
\multiput(46,34.75)(.08707865,.03370787){178}{\line(1,0){.08707865}}
%\end
%\emline(30.5,11.75)(35.25,20)
\multiput(30.5,11.75)(.03368794,.05851064){141}{\line(0,1){.05851064}}
%\end
%\emline(61.5,40.75)(53.25,45.5)
\multiput(61.5,40.75)(-.05851064,.03368794){141}{\line(-1,0){.05851064}}
%\end
%\emline(56,44)(37,35)
\multiput(56,44)(-.071161049,-.033707865){267}{\line(-1,0){.071161049}}
%\end
%\emline(28.75,16.75)(33,16)
\multiput(28.75,16.75)(.184783,-.032609){23}{\line(1,0){.184783}}
%\end
%\emline(56.5,39)(57.25,43.25)
\multiput(56.5,39)(.032609,.184783){23}{\line(0,1){.184783}}
%\end
%\emline(29,27.25)(42.25,11)
\multiput(29,27.25)(.033715013,-.041348601){393}{\line(0,-1){.041348601}}
%\end
%\emline(46,39.25)(62.25,52.5)
\multiput(46,39.25)(.041348601,.033715013){393}{\line(1,0){.041348601}}
%\end
%\emline(34.75,71.25)(38,50.25)
\multiput(34.75,71.25)(.03350515,-.21649485){97}{\line(0,-1){.21649485}}
%\end
%\emline(34.75,71.25)(31.5,60.25)
\multiput(34.75,71.25)(-.03350515,-.11340206){97}{\line(0,-1){.11340206}}
%\end
%\emline(32.75,65)(35.75,65.25)
\multiput(32.75,65)(.375,.03125){8}{\line(1,0){.375}}
%\end
\put(38.5,64.25){\line(0,-1){22.5}}
%\emline(38.5,41.5)(32,62.25)
\multiput(38.5,41.5)(-.03367876,.10751295){193}{\line(0,1){.10751295}}
%\end
%\emline(34.25,55.25)(26.75,69.75)
\multiput(34.25,55.25)(-.03363229,.06502242){223}{\line(0,1){.06502242}}
%\end
%\emline(26.5,69.75)(25.25,58.75)
\multiput(26.5,69.75)(-.0328947,-.2894737){38}{\line(0,-1){.2894737}}
%\end
%\emline(25.5,62)(37.25,45.25)
\multiput(25.5,62)(.033667622,-.047994269){349}{\line(0,-1){.047994269}}
%\end
%\emline(31.25,54)(14.75,68.5)
\multiput(31.25,54)(-.038372093,.03372093){430}{\line(-1,0){.038372093}}
%\end
%\emline(14.75,68.5)(15.5,58.25)
\multiput(14.75,68.5)(.032609,-.445652){23}{\line(0,-1){.445652}}
%\end
%\emline(15.5,61.5)(35.75,47.5)
\multiput(15.5,61.5)(.048795181,-.03373494){415}{\line(1,0){.048795181}}
%\end
%\emline(31.5,50.25)(9.5,54)
\multiput(31.5,50.25)(-.19642857,.03348214){112}{\line(-1,0){.19642857}}
%\end
%\emline(23.75,56)(4,61)
\multiput(23.75,56)(-.13255034,.03355705){149}{\line(-1,0){.13255034}}
%\end
%\emline(4,61)(11.25,51.5)
\multiput(4,61)(.03372093,-.04418605){215}{\line(0,-1){.04418605}}
%\end
%\emline(1.25,50.25)(12,43.25)
\multiput(1.25,50.25)(.05168269,-.03365385){208}{\line(1,0){.05168269}}
%\end
%\emline(9.25,45.25)(30,50.5)
\multiput(9.25,45.25)(.13301282,.03365385){156}{\line(1,0){.13301282}}
%\end
%\emline(1.5,50.25)(18.25,52.75)
\multiput(1.5,50.25)(.2233333,.0333333){75}{\line(1,0){.2233333}}
%\end
%\emline(28,50)(7.5,35.75)
\multiput(28,50)(-.048463357,-.033687943){423}{\line(-1,0){.048463357}}
%\end
%\emline(14.75,46.75)(2,35.75)
\multiput(14.75,46.75)(-.038990826,-.033639144){327}{\line(-1,0){.038990826}}
%\end
%\emline(2,35.75)(10.5,35.5)
\multiput(2,35.75)(1.0625,-.03125){8}{\line(1,0){1.0625}}
%\end
%\emline(25.5,48)(10.75,27.5)
\multiput(25.5,48)(-.033675799,-.046803653){438}{\line(0,-1){.046803653}}
%\end
%\emline(14.25,40.5)(2.75,26.5)
\multiput(14.25,40.5)(-.03372434,-.041055718){341}{\line(0,-1){.041055718}}
%\end
%\emline(2.5,26.5)(14,27.75)
\multiput(2.5,26.5)(.3026316,.0328947){38}{\line(1,0){.3026316}}
%\end
%\emline(24.5,46.5)(18,21.5)
\multiput(24.5,46.5)(-.03367876,-.12953368){193}{\line(0,-1){.12953368}}
%\end
%\emline(16.25,35.25)(10.5,16.5)
\multiput(16.25,35.25)(-.03362573,-.10964912){171}{\line(0,-1){.10964912}}
%\end
%\emline(10.5,16.5)(20,23)
\multiput(10.5,16.5)(.0492228,.03367876){193}{\line(1,0){.0492228}}
%\end
%\emline(24,44.5)(24.5,19.25)
\multiput(24,44.5)(.033333,-1.683333){15}{\line(0,-1){1.683333}}
%\end
%\emline(19.75,29.25)(20,13.75)
\multiput(19.75,29.25)(.03125,-1.9375){8}{\line(0,-1){1.9375}}
%\end
%\emline(20,13.75)(26.75,22)
\multiput(20,13.75)(.03358209,.04104478){201}{\line(0,1){.04104478}}
%\end
%\emline(24,39.75)(42.75,19)
\multiput(24,39.75)(.033723022,-.037320144){556}{\line(0,-1){.037320144}}
%\end
%\emline(24.25,35)(33.75,17.75)
\multiput(24.25,35)(.033687943,-.061170213){282}{\line(0,-1){.061170213}}
%\end
\put(26,65){\line(1,0){3.25}}
%\emline(15,63.5)(18.75,65.5)
\multiput(15,63.5)(.0625,.0333333){60}{\line(1,0){.0625}}
%\end
%\emline(6.5,47)(6.75,51.25)
\multiput(6.5,47)(.03125,.53125){8}{\line(0,1){.53125}}
%\end
%\emline(6,36)(5.5,39)
\multiput(6,36)(-.033333,.2){15}{\line(0,1){.2}}
%\end
%\emline(8.5,27.5)(6.25,30.75)
\multiput(8.5,27.5)(-.0335821,.0485075){67}{\line(0,1){.0485075}}
%\end
%\emline(40,65.5)(43.75,62.75)
\multiput(40,65.5)(.0457317,-.0335366){82}{\line(1,0){.0457317}}
%\end
%\emline(20,19.75)(23,17.75)
\multiput(20,19.75)(.05,-.0333333){60}{\line(1,0){.05}}
%\end
%\emline(38.25,16.25)(42.75,17.25)
\multiput(38.25,16.25)(.15,.033333){30}{\line(1,0){.15}}
%\end
%\emline(56.75,48)(56,52.5)
\multiput(56.75,48)(-.032609,.195652){23}{\line(0,1){.195652}}
%\end
\put(30.75,3.75){\makebox(0,0)[cc]{Figure 5.4}}
\end{picture}
\end{center}

\begin{description}

  \item[Remark:] We have been concerned to give an algorithm; that is, a systematic procedure for folding regular $\left\{ \frac{b}{a} \right\}$-gons. Of course, we do not advocate this algorithm as the best procedure in any particular case. For example, if $a < \left( \frac{1}{2} \right) b'$, it is plainly best to take the tape prepared for folding the $\left\{ \frac{b'}{a} \right\}$-gon, on which there will already be a crease line making an angle of $\left( \frac{a}{b'} \right) \pi$ with the forward direction of the top of the tape, and then to bisect this angle by folding $k$ times to produce the required angle of $\left( \frac{a}{b} \right) \pi$. This remark may seem always applicable if $a = 1$, that is when regular \emph{convex} polygons are in question; however, with $a = 1$, this simplified procedure actually coincides with our algorithm!

\end{description}



\newpage

\section*{Appendix - The Further Generation \\ Of Folding Numbers}
\addcontentsline{toc}{section}{\numberline{}Appendix - The Further Generation Of Folding Numbers}
\markboth{Appendix - Further Generation}{Appendix - Further Generation}


In Section 3, we described the set of \emph{$t$-folding numbers}, $F_{t}$, consisting of integers $s$ of the form
\begin{equation}
s = \frac{t^{xy}-1}{t^{x}-1}, \quad x \geq 1, \quad y \geq 2  \tag{A1}
\end{equation}

We report here the main results of [10], in which we generalize even further, to consider the set $F_{tu}$ of \emph{$(t,u)$-folding numbers} $s$, being integers of the form
\begin{equation}
s = \frac{t^{xy}-u^{xy}}{t^{x}-u^{x}}, \quad x \geq 1, \quad y \geq 2  \tag{A2}
\end{equation}

In (A2), $t,u$ are fixed positive integers with $t > u \geq 1$ and $\gcd(t,u) = 1$. We are content in this appendix simply to report the results; the proofs will be found in [10]. In each case, we indicate the result of Section 3 which is being generalized.

\begin{description}

  \item[Theorem A1:] If $\gcd(a,b) = d$, then $\gcd (t^{a}-u^{a}, t^{b}-u^{b}) = (t^{d}-u^{d})$. \emph{(Theorem 3.1)}

\end{description}

Now let $c$ be a positive integer $\geq 2$, prime to $t$ and $u$; and let $\frac{F}{c}$ be the set defined by the rule
\begin{equation*}
\frac{F}{c} = \{ s \in F_{tu} \text{ such that } c\ |\ s \}
\end{equation*}

Then $\frac{F}{c}$ is non-empty; indeed, we may fix $x$ in (A2) and the set of $s$ of the form (A2) with a given $x$ such that $c\ |\ s$ is non-empty.

We fix $x=1$ in (A2), let $y_{{.07532}}
\multiput(128.95,70.18)(.07532,-.03125){12}{\line(1,0){.07532}}
\multiput(130.76,69.43)(.07532,-.03125){12}{\line(1,0){.07532}}
\multiput(132.56,68.68)(.07532,-.03125){12}{\line(1,0){.07532}}
\multiput(134.37,67.93)(.07532,-.03125){12}{\line(1,0){.07532}}
%\end
%\dashline{1}(136.25,67.25)(140.75,73.75)
\multiput(136.18,67.18)(.033333,.048148){15}{\line(0,1){.048148}}
\multiput(137.18,68.62)(.033333,.048148){15}{\line(0,1){.048148}}
\multiput(138.18,70.07)(.033333,.048148){15}{\line(0,1){.048148}}
\multiput(139.18,71.51)(.033333,.048148){15}{\line(0,1){.048148}}
\multiput(140.18,72.96)(.033333,.048148){15}{\line(0,1){.048148}}
%\end
%\dashline{1}(43,86.5)(54.5,67.5)
\multiput(42.93,86.43)(.031944,-.052778){15}{\line(0,-1){.052778}}
\multiput(43.89,84.85)(.031944,-.052778){15}{\line(0,-1){.052778}}
\multiput(44.85,83.26)(.031944,-.052778){15}{\line(0,-1){.052778}}
\multiput(45.8,81.68)(.031944,-.052778){15}{\line(0,-1){.052778}}
\multiput(46.76,80.1)(.031944,-.052778){15}{\line(0,-1){.052778}}
\multiput(47.72,78.51)(.031944,-.052778){15}{\line(0,-1){.052778}}
\multiput(48.68,76.93)(.031944,-.052778){15}{\line(0,-1){.052778}}
\multiput(49.64,75.35)(.031944,-.052778){15}{\line(0,-1){.052778}}
\multiput(50.6,73.76)(.031944,-.052778){15}{\line(0,-1){.052778}}
\multiput(51.55,72.18)(.031944,-.052778){15}{\line(0,-1){.052778}}
\multiput(52.51,70.6)(.031944,-.052778){15}{\line(0,-1){.052778}}
\multiput(53.47,69.01)(.031944,-.052778){15}{\line(0,-1){.052778}}
%\end
%\dashline{1}(89.75,86.25)(112.25,67)
\multiput(89.68,86.18)(.0382,-.032683){19}{\line(1,0){.0382}}
\multiput(91.13,84.94)(.0382,-.032683){19}{\line(1,0){.0382}}
\multiput(92.58,83.7)(.0382,-.032683){19}{\line(1,0){.0382}}
\multiput(94.03,82.45)(.0382,-.032683){19}{\line(1,0){.0382}}
\multiput(95.49,81.21)(.0382,-.032683){19}{\line(1,0){.0382}}
\multiput(96.94,79.97)(.0382,-.032683){19}{\line(1,0){.0382}}
\multiput(98.39,78.73)(.0382,-.032683){19}{\line(1,0){.0382}}
\multiput(99.84,77.49)(.0382,-.032683){19}{\line(1,0){.0382}}
\multiput(101.29,76.24)(.0382,-.032683){19}{\line(1,0){.0382}}
\multiput(102.74,75)(.0382,-.032683){19}{\line(1,0){.0382}}
\multiput(104.2,73.76)(.0382,-.032683){19}{\line(1,0){.0382}}
\multiput(105.65,72.52)(.0382,-.032683){19}{\line(1,0){.0382}}
\multiput(107.1,71.28)(.0382,-.032683){19}{\line(1,0){.0382}}
\multiput(108.55,70.03)(.0382,-.032683){19}{\line(1,0){.0382}}
\multiput(110,68.79)(.0382,-.032683){19}{\line(1,0){.0382}}
\multiput(111.45,67.55)(.0382,-.032683){19}{\line(1,0){.0382}}
%\end
%\vector(8.5,95)(13.5,95)
\put(13.5,95){\vector(1,0){.07}}\put(8.5,95){\line(1,0){5}}
%\end
%\vector(9.5,61)(14.5,61)
\put(14.5,61){\vector(1,0){.07}}\put(9.5,61){\line(1,0){5}}
%\end
\put(13.5,89){\makebox(0,0)[cc]{$A_{1}$}}
\put(27.5,65){\makebox(0,0)[cc]{$A_{2}$}}
\put(43.25,89){\makebox(0,0)[cc]{$A_{3}$}}
\put(54.25,65){\makebox(0,0)[cc]{$A_{4}$}}
\put(77.25,65){\makebox(0,0)[cc]{$A_{5}$}}
\put(89.5,89){\makebox(0,0)[cc]{$A_{6}$}}
\put(113,65){\makebox(0,0)[cc]{$A_{7}$}}
\put(136,65){\makebox(0,0)[cc]{$A_{8}$}}
\put(72.25,57.75){\makebox(0,0)[cc]{Figure 2.2}}
\end{picture}
\end{center}

It is a surprising fact (for which we will indicate a line of proof later in this section) that starting with \emph{any} putative angle $\left( \frac{a}{b} \right) \pi$, where $a,b$ satisfy our conditions, we will always obtain by our rules a primary folding procedure which produces angles that get closer and closer to the putative angle.

Furthermore, notice that in our particular case, starting with a putative angle of $\frac{\pi}{7}$ at the top of the tape produced a putative angle of $\frac{3\pi}{7}$ at the bottom of the tape. And, in fact, there are also pairs of angles, at equally spaced intervals, of $\frac{2\pi}{7}$ along the top edge of the tape. Thus, we may construct \emph{all} of the regular star $7$-gons from this tape. Here are the details.

Figure 2.3(a) shows how the $D^{2} U^{1}$ tape should look after a few folds (the first few triangles of the tape, which may have been irregular, have been discarded). If the FAT-algorithm is applied to the pairs of angles measuring $\frac{\pi}{7}$, at $A_{9}, A_{12}, \cdots, A_{27}$, then $A_{30}$ will coincide with $A_{9}$, and the top edge of the tape will form a regular convex $7$-gon (see Figure 2.3(d) ). 

Similarly, if the FAT-algorithm is applied to the pairs of angles measuring $\frac{2\pi}{7}$ at suitably spaced locations where they occur along the top edge of the tape in Figure 2.3(a), say at $A_{9}, A_{15}, \cdots, A_{45}$, then the top edge of the tape can be ``woven'' through itself to form the regular convex $\left\{ \frac{7}{2} \right\}$-gon shown in Figure 2.4.

Finally, if the tape is flipped so that the bottom edge becomes the top, and the FAT-algorithm is applied to the pairs of angles measuring $\frac{3\pi}{7}$ at equally spaced locations that occur along the top, say at $A_{11}, A_{17}, \cdots, A_{47}$, then the top edge of the tape will form a regular convex $\left\{ \frac{7}{3} \right\}$-gon, as shown in Figure 2.5(a), and when the excess at each point is folded under, the more conventional-looking star $\left\{ \frac{7}{3} \right\}$-gon of Figure 2.5(b) emerges. 

Of course, if the vertices at which the FAT-algorithm is performed are selected far enough apart, it is possible to braid the tape through itself to form a $\left\{ \frac{7}{3} \right\}$-gon with a hole in the middle. However, just as in our illustration, even though the top edge would be delineating the polygon, there would be extra paper sticking out at each point (which could, of course, be folded around the point as was done in Figure 2.5(b) ). (See Figures 2.4 and 2.5.)

\begin{center}
%TeXCAD Picture [fig2.3.pic]. Options:
%\grade{\on}
%\emlines{\off}
%\epic{\off}
%\beziermacro{\on}
%\reduce{\on}
%\snapping{\off}
%\quality{8.00}
%\graddiff{0.01}
%\snapasp{1}
%\zoom{4.0000}
\unitlength .95mm % = 2.85pt
\linethickness{0.4pt}
\ifx\plotpoint\undefined\newsavebox{\plotpoint}\fi % GNUPLOT compatibility
\begin{picture}(122.85,134)(0,0)
\put(2.75,131.25){\line(1,0){107.5}}
\put(109.25,119.5){\line(-1,0){106.5}}
%\dottedline(2.5,124.75)(11.25,119.75)
\multiput(2.43,124.68)(.7955,-.4545){12}{{\rule{.4pt}{.4pt}}}
%\end
%\dottedline(11.25,119.75)(14.5,131)
\multiput(11.18,119.68)(.25,.8654){14}{{\rule{.4pt}{.4pt}}}
%\end
%\dottedline(14.5,131)(36.25,119.75)
\multiput(14.43,130.93)(.87,-.45){26}{{\rule{.4pt}{.4pt}}}
%\end
%\dottedline(36.25,119.75)(41,131)
\multiput(36.18,119.68)(.3654,.8654){14}{{\rule{.4pt}{.4pt}}}
%\end
%\dottedline(41,131)(63.25,119.75)
\multiput(40.93,130.93)(.8558,-.4327){27}{{\rule{.4pt}{.4pt}}}
%\end
%\dottedline(63.25,119.75)(67.75,131)
\multiput(63.18,119.68)(.3462,.8654){14}{{\rule{.4pt}{.4pt}}}
%\end
%\dottedline(67.75,131)(90,119.75)
\multiput(67.68,130.93)(.8558,-.4327){27}{{\rule{.4pt}{.4pt}}}
%\end
%\dottedline(90,119.75)(95.75,131)
\multiput(89.93,119.68)(.4107,.8036){15}{{\rule{.4pt}{.4pt}}}
%\end
%\dottedline(95.75,131)(110.25,124)
\multiput(95.68,130.93)(.8529,-.4118){18}{{\rule{.4pt}{.4pt}}}
%\end
%\dottedline(14.75,131)(23,119.75)
\multiput(14.68,130.93)(.55,-.75){16}{{\rule{.4pt}{.4pt}}}
%\end
%\dottedline(41.25,130.75)(51.25,119.75)
\multiput(41.18,130.68)(.5882,-.6471){18}{{\rule{.4pt}{.4pt}}}
%\end
%\dottedline(67.5,131)(78,119.75)
\multiput(67.43,130.93)(.6176,-.6618){18}{{\rule{.4pt}{.4pt}}}
%\end
%\dottedline(95.75,131)(105.5,119.75)
\multiput(95.68,130.93)(.6094,-.7031){17}{{\rule{.4pt}{.4pt}}}
%\end
\put(6.25,104.25){\line(1,0){10.75}}
%\emline(17,104.25)(41.5,93.25)
\multiput(17,104.25)(.074923547,-.033639144){327}{\line(1,0){.074923547}}
%\end
%\emline(41.5,93.25)(52,81.75)
\multiput(41.5,93.25)(.033653846,-.036858974){312}{\line(0,-1){.036858974}}
%\end
%\emline(42.25,74.75)(17.5,104.5)
\multiput(42.25,74.75)(-.033719346,.0405313351){734}{\line(0,1){.0405313351}}
%\end
\put(7.75,94){\line(1,0){18.5}}
%\dottedline(17.75,104.25)(12.25,94)
\multiput(17.68,104.18)(-.4231,-.7885){14}{{\rule{.4pt}{.4pt}}}
%\end
%\dottedline(12.25,94)(6,98.25)
\multiput(12.18,93.93)(-.6944,.4722){10}{{\rule{.4pt}{.4pt}}}
%\end
%\dashline{1}(41.5,93)(33.75,85)
\multiput(41.43,92.93)(-.032292,-.033333){20}{\line(0,-1){.033333}}
\multiput(40.14,91.6)(-.032292,-.033333){20}{\line(0,-1){.033333}}
\multiput(38.85,90.26)(-.032292,-.033333){20}{\line(0,-1){.033333}}
\multiput(37.55,88.93)(-.032292,-.033333){20}{\line(0,-1){.033333}}
\multiput(36.26,87.6)(-.032292,-.033333){20}{\line(0,-1){.033333}}
\multiput(34.97,86.26)(-.032292,-.033333){20}{\line(0,-1){.033333}}
%\end
%\dashline{1}(33.75,85)(48,85.75)
\put(33.68,84.93){\line(1,0){.95}}
\put(35.58,85.03){\line(1,0){.95}}
\put(37.48,85.13){\line(1,0){.95}}
\put(39.38,85.23){\line(1,0){.95}}
\put(41.28,85.33){\line(1,0){.95}}
\put(43.18,85.43){\line(1,0){.95}}
\put(45.08,85.53){\line(1,0){.95}}
\put(46.98,85.63){\line(1,0){.95}}
%\end
%\dashline{1}(34,84.75)(48.75,79.75)
\multiput(33.93,84.68)(.09219,-.03125){10}{\line(1,0){.09219}}
\multiput(35.77,84.05)(.09219,-.03125){10}{\line(1,0){.09219}}
\multiput(37.62,83.43)(.09219,-.03125){10}{\line(1,0){.09219}}
\multiput(39.46,82.8)(.09219,-.03125){10}{\line(1,0){.09219}}
\multiput(41.3,82.18)(.09219,-.03125){10}{\line(1,0){.09219}}
\multiput(43.15,81.55)(.09219,-.03125){10}{\line(1,0){.09219}}
\multiput(44.99,80.93)(.09219,-.03125){10}{\line(1,0){.09219}}
\multiput(46.84,80.3)(.09219,-.03125){10}{\line(1,0){.09219}}
%\end
\put(75.25,104){\line(1,0){9.25}}
%\emline(84.5,104)(112,74.25)
\multiput(84.5,104)(.0337009804,-.0364583333){816}{\line(0,-1){.0364583333}}
%\end
%\emline(100.75,68.5)(88.75,81.75)
\multiput(100.75,68.5)(-.033707865,.037219101){356}{\line(0,1){.037219101}}
%\end
%\emline(88.33,81.75)(84.33,104)
\multiput(88.33,81.75)(-.03361345,.18697479){119}{\line(0,1){.18697479}}
%\end
\put(76.25,93){\line(1,0){10.25}}
%\emline(88.5,81.75)(94.75,93)
\multiput(88.5,81.75)(.03360215,.06048387){186}{\line(0,1){.06048387}}
%\end
%\dottedline(75.75,97.5)(81.75,93)
\multiput(75.68,97.43)(.6667,-.5){10}{{\rule{.4pt}{.4pt}}}
%\end
%\dottedline(81.75,93)(84.5,104)
\multiput(81.68,92.93)(.2115,.8462){14}{{\rule{.4pt}{.4pt}}}
%\end
%\dottedline(88.25,81.75)(102.5,84.25)
\multiput(88.18,81.68)(.8906,.1563){17}{{\rule{.4pt}{.4pt}}}
%\end
%\dottedline(102.5,84.25)(106.5,70)
\multiput(102.43,84.18)(.25,-.8906){17}{{\rule{.4pt}{.4pt}}}
%\end
%\dottedline(106.5,70)(106.5,70.25)
\multiput(106.43,69.93)(0,.125){3}{{\rule{.4pt}{.4pt}}}
%\end
%\dottedline(102.25,84.5)(97.5,72.5)
\multiput(102.18,84.43)(-.3393,-.8571){15}{{\rule{.4pt}{.4pt}}}
%\end
%\emline(38,64.66)(11.75,58.41)
\multiput(38,64.66)(-.14112903,-.03360215){186}{\line(-1,0){.14112903}}
%\end
%\emline(11.75,58.41)(1,38.16)
\multiput(11.75,58.41)(-.03369906,-.063479624){319}{\line(0,-1){.063479624}}
%\end
%\emline(1,38.16)(13.25,16.66)
\multiput(1,38.16)(.033653846,-.059065934){364}{\line(0,-1){.059065934}}
%\end
%\emline(13.25,16.66)(40.5,13.41)
\multiput(13.25,16.66)(.28092784,-.03350515){97}{\line(1,0){.28092784}}
%\end
%\emline(40.5,13.41)(59.5,27.16)
\multiput(40.5,13.41)(.046568627,.03370098){408}{\line(1,0){.046568627}}
%\end
\put(59.5,27.16){\line(0,1){22.75}}
%\emline(59.5,49.91)(38,64.66)
\multiput(59.5,49.91)(-.049086758,.033675799){438}{\line(-1,0){.049086758}}
%\end
\put(43.5,55.41){\line(-5,-1){5}}
%\emline(38.5,54.41)(26.5,61.91)
\multiput(38.5,54.41)(-.05381166,.03363229){223}{\line(-1,0){.05381166}}
%\end
%\emline(38.25,54.41)(11.75,58.16)
\multiput(38.25,54.41)(-.23660714,.03348214){112}{\line(-1,0){.23660714}}
%\end
%\emline(23,56.66)(20.25,51.41)
\multiput(23,56.66)(-.0335366,-.0640244){82}{\line(0,-1){.0640244}}
%\end
%\emline(20.25,51.41)(1.5,38.16)
\multiput(20.25,51.41)(-.047709924,-.033715013){393}{\line(-1,0){.047709924}}
%\end
%\emline(9,43.16)(12.25,38.91)
\multiput(9,43.16)(.03350515,-.04381443){97}{\line(0,-1){.04381443}}
%\end
%\emline(12.25,38.91)(13.25,16.91)
\multiput(12.25,38.91)(.033333,-.733333){30}{\line(0,-1){.733333}}
%\end
\put(13,25.66){\line(1,0){6}}
%\emline(19,25.66)(40.5,13.16)
\multiput(19,25.66)(.057951482,-.033692722){371}{\line(1,0){.057951482}}
%\end
\put(31.5,18.41){\line(1,1){3.75}}
%\emline(35.75,22.16)(59.5,27.16)
\multiput(35.75,22.16)(.15939597,.03355705){149}{\line(1,0){.15939597}}
%\end
%\emline(49.75,25.41)(48.5,30.16)
\multiput(49.75,25.41)(-.0328947,.125){38}{\line(0,1){.125}}
%\end
%\emline(48.5,30.16)(59.5,49.66)
\multiput(48.5,30.16)(.033639144,.059633028){327}{\line(0,1){.059633028}}
%\end
%\emline(54.5,41.41)(50.25,43.91)
\multiput(54.5,41.41)(-.0566667,.0333333){75}{\line(-1,0){.0566667}}
%\end
%\emline(50.25,43.91)(38,64.16)
\multiput(50.25,43.91)(-.033653846,.055631868){364}{\line(0,1){.055631868}}
%\end
%\emline(50.5,44.16)(51.25,55.41)
\multiput(50.5,44.16)(.032609,.48913){23}{\line(0,1){.48913}}
%\end
%\emline(48.25,29.86)(59.5,37.86)
\multiput(48.25,29.86)(.04726891,.03361345){238}{\line(1,0){.04726891}}
%\end
%\emline(35.13,22.24)(49.88,19.99)
\multiput(35.13,22.24)(.2201493,-.0335821){67}{\line(1,0){.2201493}}
%\end
%\emline(19.25,25.41)(27.75,15.16)
\multiput(19.25,25.41)(.033730159,-.040674603){252}{\line(0,-1){.040674603}}
%\end
%\emline(12.1,38.56)(7.85,27.31)
\multiput(12.1,38.56)(-.03373016,-.08928571){126}{\line(0,-1){.08928571}}
%\end
%\emline(20,51.66)(6.5,48.91)
\multiput(20,51.66)(-.1646341,-.0335366){82}{\line(-1,0){.1646341}}
%\end
%\dashline{1}(38,54.66)(38,64.16)
\put(37.93,54.59){\line(0,1){.95}}
\put(37.93,56.49){\line(0,1){.95}}
\put(37.93,58.39){\line(0,1){.95}}
\put(37.93,60.29){\line(0,1){.95}}
\put(37.93,62.19){\line(0,1){.95}}
%\end
%\dashline{1}(20.15,51.81)(11.4,57.81)
\multiput(20.08,51.74)(-.048611,.033333){15}{\line(-1,0){.048611}}
\multiput(18.62,52.74)(-.048611,.033333){15}{\line(-1,0){.048611}}
\multiput(17.16,53.74)(-.048611,.033333){15}{\line(-1,0){.048611}}
\multiput(15.7,54.74)(-.048611,.033333){15}{\line(-1,0){.048611}}
\multiput(14.25,55.74)(-.048611,.033333){15}{\line(-1,0){.048611}}
\multiput(12.79,56.74)(-.048611,.033333){15}{\line(-1,0){.048611}}
%\end
%\dashline{1}(12,38.66)(1.25,37.91)
\put(11.93,38.59){\line(-1,0){.977}}
\put(9.98,38.45){\line(-1,0){.977}}
\put(8.02,38.32){\line(-1,0){.977}}
\put(6.07,38.18){\line(-1,0){.977}}
\put(4.11,38.04){\line(-1,0){.977}}
\put(2.16,37.91){\line(-1,0){.977}}
%\end
%\dashline{1}(19.25,25.62)(13,16.37)
\multiput(19.18,25.55)(-.032552,-.048177){16}{\line(0,-1){.048177}}
\multiput(18.14,24.01)(-.032552,-.048177){16}{\line(0,-1){.048177}}
\multiput(17.1,22.47)(-.032552,-.048177){16}{\line(0,-1){.048177}}
\multiput(16.05,20.92)(-.032552,-.048177){16}{\line(0,-1){.048177}}
\multiput(15.01,19.38)(-.032552,-.048177){16}{\line(0,-1){.048177}}
\multiput(13.97,17.84)(-.032552,-.048177){16}{\line(0,-1){.048177}}
%\end
%\dashline{1}(35.5,21.91)(40,13.41)
\multiput(35.43,21.84)(.031469,-.059441){13}{\line(0,-1){.059441}}
\multiput(36.25,20.29)(.031469,-.059441){13}{\line(0,-1){.059441}}
\multiput(37.07,18.75)(.031469,-.059441){13}{\line(0,-1){.059441}}
\multiput(37.88,17.2)(.031469,-.059441){13}{\line(0,-1){.059441}}
\multiput(38.7,15.66)(.031469,-.059441){13}{\line(0,-1){.059441}}
\multiput(39.52,14.11)(.031469,-.059441){13}{\line(0,-1){.059441}}
%\end
%\dashline{1}(59.57,26.84)(48.57,29.84)
\multiput(59.5,26.77)(-.12088,.03297){7}{\line(-1,0){.12088}}
\multiput(57.81,27.23)(-.12088,.03297){7}{\line(-1,0){.12088}}
\multiput(56.12,27.69)(-.12088,.03297){7}{\line(-1,0){.12088}}
\multiput(54.42,28.15)(-.12088,.03297){7}{\line(-1,0){.12088}}
\multiput(52.73,28.62)(-.12088,.03297){7}{\line(-1,0){.12088}}
\multiput(51.04,29.08)(-.12088,.03297){7}{\line(-1,0){.12088}}
\multiput(49.35,29.54)(-.12088,.03297){7}{\line(-1,0){.12088}}
%\end
%\dashline{1}(50.5,43.91)(59.25,49.66)
\multiput(50.43,43.84)(.049716,.03267){16}{\line(1,0){.049716}}
\multiput(52.02,44.89)(.049716,.03267){16}{\line(1,0){.049716}}
\multiput(53.61,45.93)(.049716,.03267){16}{\line(1,0){.049716}}
\multiput(55.2,46.98)(.049716,.03267){16}{\line(1,0){.049716}}
\multiput(56.79,48.02)(.049716,.03267){16}{\line(1,0){.049716}}
\multiput(58.38,49.07)(.049716,.03267){16}{\line(1,0){.049716}}
%\end
%\emline(63.5,50)(83.5,64.75)
\multiput(63.5,50)(.0456621,.033675799){438}{\line(1,0){.0456621}}
%\end
%\emline(83.5,64.75)(109.75,60)
\multiput(83.5,64.75)(.18617021,-.03368794){141}{\line(1,0){.18617021}}
%\end
%\emline(109.75,60)(122.75,39.25)
\multiput(109.75,60)(.033678756,-.053756477){386}{\line(0,-1){.053756477}}
%\end
%\emline(122.75,39.25)(112.5,18.25)
\multiput(122.75,39.25)(-.033717105,-.069078947){304}{\line(0,-1){.069078947}}
%\end
%\emline(112.5,18.25)(86.5,12.75)
\multiput(112.5,18.25)(-.15853659,-.03353659){164}{\line(-1,0){.15853659}}
%\end
%\emline(86.5,12.75)(63.5,26.5)
\multiput(86.5,12.75)(-.056372549,.03370098){408}{\line(-1,0){.056372549}}
%\end
\put(63.5,26.5){\line(0,1){23.5}}
%\emline(73,56.75)(84.5,54.25)
\multiput(73,56.75)(.1533333,-.0333333){75}{\line(1,0){.1533333}}
%\end
%\emline(84.5,54.25)(98.25,57.5)
\multiput(84.5,54.25)(.14175258,.03350515){97}{\line(1,0){.14175258}}
%\end
%\emline(95.5,62.25)(101,53.25)
\multiput(95.5,62.25)(.03353659,-.05487805){164}{\line(0,-1){.05487805}}
%\end
%\emline(101,53.25)(112.75,45.5)
\multiput(101,53.25)(.05108696,-.03369565){230}{\line(1,0){.05108696}}
%\end
%\emline(115.5,51.25)(110.25,41)
\multiput(115.5,51.25)(-.03365385,-.06570513){156}{\line(0,-1){.06570513}}
%\end
%\emline(110.25,41.25)(111.75,28.25)
\multiput(110.25,41.25)(.0333333,-.2888889){45}{\line(0,-1){.2888889}}
%\end
%\emline(117.75,28.75)(105.5,27.25)
\multiput(117.75,28.75)(-.2722222,-.0333333){45}{\line(-1,0){.2722222}}
%\end
%\emline(105.5,27.25)(94.75,19)
\multiput(105.5,27.25)(-.043877551,-.033673469){245}{\line(-1,0){.043877551}}
%\end
\put(99.5,15.5){\line(-1,0){.25}}
%\emline(99.25,15.5)(90,21.5)
\multiput(99.25,15.5)(-.05196629,.03370787){178}{\line(-1,0){.05196629}}
%\end
%\emline(90,21.5)(76,24.25)
\multiput(90,21.5)(-.1707317,.0335366){82}{\line(-1,0){.1707317}}
%\end
\put(76.75,19){\line(0,1){.25}}
%\emline(76.75,19.25)(75,28.75)
\multiput(76.75,19.25)(-.0336538,.1826923){52}{\line(0,1){.1826923}}
%\end
%\emline(75,28.75)(68.5,41.5)
\multiput(75,28.75)(-.03367876,.06606218){193}{\line(0,1){.06606218}}
%\end
%\emline(63.44,38.5)(71.69,43.5)
\multiput(63.44,38.5)(.05536913,.03355705){149}{\line(1,0){.05536913}}
%\end
\put(71.69,43.68){\line(1,2){6}}
%\emline(94,18.75)(96,19.75)
\multiput(94,18.75)(.066667,.033333){30}{\line(1,0){.066667}}
%\end
%\dashline{1}(63.5,50)(72.25,43)
\multiput(63.43,49.93)(.040509,-.032407){18}{\line(1,0){.040509}}
\multiput(64.89,48.76)(.040509,-.032407){18}{\line(1,0){.040509}}
\multiput(66.35,47.6)(.040509,-.032407){18}{\line(1,0){.040509}}
\multiput(67.8,46.43)(.040509,-.032407){18}{\line(1,0){.040509}}
\multiput(69.26,45.26)(.040509,-.032407){18}{\line(1,0){.040509}}
\multiput(70.72,44.1)(.040509,-.032407){18}{\line(1,0){.040509}}
%\end
%\dashline{1}(64.25,26.5)(75.25,28.25)
\multiput(64.18,26.43)(.18333,.02917){5}{\line(1,0){.18333}}
\multiput(66.01,26.72)(.18333,.02917){5}{\line(1,0){.18333}}
\multiput(67.85,27.01)(.18333,.02917){5}{\line(1,0){.18333}}
\multiput(69.68,27.3)(.18333,.02917){5}{\line(1,0){.18333}}
\multiput(71.51,27.6)(.18333,.02917){5}{\line(1,0){.18333}}
\multiput(73.35,27.89)(.18333,.02917){5}{\line(1,0){.18333}}
%\end
%\dashline{1}(87,13)(89.75,21.5)
\multiput(86.93,12.93)(.03056,.09444){9}{\line(0,1){.09444}}
\multiput(87.48,14.63)(.03056,.09444){9}{\line(0,1){.09444}}
\multiput(88.03,16.33)(.03056,.09444){9}{\line(0,1){.09444}}
\multiput(88.58,18.03)(.03056,.09444){9}{\line(0,1){.09444}}
\multiput(89.13,19.73)(.03056,.09444){9}{\line(0,1){.09444}}
%\end
%\dashline{1}(105.5,27.25)(112.5,18.25)
\multiput(105.43,27.18)(.033654,-.043269){16}{\line(0,-1){.043269}}
\multiput(106.51,25.8)(.033654,-.043269){16}{\line(0,-1){.043269}}
\multiput(107.58,24.41)(.033654,-.043269){16}{\line(0,-1){.043269}}
\multiput(108.66,23.03)(.033654,-.043269){16}{\line(0,-1){.043269}}
\multiput(109.74,21.64)(.033654,-.043269){16}{\line(0,-1){.043269}}
\multiput(110.81,20.26)(.033654,-.043269){16}{\line(0,-1){.043269}}
\multiput(111.89,18.87)(.033654,-.043269){16}{\line(0,-1){.043269}}
%\end
%\dashline{1}(110.35,41)(122.85,39)
\multiput(110.28,40.93)(.17857,-.02857){5}{\line(1,0){.17857}}
\multiput(112.07,40.64)(.17857,-.02857){5}{\line(1,0){.17857}}
\multiput(113.85,40.36)(.17857,-.02857){5}{\line(1,0){.17857}}
\multiput(115.64,40.07)(.17857,-.02857){5}{\line(1,0){.17857}}
\multiput(117.42,39.79)(.17857,-.02857){5}{\line(1,0){.17857}}
\multiput(119.21,39.5)(.17857,-.02857){5}{\line(1,0){.17857}}
\multiput(120.99,39.22)(.17857,-.02857){5}{\line(1,0){.17857}}
%\end
%\dashline{1}(101.25,53)(109.5,59.75)
\multiput(101.18,52.93)(.040441,.033088){17}{\line(1,0){.040441}}
\multiput(102.55,54.05)(.040441,.033088){17}{\line(1,0){.040441}}
\multiput(103.93,55.18)(.040441,.033088){17}{\line(1,0){.040441}}
\multiput(105.3,56.3)(.040441,.033088){17}{\line(1,0){.040441}}
\multiput(106.68,57.43)(.040441,.033088){17}{\line(1,0){.040441}}
\multiput(108.05,58.55)(.040441,.033088){17}{\line(1,0){.040441}}
%\end
%\dashline{1}(84.75,54.25)(84,65)
\put(84.68,54.18){\line(0,1){.977}}
\put(84.54,56.13){\line(0,1){.977}}
\put(84.41,58.09){\line(0,1){.977}}
\put(84.27,60.04){\line(0,1){.977}}
\put(84.13,62){\line(0,1){.977}}
\put(84,63.95){\line(0,1){.977}}
%\end
%\qbezvec(56.25,55.75)(55.38,62.38)(64,57.5)
\put(64,57.5){\vector(2,-1){.07}}\qbezier(56.25,55.75)(55.38,62.38)(64,57.5)
%\end
\put(57,110){\makebox(0,0)[cc]{(a)}}
\put(18.25,79.01){\makebox(0,0)[cc]{(b)}}
\put(80.58,78.98){\makebox(0,0)[cc]{(c)}}
\put(60.75,15.25){\makebox(0,0)[cc]{(d)}}
\put(14.5,133.5){\makebox(0,0)[cc]{$A_{9}$}}
\put(23,116.75){\makebox(0,0)[cc]{$A_{10}$}}
\put(36.5,116.5){\makebox(0,0)[cc]{$A_{11}$}}
\put(41,133.75){\makebox(0,0)[cc]{$A_{12}$}}
\put(50.75,116.5){\makebox(0,0)[cc]{$A_{13}$}}
\put(63.5,116.5){\makebox(0,0)[cc]{$A_{14}$}}
\put(67.75,134){\makebox(0,0)[cc]{$A_{15}$}}
\put(78.5,116.5){\makebox(0,0)[cc]{$A_{16}$}}
\put(90.25,116.75){\makebox(0,0)[cc]{$A_{17}$}}
\put(95.75,133.75){\makebox(0,0)[cc]{$A_{18}$}}
\put(105.25,116.5){\makebox(0,0)[cc]{$A_{19}$}}
\put(17.44,107.5){\makebox(0,0)[cc]{$A_{12}$}}
\put(24.25,92){\makebox(0,0)[cc]{$A_{13}$}}
\put(84.33,106.8){\makebox(0,0)[cc]{$A_{12}$}}
\put(106.25,86.5){\makebox(0,0)[cc]{$A_{15}$}}
\put(94.25,71.5){\makebox(0,0)[cc]{$A_{16}$}}
\put(15,126){\makebox(0,0)[cc]{$\frac{2\pi}{7}$}}
\put(21.25,125.25){\makebox(0,0)[cc]{$\frac{\pi}{7}$}}
\put(14.26,121.92){\makebox(0,0)[cc]{$\frac{3\pi}{7}$}}
\put(25.3,121.67){\makebox(0,0)[cc]{$\frac{5\pi}{7}$}}
\put(60.75,3){\makebox(0,0)[cc]{Figure 2.3}}
\end{picture}
\end{center}

\vspace{0.5cm}

\begin{center}
%TeXCAD Picture [fig2.4.pic]. Options:
%\grade{\on}
%\emlines{\off}
%\epic{\off}
%\beziermacro{\on}
%\reduce{\on}
%\snapping{\off}
%\quality{8.00}
%\graddiff{0.01}
%\snapasp{1}
%\zoom{4.0000}
\unitlength .95mm % = 2.85pt
\linethickness{0.4pt}
\ifx\plotpoint\undefined\newsavebox{\plotpoint}\fi % GNUPLOT compatibility
\begin{picture}(125.25,55)(0,0)
%\emline(19.5,53.75)(17,43.25)
\multiput(19.5,53.75)(-.0333333,-.14){75}{\line(0,-1){.14}}
%\end
%\emline(17,43.25)(5.25,39)
\multiput(17,43.25)(-.09325397,-.03373016){126}{\line(-1,0){.09325397}}
%\end
%\emline(5.25,39)(13,29.75)
\multiput(5.25,39)(.03369565,-.04021739){230}{\line(0,-1){.04021739}}
%\end
%\emline(13,29.75)(11,16)
\multiput(13,29.75)(-.0333333,-.2291667){60}{\line(0,-1){.2291667}}
%\end
%\emline(11,16)(25,15.5)
\multiput(11,16)(.933333,-.033333){15}{\line(1,0){.933333}}
%\end
%\emline(24.65,16.03)(32.4,7.28)
\multiput(24.65,16.03)(.03369565,-.03804348){230}{\line(0,-1){.03804348}}
%\end
%\emline(32.5,7.25)(41.5,15.5)
\multiput(32.5,7.25)(.036734694,.033673469){245}{\line(1,0){.036734694}}
%\end
%\emline(57.25,15)(53.75,27.75)
\multiput(57.25,15)(-.03365385,.12259615){104}{\line(0,1){.12259615}}
%\end
%\emline(53.75,27.75)(62.5,36.25)
\multiput(53.75,27.75)(.034722222,.033730159){252}{\line(1,0){.034722222}}
%\end
\put(62.5,36.25){\line(-5,2){12.5}}
%\emline(50,41.25)(47.25,53.75)
\multiput(50,41.25)(-.0335366,.152439){82}{\line(0,1){.152439}}
%\end
%\emline(50.88,36.83)(48.13,49.33)
\multiput(50.88,36.83)(-.0335366,.152439){82}{\line(0,1){.152439}}
%\end
%\emline(47.25,53.75)(32,49)
\multiput(47.25,53.75)(-.10815603,-.03368794){141}{\line(-1,0){.10815603}}
%\end
%\emline(32,49)(19.5,53.75)
\multiput(32,49)(-.08865248,.03368794){141}{\line(-1,0){.08865248}}
%\end
%\emline(19.75,53.75)(22.75,43.75)
\multiput(19.75,53.75)(.0337079,-.1123596){89}{\line(0,-1){.1123596}}
%\end
%\emline(22.75,43.75)(28.75,50.25)
\multiput(22.75,43.75)(.03370787,.03651685){178}{\line(0,1){.03651685}}
%\end
%\emline(23.25,43.75)(41.25,36.5)
\multiput(23.25,43.75)(.08372093,-.03372093){215}{\line(1,0){.08372093}}
%\end
%\emline(41.25,36.5)(38.5,45.5)
\multiput(41.25,36.5)(-.0335366,.1097561){82}{\line(0,1){.1097561}}
%\end
%\emline(38.5,45.5)(47.25,53.75)
\multiput(38.5,45.5)(.035714286,.033673469){245}{\line(1,0){.035714286}}
%\end
\put(38.75,45.5){\line(1,0){10.5}}
%\emline(41.5,36.25)(43,31.75)
\multiput(41.5,36.25)(.0333333,-.1){45}{\line(0,-1){.1}}
%\end
%\emline(43,31.75)(49.5,37)
\multiput(43,31.75)(.04166667,.03365385){156}{\line(1,0){.04166667}}
%\end
%\emline(49.5,37)(62.25,36.25)
\multiput(49.5,37)(.554348,-.032609){23}{\line(1,0){.554348}}
%\end
\put(49.75,37){\line(1,-1){6.5}}
%\emline(42.75,31.5)(36.25,24.75)
\multiput(42.75,31.5)(-.03367876,-.03497409){193}{\line(0,-1){.03497409}}
%\end
\put(36.25,24.75){\line(1,0){11.5}}
%\emline(47.75,24.75)(57,15)
\multiput(47.75,24.75)(.033636364,-.035454545){275}{\line(0,-1){.035454545}}
%\end
\put(53.75,27.75){\line(-1,-1){4}}
%\emline(49.75,23.75)(49.25,23.25)
\multiput(49.75,23.75)(-.033333,-.033333){15}{\line(0,-1){.033333}}
%\end
%\emline(48,24.25)(47,15.75)
\multiput(48,24.25)(-.033333,-.283333){30}{\line(0,-1){.283333}}
%\end
\put(36.5,24.75){\line(-1,0){6.75}}
%\emline(29.75,24.5)(36.5,18.25)
\multiput(29.75,24.5)(.03629032,-.03360215){186}{\line(1,0){.03629032}}
%\end
%\emline(36.5,18.25)(32.75,7.5)
\multiput(36.5,18.25)(-.03348214,-.09598214){112}{\line(0,-1){.09598214}}
%\end
%\emline(36.25,18)(27,13.25)
\multiput(36.25,18)(-.06560284,-.03368794){141}{\line(-1,0){.06560284}}
%\end
%\emline(21.75,20.25)(25.75,15)
\multiput(21.75,20.25)(.03361345,-.04411765){119}{\line(0,-1){.04411765}}
%\end
%\emline(11.5,16.25)(22,20.25)
\multiput(11.5,16.25)(.08823529,.03361345){119}{\line(1,0){.08823529}}
%\end
%\emline(22,20.25)(25.25,37)
\multiput(22,20.25)(.03350515,.17268041){97}{\line(0,1){.17268041}}
%\end
%\emline(21.75,20.5)(12.25,25)
\multiput(21.75,20.5)(-.07089552,.03358209){134}{\line(-1,0){.07089552}}
%\end
%\emline(13,29.25)(14,36)
\multiput(13,29.25)(.033333,.225){30}{\line(0,1){.225}}
%\end
%\emline(14,36)(5.5,39)
\multiput(14,36)(-.0955056,.0337079){89}{\line(-1,0){.0955056}}
%\end
%\emline(13.9,35.93)(35.65,38.68)
\multiput(13.9,35.93)(.2652439,.0335366){82}{\line(1,0){.2652439}}
%\end
%\emline(16.57,43.18)(22.32,45.18)
\multiput(16.57,43.18)(.0958333,.0333333){60}{\line(1,0){.0958333}}
%\end
%\emline(29.75,24.5)(24,30.75)
\multiput(29.75,24.5)(-.03362573,.03654971){171}{\line(0,1){.03654971}}
%\end
%\emline(24,30.75)(24.25,30.25)
\multiput(24,30.75)(.03125,-.0625){8}{\line(0,-1){.0625}}
%\end
%\dottedline(27,47.75)(35.75,38.75)
\multiput(26.93,47.68)(.625,-.6429){15}{{\rule{.4pt}{.4pt}}}
%\end
%\dottedline(35.75,38.75)(39,43.25)
\multiput(35.68,38.68)(.4643,.6429){8}{{\rule{.4pt}{.4pt}}}
%\end
%\emline(32.25,48.75)(38.75,46.25)
\multiput(32.25,48.75)(.0866667,-.0333333){75}{\line(1,0){.0866667}}
%\end
%\dottedline(44.5,45.25)(42,34.75)
\multiput(44.43,45.18)(-.2083,-.875){13}{{\rule{.4pt}{.4pt}}}
%\end
%\dottedline(42,34.75)(47.5,35.5)
\multiput(41.93,34.68)(.9167,.125){7}{{\rule{.4pt}{.4pt}}}
%\end
%\dottedline(53.5,33.25)(40.25,29)
\multiput(53.43,33.18)(-.8833,-.2833){16}{{\rule{.4pt}{.4pt}}}
%\end
%\dottedline(40.25,29)(45.25,25)
\multiput(40.18,28.93)(.625,-.5){9}{{\rule{.4pt}{.4pt}}}
%\end
%\dottedline(47,19.75)(34,24.75)
\multiput(46.93,19.68)(-.8667,.3333){16}{{\rule{.4pt}{.4pt}}}
%\end
%\dottedline(34,24.75)(33.5,21.75)
\multiput(33.93,24.68)(-.125,-.75){5}{{\rule{.4pt}{.4pt}}}
%\end
%\dottedline(31.25,16.25)(26.75,27.75)
\multiput(31.18,16.18)(-.3462,.8846){14}{{\rule{.4pt}{.4pt}}}
%\end
%\dottedline(26.75,27.75)(23,25.25)
\multiput(26.68,27.68)(-.75,-.5){6}{{\rule{.4pt}{.4pt}}}
%\end
%\dottedline(17.25,22.5)(24.25,33.25)
\multiput(17.18,22.43)(.5,.7679){15}{{\rule{.4pt}{.4pt}}}
%\end
%\dottedline(24.25,33.25)(20,36.5)
\multiput(24.18,33.18)(-.6071,.4643){8}{{\rule{.4pt}{.4pt}}}
%\end
%\dottedline(15.75,38.75)(28.5,37.75)
\multiput(15.68,38.68)(.9808,-.0769){14}{{\rule{.4pt}{.4pt}}}
%\end
%\dottedline(28.5,37.75)(25.25,42.75)
\multiput(28.43,37.68)(-.4643,.7143){8}{{\rule{.4pt}{.4pt}}}
%\end
%\emline(15,36.25)(14.5,42.25)
\multiput(15,36.25)(-.033333,.4){15}{\line(0,1){.4}}
%\end
\put(14.5,42.25){\line(0,-1){.25}}
%\emline(81,55)(77.75,41.75)
\multiput(81,55)(-.03350515,-.13659794){97}{\line(0,-1){.13659794}}
%\end
%\emline(77.75,41.75)(65.5,35.75)
\multiput(77.75,41.75)(-.06882022,-.03370787){178}{\line(-1,0){.06882022}}
%\end
%\emline(65.5,35.75)(74.25,27.5)
\multiput(65.5,35.75)(.035714286,-.033673469){245}{\line(1,0){.035714286}}
%\end
%\emline(74.25,27.5)(71.75,17)
\multiput(74.25,27.5)(-.0333333,-.14){75}{\line(0,-1){.14}}
%\end
\put(71.75,17){\line(0,1){.25}}
\put(71.75,17.25){\line(1,0){13.75}}
%\emline(85.5,17.25)(97,7)
\multiput(85.5,17.25)(.037828947,-.033717105){304}{\line(1,0){.037828947}}
%\end
%\emline(97,7)(104.75,16.5)
\multiput(97,7)(.03369565,.04130435){230}{\line(0,1){.04130435}}
%\end
%\emline(119.75,16.25)(115.25,27)
\multiput(119.75,16.25)(-.03358209,.08022388){134}{\line(0,1){.08022388}}
%\end
%\emline(115.25,27)(125.25,37.25)
\multiput(115.25,27)(.033670034,.034511785){297}{\line(0,1){.034511785}}
%\end
%\emline(125.25,37.25)(110.75,43)
\multiput(125.25,37.25)(-.08479532,.03362573){171}{\line(-1,0){.08479532}}
%\end
%\emline(110.75,43)(106.5,54.25)
\multiput(110.75,43)(-.03373016,.08928571){126}{\line(0,1){.08928571}}
%\end
%\emline(112.34,38.93)(108.09,50.18)
\multiput(112.34,38.93)(-.03373016,.08928571){126}{\line(0,1){.08928571}}
%\end
\put(84,44.25){\line(-5,-2){7.5}}
%\emline(79,46.5)(101.25,37.25)
\multiput(79,46.5)(.080909091,-.033636364){275}{\line(1,0){.080909091}}
%\end
%\emline(101.25,37.25)(96.75,50.25)
\multiput(101.25,37.25)(-.03358209,.09701493){134}{\line(0,1){.09701493}}
%\end
%\emline(103.9,29.65)(99.4,42.65)
\multiput(103.9,29.65)(-.03358209,.09701493){134}{\line(0,1){.09701493}}
%\end
%\emline(103.75,29)(114.75,41.25)
\multiput(103.75,29)(.033639144,.037461774){327}{\line(0,1){.037461774}}
%\end
\put(114.75,41.25){\line(0,-1){.25}}
%\emline(104,29)(100,25.25)
\multiput(104,29)(-.03571429,-.03348214){112}{\line(-1,0){.03571429}}
%\end
\put(116.65,23.75){\line(-1,0){25.5}}
%\emline(100.25,25.75)(98.5,24)
\multiput(100.25,25.75)(-.0336538,-.0336538){52}{\line(0,-1){.0336538}}
%\end
%\emline(102.25,13.75)(85,29.25)
\multiput(102.25,13.75)(-.0375,.033695652){460}{\line(-1,0){.0375}}
%\end
%\emline(81.5,17.5)(86.25,35)
\multiput(81.5,17.5)(.03368794,.12411348){141}{\line(0,1){.12411348}}
%\end
%\emline(95.25,39.25)(73,28.5)
\multiput(95.25,39.25)(-.069749216,-.03369906){319}{\line(-1,0){.069749216}}
%\end
%\emline(74.75,29.5)(74,27.25)
\multiput(74.75,29.5)(-.032609,-.097826){23}{\line(0,-1){.097826}}
%\end
%\emline(82.5,20.25)(86,17)
\multiput(82.5,20.25)(.03608247,-.03350515){97}{\line(1,0){.03608247}}
%\end
%\emline(115.25,26.75)(112.25,23.75)
\multiput(115.25,26.75)(-.0337079,-.0337079){89}{\line(0,-1){.0337079}}
%\end
%\dottedline(79.5,47)(98.25,46)
\multiput(79.43,46.93)(.9868,-.0526){20}{{\rule{.4pt}{.4pt}}}
%\end
%\dottedline(98.75,43.75)(91.25,41.25)
\multiput(98.68,43.68)(-.8333,-.2778){10}{{\rule{.4pt}{.4pt}}}
%\end
%\dottedline(91.25,41.25)(93,38)
\multiput(91.18,41.18)(.35,-.65){6}{{\rule{.4pt}{.4pt}}}
%\end
%\dottedline(99,38.5)(100.25,39.75)
\multiput(98.93,38.43)(.417,.417){4}{{\rule{.4pt}{.4pt}}}
%\end
%\dottedline(100.25,39.75)(110.25,36)
\multiput(100.18,39.68)(.8333,-.3125){13}{{\rule{.4pt}{.4pt}}}
%\end
%\dottedline(111.5,38)(97,50.25)
\multiput(111.43,37.93)(-.725,.6125){21}{{\rule{.4pt}{.4pt}}}
%\end
%\dottedline(115,41)(110.25,24)
\multiput(114.93,40.93)(-.25,-.8947){20}{{\rule{.4pt}{.4pt}}}
%\end
%\dottedline(108.25,24.25)(106.75,32.5)
\multiput(108.18,24.18)(-.1667,.9167){10}{{\rule{.4pt}{.4pt}}}
%\end
%\dottedline(106.75,32.5)(103.25,32)
\multiput(106.68,32.43)(-.875,-.125){5}{{\rule{.4pt}{.4pt}}}
%\end
%\dottedline(101.25,26.25)(101.75,24)
\multiput(101.18,26.18)(.167,-.75){4}{{\rule{.4pt}{.4pt}}}
%\end
%\dottedline(101.75,24)(96.75,18.75)
\multiput(101.68,23.93)(-.5556,-.5833){10}{{\rule{.4pt}{.4pt}}}
%\end
%\dottedline(97.75,18)(116.5,23.75)
\multiput(97.68,17.93)(.9375,.2875){21}{{\rule{.4pt}{.4pt}}}
%\end
%\dottedline(94.18,23.82)(94.43,20.82)
\multiput(94.11,23.75)(.0625,-.75){5}{{\rule{.4pt}{.4pt}}}
%\end
%\dottedline(94.43,20.82)(83.18,22.57)
\multiput(94.36,20.75)(-.9375,.1458){13}{{\rule{.4pt}{.4pt}}}
%\end
%\dottedline(82.5,21.75)(102,13.75)
\multiput(82.43,21.68)(.8864,-.3636){23}{{\rule{.4pt}{.4pt}}}
%\end
%\dottedline(81.25,17.25)(78,30.75)
\multiput(81.18,17.18)(-.2167,.9){16}{{\rule{.4pt}{.4pt}}}
%\end
%\dottedline(80,31.75)(84,26.5)
\multiput(79.93,31.68)(.5,-.6563){9}{{\rule{.4pt}{.4pt}}}
%\end
%\dottedline(84,26.5)(87.5,27)
\multiput(83.93,26.43)(.875,.125){5}{{\rule{.4pt}{.4pt}}}
%\end
%\dottedline(73.25,28.75)(85,44)
\multiput(73.18,28.68)(.5875,.7625){21}{{\rule{.4pt}{.4pt}}}
%\end
%\dottedline(86.75,43)(82.5,33.25)
\multiput(86.68,42.93)(-.3864,-.8864){12}{{\rule{.4pt}{.4pt}}}
%\end
%\dottedline(82.5,33.25)(85,31.75)
\multiput(82.43,33.18)(.625,-.375){5}{{\rule{.4pt}{.4pt}}}
%\end
%\qbezvec(58.75,47.25)(57.88,54.38)(66.5,49)
\put(66.5,49){\vector(3,-2){.07}}\qbezier(58.75,47.25)(57.88,54.38)(66.5,49)
%\end
%\emline(119.69,16.22)(99.36,16.4)
\multiput(119.69,16.22)(-3.38822,.02946){6}{\line(-1,0){3.38822}}
%\end
%\emline(57.23,14.89)(35.84,16.13)
\multiput(57.23,14.89)(-.5781076,.0334442){37}{\line(-1,0){.5781076}}
%\end
%\emline(81.25,55)(97.5,48)
\multiput(81.25,55)(.078125,-.03365385){208}{\line(1,0){.078125}}
%\end
%\emline(106.75,54)(94.25,49.5)
\multiput(106.75,54)(-.09328358,-.03358209){134}{\line(-1,0){.09328358}}
%\end
\put(65.25,3.75){\makebox(0,0)[cc]{Figure 2.4}}
\end{picture}
\end{center}

\begin{center}
%TeXCAD Picture [fig2.5.pic]. Options:
%\grade{\on}
%\emlines{\off}
%\epic{\off}
%\beziermacro{\on}
%\reduce{\on}
%\snapping{\off}
%\quality{8.00}
%\graddiff{0.01}
%\snapasp{1}
%\zoom{4.0000}
\unitlength 1mm % = 2.85pt
\linethickness{0.4pt}
\ifx\plotpoint\undefined\newsavebox{\plotpoint}\fi % GNUPLOT compatibility
\begin{picture}(102.6,61.94)(0,0)
%\emline(21.6,60.19)(10.6,58.94)
\multiput(21.6,60.19)(-.2894737,-.0328947){38}{\line(-1,0){.2894737}}
%\end
%\emline(10.6,58.94)(24.1,31.19)
\multiput(10.6,58.94)(.033665835,-.069201995){401}{\line(0,-1){.069201995}}
%\end
%\emline(24.1,31.19)(30.6,39.69)
\multiput(24.1,31.19)(.03367876,.04404145){193}{\line(0,1){.04404145}}
%\end
%\emline(30.6,39.69)(21.6,59.94)
\multiput(30.6,39.69)(-.033707865,.075842697){267}{\line(0,1){.075842697}}
%\end
\put(21.6,59.94){\line(1,0){.25}}
%\emline(19.95,39.5)(7.45,50.25)
\multiput(19.95,39.5)(-.039184953,.03369906){319}{\line(-1,0){.039184953}}
%\end
%\emline(7.45,50.25)(2.7,43.25)
\multiput(7.45,50.25)(-.03368794,-.04964539){141}{\line(0,-1){.04964539}}
%\end
%\emline(2.99,43.25)(21.24,36.75)
\multiput(2.99,43.25)(.09455959,-.03367876){193}{\line(1,0){.09455959}}
%\end
%\emline(21.6,35.94)(4.85,34.69)
\multiput(21.6,35.94)(-.4407895,-.0328947){38}{\line(-1,0){.4407895}}
%\end
%\emline(4.85,34.69)(8.35,26.94)
\multiput(4.85,34.69)(.03365385,-.07451923){104}{\line(0,-1){.07451923}}
%\end
%\emline(8.35,26.94)(22.35,34.19)
\multiput(8.35,26.94)(.06511628,.03372093){215}{\line(1,0){.06511628}}
%\end
%\emline(19.6,32.44)(19.1,24.19)
\multiput(19.6,32.44)(-.033333,-.55){15}{\line(0,-1){.55}}
%\end
%\emline(19.1,24.19)(18.85,23.44)
\multiput(19.1,24.19)(-.03125,-.09375){8}{\line(0,-1){.09375}}
%\end
%\emline(18.85,23.44)(24.6,31.69)
\multiput(18.85,23.44)(.03362573,.04824561){171}{\line(0,1){.04824561}}
%\end
%\emline(30.7,39.69)(28.95,21.19)
\multiput(30.7,39.69)(-.0336538,-.3557692){52}{\line(0,-1){.3557692}}
%\end
%\emline(19.1,23.19)(29.1,21.69)
\multiput(19.1,23.19)(.2222222,-.0333333){45}{\line(1,0){.2222222}}
%\end
%\emline(30.1,32.19)(38.6,26.69)
\multiput(30.1,32.19)(.05182927,-.03353659){164}{\line(1,0){.05182927}}
%\end
%\emline(38.6,26.69)(30.85,39.44)
\multiput(38.6,26.69)(-.03369565,.05543478){230}{\line(0,1){.05543478}}
%\end
%\emline(38.6,26.44)(47.1,30.94)
\multiput(38.6,26.44)(.06343284,.03358209){134}{\line(1,0){.06343284}}
%\end
%\emline(47.1,30.94)(28.6,44.19)
\multiput(47.1,30.94)(-.047073791,.033715013){393}{\line(-1,0){.047073791}}
%\end
%\emline(28.6,44.19)(47.1,40.19)
\multiput(28.6,44.19)(.15546218,-.03361345){119}{\line(1,0){.15546218}}
%\end
%\emline(47.1,40.19)(38.1,37.44)
\multiput(47.1,40.19)(-.1097561,-.0335366){82}{\line(-1,0){.1097561}}
%\end
%\emline(46.94,40.27)(48.19,48.52)
\multiput(46.94,40.27)(.0328947,.2171053){38}{\line(0,1){.2171053}}
%\end
%\emline(48.19,48.52)(27.44,47.02)
\multiput(48.19,48.52)(-.4611111,-.0333333){45}{\line(-1,0){.4611111}}
%\end
%\emline(27.4,47.14)(41.9,54.14)
\multiput(27.4,47.14)(.06971154,.03365385){208}{\line(1,0){.06971154}}
%\end
%\emline(41.76,54.14)(37.01,47.89)
\multiput(41.76,54.14)(-.03368794,-.04432624){141}{\line(0,-1){.04432624}}
%\end
%\emline(41.61,54.19)(34.36,60.44)
\multiput(41.61,54.19)(-.03897849,.03360215){186}{\line(-1,0){.03897849}}
%\end
%\emline(34.51,60.14)(25.76,50.14)
\multiput(34.51,60.14)(-.033653846,-.038461538){260}{\line(0,-1){.038461538}}
%\end
%\emline(15.6,47.94)(7.1,50.44)
\multiput(15.6,47.94)(-.1133333,.0333333){75}{\line(-1,0){.1133333}}
%\end
%\emline(4.85,34.69)(12.85,39.44)
\multiput(4.85,34.69)(.05673759,.03368794){141}{\line(1,0){.05673759}}
%\end
%\emline(21.35,60.19)(19.6,40.19)
\multiput(21.35,60.19)(-.0336538,-.3846154){52}{\line(0,-1){.3846154}}
%\end
%\dashline{1}(47.1,30.94)(50.35,28.44)
\multiput(47.03,30.87)(.043333,-.033333){15}{\line(1,0){.043333}}
\multiput(48.33,29.87)(.043333,-.033333){15}{\line(1,0){.043333}}
\multiput(49.63,28.87)(.043333,-.033333){15}{\line(1,0){.043333}}
%\end
%\dottedline(30.6,39.94)(19.6,41.19)
\multiput(30.53,39.87)(-.9167,.1042){13}{{\rule{.4pt}{.4pt}}}
%\end
%\emline(67.85,58.44)(80.35,18.94)
\multiput(67.85,58.44)(.033692722,-.106469003){371}{\line(0,-1){.106469003}}
%\end
%\emline(80.35,18.94)(85.1,35.69)
\multiput(80.35,18.94)(.03368794,.11879433){141}{\line(0,1){.11879433}}
%\end
%\emline(86.6,31.94)(99.35,26.94)
\multiput(86.6,31.94)(.08557047,-.03355705){149}{\line(1,0){.08557047}}
%\end
\put(99.35,26.94){\line(-1,1){31.5}}
%\emline(86.6,31.94)(84.1,33.19)
\multiput(86.6,31.94)(-.0657895,.0328947){38}{\line(-1,0){.0657895}}
%\end
%\emline(89.35,37.19)(102.6,44.69)
\multiput(89.35,37.19)(.05941704,.03363229){223}{\line(1,0){.05941704}}
%\end
\put(102.6,44.69){\line(-1,0){20.75}}
%\emline(88.35,44.44)(91.6,56.69)
\multiput(88.35,44.44)(.03350515,.12628866){97}{\line(0,1){.12628866}}
%\end
%\emline(91.6,56.69)(80.1,46.19)
\multiput(91.6,56.69)(-.036858974,-.033653846){312}{\line(-1,0){.036858974}}
%\end
%\emline(72.35,44.19)(55.85,43.44)
\multiput(72.35,44.19)(-.717391,-.032609){23}{\line(-1,0){.717391}}
%\end
%\emline(55.85,43.44)(74.85,36.19)
\multiput(55.85,43.44)(.08837209,-.03372093){215}{\line(1,0){.08837209}}
%\end
%\emline(71.6,37.19)(61.6,26.44)
\multiput(71.6,37.19)(-.033670034,-.036195286){297}{\line(0,-1){.036195286}}
%\end
%\emline(61.6,26.44)(76.1,31.94)
\multiput(61.6,26.44)(.08841463,.03353659){164}{\line(1,0){.08841463}}
%\end
%\emline(67.6,43.94)(73.6,39.19)
\multiput(67.6,43.94)(.04255319,-.03368794){141}{\line(1,0){.04255319}}
%\end
%\emline(73.6,39.19)(76.85,49.94)
\multiput(73.6,39.19)(.03350515,.11082474){97}{\line(0,1){.11082474}}
%\end
\put(76.85,49.94){\line(-1,0){.5}}
%\emline(71.35,48.19)(75.1,51.69)
\multiput(71.35,48.19)(.03605769,.03365385){104}{\line(1,0){.03605769}}
%\end
%\emline(75.1,51.69)(70.6,50.19)
\multiput(75.1,51.69)(-.1,-.0333333){45}{\line(-1,0){.1}}
%\end
%\emline(65.1,40.19)(65.85,44.19)
\multiput(65.1,40.19)(.032609,.173913){23}{\line(0,1){.173913}}
%\end
\put(65.85,44.19){\line(1,-4){1.25}}
%\emline(70.85,30.19)(68.1,33.69)
\multiput(70.85,30.19)(-.0335366,.0426829){82}{\line(0,1){.0426829}}
%\end
%\emline(68.1,33.69)(72.6,30.94)
\multiput(68.1,33.69)(.054878,-.0335366){82}{\line(1,0){.054878}}
%\end
%\emline(69.35,34.69)(75.6,34.19)
\multiput(69.35,34.69)(.416667,-.033333){15}{\line(1,0){.416667}}
%\end
%\emline(77.1,28.94)(85.1,35.94)
\multiput(77.1,28.94)(.03846154,.03365385){208}{\line(1,0){.03846154}}
%\end
%\emline(82.6,28.94)(77.85,26.94)
\multiput(82.6,28.94)(-.0791667,-.0333333){60}{\line(-1,0){.0791667}}
%\end
%\emline(77.85,26.94)(82.35,26.44)
\multiput(77.85,26.94)(.3,-.033333){15}{\line(1,0){.3}}
%\end
%\emline(91.1,35.44)(89.35,30.94)
\multiput(91.1,35.44)(-.0336538,-.0865385){52}{\line(0,-1){.0865385}}
%\end
%\emline(89.35,30.94)(92.1,34.19)
\multiput(89.35,30.94)(.0335366,.0396341){82}{\line(0,1){.0396341}}
%\end
%\emline(84.1,42.44)(92.35,39.19)
\multiput(84.1,42.44)(.08505155,-.03350515){97}{\line(1,0){.08505155}}
%\end
%\emline(90.35,44.44)(94.35,40.19)
\multiput(90.35,44.44)(.03361345,-.03571429){119}{\line(0,-1){.03571429}}
%\end
%\emline(94.35,40.19)(92.6,44.44)
\multiput(94.35,40.19)(-.0336538,.0817308){52}{\line(0,1){.0817308}}
%\end
%\emline(84.6,50.69)(89.35,48.44)
\multiput(84.6,50.69)(.0708955,-.0335821){67}{\line(1,0){.0708955}}
%\end
%\emline(89.35,48.44)(83.35,48.94)
\multiput(89.35,48.44)(-.4,.033333){15}{\line(-1,0){.4}}
%\end
%\emline(81.85,44.94)(89.1,46.44)
\multiput(81.85,44.94)(.1611111,.0333333){45}{\line(1,0){.1611111}}
%\end
%\dottedline(85.1,35.94)(74.35,38.69)
\multiput(85.03,35.87)(-.8958,.2292){13}{{\rule{.4pt}{.4pt}}}
%\end
%\qbezvec(48.1,53.19)(49.98,61.57)(58.35,55.44)
\put(58.35,55.44){\vector(3,-2){.07}}\qbezier(48.1,53.19)(49.98,61.57)(58.35,55.44)
%\end
%\emline(82.85,43.19)(86.1,31.69)
\multiput(82.85,43.19)(.03350515,-.1185567){97}{\line(0,-1){.1185567}}
%\end
%\dashline{1}(7.1,61.94)(28.6,43.69)
\multiput(7.03,61.87)(.037719,-.032018){19}{\line(1,0){.037719}}
\multiput(8.46,60.66)(.037719,-.032018){19}{\line(1,0){.037719}}
\multiput(9.9,59.44)(.037719,-.032018){19}{\line(1,0){.037719}}
\multiput(11.33,58.22)(.037719,-.032018){19}{\line(1,0){.037719}}
\multiput(12.76,57.01)(.037719,-.032018){19}{\line(1,0){.037719}}
\multiput(14.2,55.79)(.037719,-.032018){19}{\line(1,0){.037719}}
\multiput(15.63,54.57)(.037719,-.032018){19}{\line(1,0){.037719}}
\multiput(17.06,53.36)(.037719,-.032018){19}{\line(1,0){.037719}}
\multiput(18.5,52.14)(.037719,-.032018){19}{\line(1,0){.037719}}
\multiput(19.93,50.92)(.037719,-.032018){19}{\line(1,0){.037719}}
\multiput(21.36,49.71)(.037719,-.032018){19}{\line(1,0){.037719}}
\multiput(22.8,48.49)(.037719,-.032018){19}{\line(1,0){.037719}}
\multiput(24.23,47.27)(.037719,-.032018){19}{\line(1,0){.037719}}
\multiput(25.66,46.06)(.037719,-.032018){19}{\line(1,0){.037719}}
\multiput(27.1,44.84)(.037719,-.032018){19}{\line(1,0){.037719}}
%\end
\put(52.6,5.39){\makebox(0,0)[cc]{Figure 2.5}}
\put(28.72,14.3){\makebox(0,0)[cc]{(a)}}
\put(75.6,14.39){\makebox(0,0)[cc]{(b)}}
\end{picture}
\end{center}

\begin{center}
%TeXCAD Picture [fig2.6.pic]. Options:
%\grade{\on}
%\emlines{\off}
%\epic{\off}
%\beziermacro{\on}
%\reduce{\on}
%\snapping{\off}
%\quality{8.00}
%\graddiff{0.01}
%\snapasp{1}
%\zoom{4.0000}
\unitlength 1mm % = 2.85pt
\linethickness{0.4pt}
\ifx\plotpoint\undefined\newsavebox{\plotpoint}\fi % GNUPLOT compatibility
\begin{picture}(88.25,27.25)(0,0)
\put(4,23.25){\line(1,0){84.25}}
\put(88.25,12.5){\line(-1,0){84}}
%\dashline{1}(14.5,22.75)(37,12.25)
\multiput(14.43,22.68)(.07212,-.03365){12}{\line(1,0){.07212}}
\multiput(16.16,21.87)(.07212,-.03365){12}{\line(1,0){.07212}}
\multiput(17.89,21.06)(.07212,-.03365){12}{\line(1,0){.07212}}
\multiput(19.62,20.26)(.07212,-.03365){12}{\line(1,0){.07212}}
\multiput(21.35,19.45)(.07212,-.03365){12}{\line(1,0){.07212}}
\multiput(23.08,18.64)(.07212,-.03365){12}{\line(1,0){.07212}}
\multiput(24.81,17.83)(.07212,-.03365){12}{\line(1,0){.07212}}
\multiput(26.55,17.03)(.07212,-.03365){12}{\line(1,0){.07212}}
\multiput(28.28,16.22)(.07212,-.03365){12}{\line(1,0){.07212}}
\multiput(30.01,15.41)(.07212,-.03365){12}{\line(1,0){.07212}}
\multiput(31.74,14.6)(.07212,-.03365){12}{\line(1,0){.07212}}
\multiput(33.47,13.8)(.07212,-.03365){12}{\line(1,0){.07212}}
\multiput(35.2,12.99)(.07212,-.03365){12}{\line(1,0){.07212}}
%\end
%\dashline{1}(37,12.25)(67.5,23)
\multiput(36.93,12.18)(.09242,.03258){10}{\line(1,0){.09242}}
\multiput(38.78,12.83)(.09242,.03258){10}{\line(1,0){.09242}}
\multiput(40.63,13.48)(.09242,.03258){10}{\line(1,0){.09242}}
\multiput(42.48,14.13)(.09242,.03258){10}{\line(1,0){.09242}}
\multiput(44.32,14.79)(.09242,.03258){10}{\line(1,0){.09242}}
\multiput(46.17,15.44)(.09242,.03258){10}{\line(1,0){.09242}}
\multiput(48.02,16.09)(.09242,.03258){10}{\line(1,0){.09242}}
\multiput(49.87,16.74)(.09242,.03258){10}{\line(1,0){.09242}}
\multiput(51.72,17.39)(.09242,.03258){10}{\line(1,0){.09242}}
\multiput(53.57,18.04)(.09242,.03258){10}{\line(1,0){.09242}}
\multiput(55.41,18.69)(.09242,.03258){10}{\line(1,0){.09242}}
\multiput(57.26,19.35)(.09242,.03258){10}{\line(1,0){.09242}}
\multiput(59.11,20)(.09242,.03258){10}{\line(1,0){.09242}}
\multiput(60.96,20.65)(.09242,.03258){10}{\line(1,0){.09242}}
\multiput(62.81,21.3)(.09242,.03258){10}{\line(1,0){.09242}}
\multiput(64.66,21.95)(.09242,.03258){10}{\line(1,0){.09242}}
\multiput(66.51,22.6)(.09242,.03258){10}{\line(1,0){.09242}}
%\end
%\dashline{1}(67.5,23)(86.5,14.5)
\multiput(67.43,22.93)(.07197,-.0322){12}{\line(1,0){.07197}}
\multiput(69.16,22.16)(.07197,-.0322){12}{\line(1,0){.07197}}
\multiput(70.88,21.38)(.07197,-.0322){12}{\line(1,0){.07197}}
\multiput(72.61,20.61)(.07197,-.0322){12}{\line(1,0){.07197}}
\multiput(74.34,19.84)(.07197,-.0322){12}{\line(1,0){.07197}}
\multiput(76.07,19.07)(.07197,-.0322){12}{\line(1,0){.07197}}
\multiput(77.79,18.29)(.07197,-.0322){12}{\line(1,0){.07197}}
\multiput(79.52,17.52)(.07197,-.0322){12}{\line(1,0){.07197}}
\multiput(81.25,16.75)(.07197,-.0322){12}{\line(1,0){.07197}}
\multiput(82.98,15.98)(.07197,-.0322){12}{\line(1,0){.07197}}
\multiput(84.7,15.2)(.07197,-.0322){12}{\line(1,0){.07197}}
%\end
%\dashline{1}(14,22.75)(16,19)
\multiput(13.93,22.68)(.03333,-.0625){12}{\line(0,-1){.0625}}
\multiput(14.73,21.18)(.03333,-.0625){12}{\line(0,-1){.0625}}
\multiput(15.53,19.68)(.03333,-.0625){12}{\line(0,-1){.0625}}
%\end
%\dashline{1}(14.25,22.75)(18,19)
\multiput(14.18,22.68)(.032895,-.032895){19}{\line(0,-1){.032895}}
\multiput(15.43,21.43)(.032895,-.032895){19}{\line(1,0){.032895}}
\multiput(16.68,20.18)(.032895,-.032895){19}{\line(0,-1){.032895}}
%\end
%\dashline{1}(37.25,12.25)(38.25,17.75)
\multiput(37.18,12.18)(.02857,.15714){5}{\line(0,1){.15714}}
\multiput(37.47,13.75)(.02857,.15714){5}{\line(0,1){.15714}}
\multiput(37.75,15.32)(.02857,.15714){5}{\line(0,1){.15714}}
\multiput(38.04,16.89)(.02857,.15714){5}{\line(0,1){.15714}}
%\end
%\dashline{1}(37.5,12.75)(42,15)
\multiput(37.43,12.68)(.0625,.03125){12}{\line(1,0){.0625}}
\multiput(38.93,13.43)(.0625,.03125){12}{\line(1,0){.0625}}
\multiput(40.43,14.18)(.0625,.03125){12}{\line(1,0){.0625}}
%\end
%\dashline{1}(37.75,13)(40,16.25)
\multiput(37.68,12.93)(.032143,.046429){14}{\line(0,1){.046429}}
\multiput(38.58,14.23)(.032143,.046429){14}{\line(0,1){.046429}}
\multiput(39.48,15.53)(.032143,.046429){14}{\line(0,1){.046429}}
%\end
%\dashline{1}(67.25,23)(73,17)
\multiput(67.18,22.93)(.031944,-.033333){18}{\line(0,-1){.033333}}
\multiput(68.33,21.73)(.031944,-.033333){18}{\line(0,-1){.033333}}
\multiput(69.48,20.53)(.031944,-.033333){18}{\line(0,-1){.033333}}
\multiput(70.63,19.33)(.031944,-.033333){18}{\line(0,-1){.033333}}
\multiput(71.78,18.13)(.031944,-.033333){18}{\line(0,-1){.033333}}
%\end
%\dashline{1}(67.75,23)(72,19.5)
\multiput(67.68,22.93)(.040476,-.033333){15}{\line(1,0){.040476}}
\multiput(68.89,21.93)(.040476,-.033333){15}{\line(1,0){.040476}}
\multiput(70.11,20.93)(.040476,-.033333){15}{\line(1,0){.040476}}
\multiput(71.32,19.93)(.040476,-.033333){15}{\line(1,0){.040476}}
%\end
\put(2.25,24.75){\vector(1,0){4.5}}
\put(7.5,9.75){\vector(1,0){4.25}}
\put(21.5,21.5){\makebox(0,0)[cc]{$u_{k}$}}
\put(27.5,14.75){\makebox(0,0)[cc]{$u_{k}$}}
\put(47.75,14.75){\makebox(0,0)[cc]{$u_{k+1}$}}
\put(9,27.25){\makebox(0,0)[cc]{Fold down $n$ times}}
\put(35.25,7.75){\makebox(0,0)[cc]{Fold up $n$ times}}
\put(44.75,2.75){\makebox(0,0)[cc]{Figure 2.6}}
\end{picture}

\end{center}

We are now in the lovely position of having many answerable questions. Here are just a few of the possibilities:
\begin{enumerate}[\hspace{1cm}(1)]
  \item Can we prove in the above example, or in general, that the actual angles on the tape really converge to the putative angle we originally sought?
  \item What happens if you consider general folding procedures with other periods, such as $D^{3} U^{3}$, $D^{4} U^{2}$, or $D^{3} U^{1} D^{1} U^{3} D^{1} U^{1}$? (The period is determined when the exponents repeat, so these examples have periods $1$, $2$, and $3$, respectively.)
  \item How many different regular convex polygons is it possible to fold with a $2$-period $D^{m} U^{n}$ folding procedure? Which ones are they? How do we recognize them? How do we determine the appropriate $m$ and $n$ in order to get them?
  \item Is there an easiest folding procedure that will produce a regular convex $b$-gon for every odd number $b > 3$? If so, how do we find it?
  \item How many of the regular star polygons is it possible to construct by folding paper?
\end{enumerate}

The answers to all of these questions (and, in fact, many more) are \emph{known}! In the remainder of this section, we will first indicate how to answer questions 1 and 2, and then illustrate a method for dealing with question 4. Question 3 is discussed in [4], and we show in the last section of this paper that the answer to question 5 is the best possible; that is, in theory, we can construct by folding paper, \emph{all} regular star polygons. Furthermore (in answer to question 4), we indicate how the procedures we use give the best possible method of constructing the regular star polygons.

Let us begin to answer these questions. It turns out that the answers to questions 1 and 2 can be obtained by a straightforward use of the following:

\begin{description}

  \item[Lemma 2.1:] Let the sequence $\{ x_{k} \}$, $k = 0,1,2,\cdots$, be generated by the recurrence relation
    \begin{equation}
    x_{k} + ax_{k+1} = b,\ k = 0,1,2,\cdots  \tag{2.1}
    \end{equation}
    So if $|a| > 1$, then $x_{k} \to \frac{b}{(1+a)}$.

  \item[Proof:] Set $x_{k} = \frac{b}{(l+a)} + y_{k}$. Then $y_{k} + ay_{k+1} = 0$. It follows that $y_{k} = \left( \frac{-l}{a} \right)^{k} y_{0}$. If $|a| > 1$, $\left( \frac{-1}{a} \right)^{k} \to 0$, so that $y_{k} \to 0$ as $k \to \infty$. Hence $x_{k} \to \frac{b}{(l+a)}$.

\end{description}

This result may also be viewed as a special case of the well-known contraction mapping principle (see [11]). We point out that it is significant that neither the convergence nor the limit depends on the initial value $x_{0}$. And, as we will soon demonstrate, the proof of the lemma tells us that the convergence of our folding procedure is rapid, since in all cases $|a|$ will be a power of $2$.

Let us begin our discussion of convergence by looking at the general $1$-period folding procedure $D^{n} U^{n}$. A typical portion of the tape would appear as illustrated in Figure 2.6. Then, if the folding process had been started with an arbitrary angle $u_{0}$, we would have
\begin{equation}
u_{k} + 2^{n} u_{k+1} = \pi,\ k = 0,1,2,\cdots  \tag{2.2}
\end{equation}
with $|u_{0}| < \frac{\pi}{2}$. Equation (2.2) is of the form (2.1), so it follows immediately from Lemma 2.1 that
\begin{equation}
u_{k} \to \frac{\pi}{2^{n}+1} \text{ as } k \to \infty  \tag{2.3}
\end{equation}

Furthermore, we can readily see that if the original fold differed from the putative angle of $\frac{\pi}{(2^{n}+l)}$ by an error $E_{0}$, then the error at the $k$th stage of the folding process would be given by
\begin{equation*}
E_{k} = \frac{E_{0}}{2^{nk}}
\end{equation*}

We see from (2.3) that the $D^{n} U^{n}$ folding procedure produces tape from which we may construct $n$ regular convex $(2^{n}+1)$-gons. In fact, because the measure of the smallest acute angle at the bottom and the top of this tape is the same, we may construct the $(2^{n}+1)$-gons by special methods which are not as comprehensive as our FAT-algorithm. Figure 2.7 shows how $D^{2} U^{2}$ tape may be folded along just the short lines of the folded tape to form the outline of a small regular convex pentagon (a); and along just the long lines of the folded tape to form the outline of a slightly larger regular convex pentagon (b); and finally, in (c) we show the regular convex pentagon formed by executing the FAT-algorithm. 

\begin{center}
%TeXCAD Options
%\grade{\on}
%\emlines{\off}
%\epic{\off}
%\beziermacro{\on}
%\reduce{\on}
%\snapping{\off}
%\quality{8.00}
%\graddiff{0.01}
%\snapasp{1}
%\zoom{4.0000}
\unitlength .6mm % = 2.85pt
\linethickness{0.4pt}
\ifx\plotpoint\undefined\newsavebox{\plotpoint}\fi % GNUPLOT compatibility
\begin{picture}(190,100)(10,25)
\put(17.5,117.75){\line(1,0){135}}
\put(154,104){\line(-1,0){136.25}}
%\emline(40.75,98)(25,90.5)
\multiput(40.75,98)(-.0706278,-.03363229){223}{\line(-1,0){.0706278}}
%\end
\put(25,90.5){\line(1,0){27}}
%\emline(52,90.5)(40.75,98.25)
\multiput(52,90.5)(-.04891304,.03369565){230}{\line(-1,0){.04891304}}
%\end
%\emline(52,90.75)(48,76.25)
\multiput(52,90.75)(-.03361345,-.12184874){119}{\line(0,-1){.12184874}}
%\end
\put(48,76.25){\line(-1,0){17.75}}
%\emline(30.25,76.25)(25,90.5)
\multiput(30.25,76.25)(-.03365385,.09134615){156}{\line(0,1){.09134615}}
%\end
%\emline(80.5,91.75)(58.5,78.5)
\multiput(80.5,91.75)(-.055979644,-.033715013){393}{\line(-1,0){.055979644}}
%\end
%\emline(86.25,78.75)(80.5,91.5)
\multiput(86.25,78.75)(-.03362573,.0745614){171}{\line(0,1){.0745614}}
%\end
%\emline(80.5,91.5)(101.75,79.25)
\multiput(80.5,91.5)(.058379121,-.033653846){364}{\line(1,0){.058379121}}
%\end
%\emline(101.75,79.25)(81.5,65.5)
\multiput(101.75,79.25)(-.049632353,-.03370098){408}{\line(-1,0){.049632353}}
%\end
%\emline(81.5,65.5)(73.5,70.25)
\multiput(81.5,65.5)(-.05673759,.03368794){141}{\line(-1,0){.05673759}}
%\end
%\emline(73.5,70.25)(75.5,78.75)
\multiput(73.5,70.25)(.0333333,.1416667){60}{\line(0,1){.1416667}}
%\end
%\emline(73.25,70)(69.5,56.75)
\multiput(73.25,70)(-.03348214,-.11830357){112}{\line(0,-1){.11830357}}
%\end
%\emline(69.5,56.75)(58.5,78)
\multiput(69.5,56.75)(-.033639144,.064984709){327}{\line(0,1){.064984709}}
%\end
%\emline(58.25,78.75)(59.5,76.25)
\multiput(58.25,78.75)(.0328947,-.0657895){38}{\line(0,-1){.0657895}}
%\end
\put(69.25,56.5){\line(1,0){23.75}}
%\emline(93,56.5)(102,79.25)
\multiput(93,56.5)(.033707865,.085205993){267}{\line(0,1){.085205993}}
%\end
\put(86.5,78.5){\line(1,-2){3.5}} \put(86.25,78.75){\line(-1,0){28}}
%\emline(81.5,65.5)(93.25,56.5)
\multiput(81.5,65.5)(.044007491,-.033707865){267}{\line(1,0){.044007491}}
%\end
%\dottedline(18.75,113.25)(27,104)
\multiput(18.68,113.18)(.5893,-.6607){15}{{\rule{.4pt}{.4pt}}}
%\end
%\dottedline(27,104)(31.75,117.5)
\multiput(26.93,103.93)(.3167,.9){16}{{\rule{.4pt}{.4pt}}}
%\end
%\dottedline(27.25,104)(48.75,117.75)
\multiput(27.18,103.93)(.8269,.5288){27}{{\rule{.4pt}{.4pt}}}
%\end
%\dottedline(48.75,117.75)(71.25,104.25)
\multiput(48.68,117.68)(.8333,-.5){28}{{\rule{.4pt}{.4pt}}}
%\end
%\dottedline(71.5,104)(93,117.5)
\multiput(71.43,103.93)(.8269,.5192){27}{{\rule{.4pt}{.4pt}}}
%\end
%\dottedline(93,117.5)(115.5,103.75)
\multiput(92.93,117.43)(.8333,-.5093){28}{{\rule{.4pt}{.4pt}}}
%\end
%\dottedline(49,117.5)(55,104)
\multiput(48.93,117.43)(.375,-.8438){17}{{\rule{.4pt}{.4pt}}}
%\end
%\dottedline(92.75,117)(98,104)
\multiput(92.68,116.93)(.35,-.8667){16}{{\rule{.4pt}{.4pt}}}
%\end
%\dottedline(71.5,103.75)(75.25,117.25)
\multiput(71.43,103.68)(.25,.9){16}{{\rule{.4pt}{.4pt}}}
%\end
%\dottedline(115.5,104.25)(137,117.75)
\multiput(115.43,104.18)(.8269,.5192){27}{{\rule{.4pt}{.4pt}}}
%\end
%\dottedline(137,117.75)(159.5,104)
\multiput(136.93,117.68)(.8333,-.5093){28}{{\rule{.4pt}{.4pt}}}
%\end
%\dottedline(136.75,117.25)(142,104.25)
\multiput(136.68,117.18)(.35,-.8667){16}{{\rule{.4pt}{.4pt}}}
%\end
%\dottedline(115.5,104)(119.25,117.5)
\multiput(115.43,103.93)(.25,.9){16}{{\rule{.4pt}{.4pt}}}
%\end
\put(152.75,117.75){\line(1,0){38.75}}
\put(154.5,104){\line(1,0){36.25}}
\put(190.75,104){\line(1,0){6.75}}
\put(191.5,117.75){\line(1,0){4.25}}
%\dottedline(159.25,104)(178.75,117.5)
\multiput(159.18,103.93)(.78,.54){26}{{\rule{.4pt}{.4pt}}}
%\end
%\dottedline(178.75,117.5)(196.25,108.75)
\multiput(178.68,117.43)(.875,-.4375){21}{{\rule{.4pt}{.4pt}}}
%\end
%\dottedline(178.75,117.5)(186.5,104.5)
\multiput(178.68,117.43)(.4844,-.8125){17}{{\rule{.4pt}{.4pt}}}
%\end
%\dottedline(159.25,104.25)(164.75,117.75)
\multiput(159.18,104.18)(.3667,.9){16}{{\rule{.4pt}{.4pt}}}
%\end
%\emline(156.75,99)(121.5,77.25)
\multiput(156.75,99)(-.054651163,-.03372093){645}{\line(-1,0){.054651163}}
%\end
%\emline(121.5,77.25)(136.75,41.75)
\multiput(121.5,77.25)(.033738938,-.078539823){452}{\line(0,-1){.078539823}}
%\end
\put(136.75,41.75){\line(1,0){43.25}}
%\emline(180,41.75)(192,76.5)
\multiput(180,41.75)(.033707865,.09761236){356}{\line(0,1){.09761236}}
%\end
%\emline(192,76.5)(157.25,99)
\multiput(192,76.5)(-.0520989505,.0337331334){667}{\line(-1,0){.0520989505}}
%\end
%\emline(157,99)(165.75,76.5)
\multiput(157,99)(.033653846,-.086538462){260}{\line(0,-1){.086538462}}
%\end
%\emline(165.75,76.5)(178.75,68.5)
\multiput(165.75,76.5)(.05462185,-.03361345){238}{\line(1,0){.05462185}}
%\end
%\emline(178.75,68.5)(192,76.5)
\multiput(178.75,68.5)(.05567227,.03361345){238}{\line(1,0){.05567227}}
%\end
%\emline(179,68.5)(169.5,62.25)
\multiput(179,68.5)(-.05107527,-.03360215){186}{\line(-1,0){.05107527}}
%\end
\put(186.75,61.75){\line(-1,0){17.75}}
%\emline(169,61.75)(171.25,63.25)
\multiput(169,61.75)(.05,.0333333){45}{\line(1,0){.05}}
%\end
%\emline(168.75,61.5)(166.25,50.75)
\multiput(168.75,61.5)(-.0333333,-.1433333){75}{\line(0,-1){.1433333}}
%\end
%\emline(166.25,50.75)(180,41.75)
\multiput(166.25,50.75)(.051498127,-.033707865){267}{\line(1,0){.051498127}}
%\end
%\emline(166.5,50.5)(157.75,56.25)
\multiput(166.5,50.5)(-.05116959,.03362573){171}{\line(-1,0){.05116959}}
%\end
%\emline(157.75,56.25)(164,42)
\multiput(157.75,56.25)(.03360215,-.0766129){186}{\line(0,-1){.0766129}}
%\end
\put(157.5,56){\line(-1,0){14.25}}
%\emline(143.25,56)(137,42)
\multiput(143.25,56)(-.03360215,-.07526882){186}{\line(0,-1){.07526882}}
%\end
\put(143.25,55.5){\line(1,2){4}}
%\emline(147.25,63.5)(140,77)
\multiput(147.25,63.5)(-.03372093,.0627907){215}{\line(0,1){.0627907}}
%\end
\put(140,77){\line(-1,0){18.5}} \put(140.25,77){\line(1,0){9.75}}
%\emline(162.75,84.25)(150,77)
\multiput(162.75,84.25)(-.05930233,-.03372093){215}{\line(-1,0){.05930233}}
%\end
%\emline(150,77.25)(135,85.75)
\multiput(150,77.25)(-.05952381,.033730159){252}{\line(-1,0){.05952381}}
%\end
%\emline(147.25,63.5)(131,55.75)
\multiput(147.25,63.5)(-.07065217,-.03369565){230}{\line(-1,0){.07065217}}
%\end
%\emline(165.75,76.25)(171.5,90)
\multiput(165.75,76.25)(.03362573,.08040936){171}{\line(0,1){.08040936}}
%\end
%\dottedline(150.25,77.5)(157,99.25)
\multiput(150.18,77.43)(.2935,.9457){24}{{\rule{.4pt}{.4pt}}}
%\end
%\dottedline(150.25,77.25)(143.5,90.5)
\multiput(150.18,77.18)(-.4219,.8281){17}{{\rule{.4pt}{.4pt}}}
%\end
%\dottedline(147,63.75)(121.5,77)
\multiput(146.93,63.68)(-.85,.4417){31}{{\rule{.4pt}{.4pt}}}
%\end
%\dottedline(146.75,63.75)(127.5,63.5)
\multiput(146.68,63.68)(-.9625,-.0125){21}{{\rule{.4pt}{.4pt}}}
%\end
%\dottedline(165.5,76.5)(179,84.75)
\multiput(165.43,76.43)(.7941,.4853){18}{{\rule{.4pt}{.4pt}}}
%\end
%\dottedline(165.75,76.5)(191.75,76.75)
\multiput(165.68,76.43)(.963,.0093){28}{{\rule{.4pt}{.4pt}}}
%\end
%\dottedline(169.25,61.75)(184.5,55.25)
\multiput(169.18,61.68)(.8472,-.3611){19}{{\rule{.4pt}{.4pt}}}
%\end
%\dottedline(169.5,61.25)(180,42)
\multiput(169.43,61.18)(.4565,-.837){24}{{\rule{.4pt}{.4pt}}}
%\end
%\dottedline(158,56)(137.25,41.75)
\multiput(157.93,55.93)(-.7981,-.5481){27}{{\rule{.4pt}{.4pt}}}
%\end
%\dottedline(157.75,55.5)(154,41.75)
\multiput(157.68,55.43)(-.25,-.9167){16}{{\rule{.4pt}{.4pt}}}
%\end
%\dottedline(25,90.25)(48.25,76.25)
\multiput(24.93,90.18)(.8017,-.4828){30}{{\rule{.4pt}{.4pt}}}
%\end
%\dottedline(41,98)(44.5,90.75)
\multiput(40.93,97.93)(.3889,-.8056){10}{{\rule{.4pt}{.4pt}}}
%\end
\put(100,100){\makebox(0,0)[cc]{The $D^{2} U^{2}$ folded tape}}
\put(38.25,72.5){\makebox(0,0)[cc]{(a)}}
\put(80,51.5){\makebox(0,0)[cc]{(b)}}
\put(158.25,38.25){\makebox(0,0)[cc]{(c)}}
\put(102.25,30.75){\makebox(0,0)[cc]{Figure 2.7}}
\end{picture}
\end{center}

We feel certain that both the ancient Greeks and Gauss would have appreciated the fact that when $n = 1,2,3,4,8,16,32$, we obtain, in theory, easy constructions for regular $3$, $5$, $9$, $17$, $257$, $65537$, and $(2^{32}+1)$-gons, respectively.

The Greeks first raised (and solved) the question of how to construct (with straight-edge and compass) regular $N$-gons of the form $2^{n} N_{0}$, where $N_{0} = 3$, $5$, or $15$, and $n \geq 0$; while Guass (1777-1855) proved that regular $N$-gons are constructible with a straight-edge and compass if and only if $N$ is of the form $2^{n} N_{0}$ with $N_{0}$ a product of distinct Fermat primes; that is, a product of primes of the form $2^{2^{m}} + 1$.

We would argue that in practice, the approximations obtained by folding paper are more accurate than the real-world constructions with a straight-edge and compass -- for the latter are only perfect in the mind. In both cases, what you actually get is a function of human skill, but the skill element in paper-folding is much less severe.

Now, suppose we look at the general $2$-period folding procedure, $D^{m} U^{n}$. In this case, a typical portion of the tape would appear as shown in Figure 2.8. 

\begin{center}
%TeXCAD Picture [fig2.8.pic]. Options:
%\grade{\on}
%\emlines{\off}
%\epic{\off}
%\beziermacro{\on}
%\reduce{\on}
%\snapping{\off}
%\quality{8.00}
%\graddiff{0.01}
%\snapasp{1}
%\zoom{4.0000}
\unitlength 1mm % = 2.85pt
\linethickness{0.4pt}
\ifx\plotpoint\undefined\newsavebox{\plotpoint}\fi % GNUPLOT compatibility
\begin{picture}(99,31)(0,0)
\put(10,26){\line(1,0){88.75}}
\put(99,14.5){\line(-1,0){88}}
%\dottedline(20,25.75)(38.75,14.5)
\multiput(19.93,25.68)(.8152,-.4891){24}{{\rule{.4pt}{.4pt}}}
%\end
%\dottedline(38.75,14.5)(63.25,26)
\multiput(38.68,14.43)(.875,.4107){29}{{\rule{.4pt}{.4pt}}}
%\end
%\dottedline(63.25,26)(88.75,14.25)
\multiput(63.18,25.93)(.8793,-.4052){30}{{\rule{.4pt}{.4pt}}}
%\end
%\dottedline(64,25.75)(65.5,19.75)
\multiput(63.93,25.68)(.1875,-.75){9}{{\rule{.4pt}{.4pt}}}
%\end
%\dottedline(64.75,25)(67,21.75)
\multiput(64.68,24.93)(.45,-.65){6}{{\rule{.4pt}{.4pt}}}
%\end
%\dottedline(38.5,14.25)(41.75,20.75)
\multiput(38.43,14.18)(.4063,.8125){9}{{\rule{.4pt}{.4pt}}}
%\end
%\dottedline(39,14.75)(45.5,19.5)
\multiput(38.93,14.68)(.7222,.5278){10}{{\rule{.4pt}{.4pt}}}
%\end
%\dottedline(20,25.75)(22.75,20.5)
\multiput(19.93,25.68)(.3929,-.75){8}{{\rule{.4pt}{.4pt}}}
%\end
%\dottedline(20.5,25.25)(25,20.75)
\multiput(20.43,25.18)(.5625,-.5625){9}{{\rule{.4pt}{.4pt}}}
%\end
\put(31.75,16.5){\makebox(0,0)[cc]{$u_{k}$}}
\put(48.5,16.75){\makebox(0,0)[cc]{$v_{k}$}}
\put(52.75,23.75){\makebox(0,0)[cc]{$v_{k}$}}
\put(73,23.75){\makebox(0,0)[cc]{$u_{k+1}$}}
\put(11.75,29.25){\makebox(0,0)[cc]{Fold down $m$ times}}
\put(56,10.25){\makebox(0,0)[cc]{Fold up $m$ times}}
\put(50.25,2.75){\makebox(0,0)[cc]{Figure 2.8}}
\end{picture}
\end{center}

If the folding process had been started with an arbitrary acute angle $u_{0}$ at the top of the tape, we would have
\begin{align*}
u_{k} + 2^{n} v_{k} & = \pi \\
v_{k} + 2^{m} u_{k+1} & = \pi
\end{align*}
and hence, it follows that
\begin{equation}
u_{k} - 2^{m+n} u_{k+1} = \pi (1 - 2^{n}) \text{ with } u_{0} < \frac{\pi}{2},\ k = 0,1,2,\cdots  \tag{2.4}
\end{equation}
and thus, again using Lemma 2.1, we see that
\begin{equation}
u_{k} \to \frac{2^{n}-1}{2^{n+m}-1} \pi \text{ as } k \to \infty  \tag{2.5}
\end{equation}
so that the FAT-algorithm will produce from this tape a star $\left\{ \frac{b}{a} \right\}$-gon, with $b = 2^{m+n}-1$, $a = 2^{n}-1$. By symmetry, we can infer that
\begin{equation}
v_{k} \to \frac{2^{m}-1}{2^{n+m}-1} \pi \text{ as } k \to \infty  \tag{2.6}
\end{equation}

Furthermore, if we assume an initial error of $E_{0}$, then we know that the error at the $k$th stage (when the folding $D^{m} U^{n}$ has been done exactly $k$ times) will be given by $E_{k} = \frac{E_{0}}{2^{(n+m)k}}$. Hence, we see that in the case of our $D^{2} U^{1}$ folding, any initial error $E_{0}$ is reduced by a factor of $\frac{1}{2^{3}}$ between consecutive appearances of any one of the putative angles.

It should now be clear why we advised throwing away the first part of the tape -- but likewise, it should also be clear that it is never necessary to throw away very much of the tape. In practice, convergence is very rapid indeed, and if one made it a rule of thumb to always throw away the first 20 fold lines on the tape for any iterative folding procedure, it would turn out to be a very conservative rule ($\frac{(\pi/2)}{2^{20}}$ is a very small error!).

The above technique can be used to prove the convergence for any given folding sequence of general period, say $r$. This is fine, once we have identified the folding sequence, but the question that is more likely to arise for the paper-folder is: How do we know which sequence to fold in order to produce a particular regular star polygon with the FAT-algorithm? Let us demonstrate with an example, which is sufficiently general that it will apply to any regular star $\left\{ \frac{b}{a} \right\}$-gon where $a,b$ are mutually prime odd integers with $a < \left( \frac{b}{2} \right)$.

Suppose we desperately want to construct a regular star $\left\{ \frac{11}{3} \right\}$-gon. Then of course, $b = 11$, $a = 3$, and it is clear that if we begin with an angle of $\left( \frac{3}{ll} \right) \pi$ at $A_{0}$ and adhere to the Folding Rules, we will obtain a strip of tape with the angles and creases (indicated by dotted lines) shown in Figure 2.9. 

We could denote this folding process as $D^{1} U^{3} D^{1} U^{1} D^{3} U^{1}$, interpreted in the obvious way on the tape -- that is, the first exponent ``$1$'' refers to the one bisection (producing a line in a downward direction) at the vertices $A_{10n}$ (for $n = 0,1,2,\cdots$) on the top of the tape; similarly, the ``$3$'' refers to the $3$ bisections (producing creases in an upward direction) made at the bottom of the tape through the vertices $A_{10n+1}$; etc. However, since the folding process is \emph{duplicated} halfway through, we can abbreviate the notation and write simply $\{ 1,3,1 \}$, with the understanding that we alternate folding from the top and bottom of the tape as described, with the number of bisections at each vertex running, in order, through the values $1,3,1,\cdots$. We call this a \emph{primary folding procedure of period $3$}.

\begin{center}
%TeXCAD Options
%\grade{\on}
%\emlines{\off}
%\epic{\off}
%\beziermacro{\on}
%\reduce{\on}
%\snapping{\off}
%\quality{8.00}
%\graddiff{0.01}
%\snapasp{1}
%\zoom{4.0000}
\unitlength .78mm % = 2.85pt
\linethickness{0.4pt}
\ifx\plotpoint\undefined\newsavebox{\plotpoint}\fi % GNUPLOT compatibility
\begin{picture}(155,70)(5,70)
\thicklines 
\put(29.25,131.75){\line(1,0){102}}
\put(29.25,120.25){\line(1,0){102}}
\put(14.75,106.25){\line(1,0){141.5}}
\put(13.75,93){\line(1,0){143}}
%\dashline{1}(17,93)(25.75,106.25)
\multiput(16.93,92.93)(.032169,.048713){16}{\line(0,1){.048713}}
\multiput(17.96,94.49)(.032169,.048713){16}{\line(0,1){.048713}}
\multiput(18.99,96.05)(.032169,.048713){16}{\line(0,1){.048713}}
\multiput(20.02,97.61)(.032169,.048713){16}{\line(0,1){.048713}}
\multiput(21.05,99.17)(.032169,.048713){16}{\line(0,1){.048713}}
\multiput(22.08,100.72)(.032169,.048713){16}{\line(0,1){.048713}}
\multiput(23.11,102.28)(.032169,.048713){16}{\line(0,1){.048713}}
\multiput(24.14,103.84)(.032169,.048713){16}{\line(0,1){.048713}}
\multiput(25.17,105.4)(.032169,.048713){16}{\line(0,1){.048713}}
%\end
%\dashline{1}(25.75,106.25)(35.75,93)
\multiput(25.68,106.18)(.03268,-.043301){17}{\line(0,-1){.043301}}
\multiput(26.79,104.71)(.03268,-.043301){17}{\line(0,-1){.043301}}
\multiput(27.9,103.24)(.03268,-.043301){17}{\line(0,-1){.043301}}
\multiput(29.01,101.76)(.03268,-.043301){17}{\line(0,-1){.043301}}
\multiput(30.12,100.29)(.03268,-.043301){17}{\line(0,-1){.043301}}
\multiput(31.24,98.82)(.03268,-.043301){17}{\line(0,-1){.043301}}
\multiput(32.35,97.35)(.03268,-.043301){17}{\line(0,-1){.043301}}
\multiput(33.46,95.87)(.03268,-.043301){17}{\line(0,-1){.043301}}
\multiput(34.57,94.4)(.03268,-.043301){17}{\line(0,-1){.043301}}
%\end
%\dashline{1}(35.75,93)(44,106)
\multiput(35.68,92.93)(.032353,.05098){15}{\line(0,1){.05098}}
\multiput(36.65,94.46)(.032353,.05098){15}{\line(0,1){.05098}}
\multiput(37.62,95.99)(.032353,.05098){15}{\line(0,1){.05098}}
\multiput(38.59,97.52)(.032353,.05098){15}{\line(0,1){.05098}}
\multiput(39.56,99.05)(.032353,.05098){15}{\line(0,1){.05098}}
\multiput(40.53,100.58)(.032353,.05098){15}{\line(0,1){.05098}}
\multiput(41.5,102.11)(.032353,.05098){15}{\line(0,1){.05098}}
\multiput(42.47,103.64)(.032353,.05098){15}{\line(0,1){.05098}}
\multiput(43.44,105.17)(.032353,.05098){15}{\line(0,1){.05098}}
%\end
%\dashline{1}(59.5,106.25)(36.25,93)
\multiput(59.43,106.18)(-.055357,-.031548){15}{\line(-1,0){.055357}}
\multiput(57.77,105.23)(-.055357,-.031548){15}{\line(-1,0){.055357}}
\multiput(56.11,104.29)(-.055357,-.031548){15}{\line(-1,0){.055357}}
\multiput(54.45,103.34)(-.055357,-.031548){15}{\line(-1,0){.055357}}
\multiput(52.79,102.39)(-.055357,-.031548){15}{\line(-1,0){.055357}}
\multiput(51.13,101.45)(-.055357,-.031548){15}{\line(-1,0){.055357}}
\multiput(49.47,100.5)(-.055357,-.031548){15}{\line(-1,0){.055357}}
\multiput(47.8,99.55)(-.055357,-.031548){15}{\line(-1,0){.055357}}
\multiput(46.14,98.61)(-.055357,-.031548){15}{\line(-1,0){.055357}}
\multiput(44.48,97.66)(-.055357,-.031548){15}{\line(-1,0){.055357}}
\multiput(42.82,96.72)(-.055357,-.031548){15}{\line(-1,0){.055357}}
\multiput(41.16,95.77)(-.055357,-.031548){15}{\line(-1,0){.055357}}
\multiput(39.5,94.82)(-.055357,-.031548){15}{\line(-1,0){.055357}}
\multiput(37.84,93.88)(-.055357,-.031548){15}{\line(-1,0){.055357}}
%\end
%\dashline{1}(36.25,93)(87,106)
\multiput(36.18,92.93)(.11969,.03066){8}{\line(1,0){.11969}}
\multiput(38.09,93.42)(.11969,.03066){8}{\line(1,0){.11969}}
\multiput(40.01,93.91)(.11969,.03066){8}{\line(1,0){.11969}}
\multiput(41.92,94.4)(.11969,.03066){8}{\line(1,0){.11969}}
\multiput(43.84,94.89)(.11969,.03066){8}{\line(1,0){.11969}}
\multiput(45.76,95.38)(.11969,.03066){8}{\line(1,0){.11969}}
\multiput(47.67,95.87)(.11969,.03066){8}{\line(1,0){.11969}}
\multiput(49.59,96.36)(.11969,.03066){8}{\line(1,0){.11969}}
\multiput(51.5,96.85)(.11969,.03066){8}{\line(1,0){.11969}}
\multiput(53.42,97.34)(.11969,.03066){8}{\line(1,0){.11969}}
\multiput(55.33,97.84)(.11969,.03066){8}{\line(1,0){.11969}}
\multiput(57.25,98.33)(.11969,.03066){8}{\line(1,0){.11969}}
\multiput(59.16,98.82)(.11969,.03066){8}{\line(1,0){.11969}}
\multiput(61.08,99.31)(.11969,.03066){8}{\line(1,0){.11969}}
\multiput(62.99,99.8)(.11969,.03066){8}{\line(1,0){.11969}}
\multiput(64.91,100.29)(.11969,.03066){8}{\line(1,0){.11969}}
\multiput(66.82,100.78)(.11969,.03066){8}{\line(1,0){.11969}}
\multiput(68.74,101.27)(.11969,.03066){8}{\line(1,0){.11969}}
\multiput(70.65,101.76)(.11969,.03066){8}{\line(1,0){.11969}}
\multiput(72.57,102.25)(.11969,.03066){8}{\line(1,0){.11969}}
\multiput(74.48,102.74)(.11969,.03066){8}{\line(1,0){.11969}}
\multiput(76.4,103.23)(.11969,.03066){8}{\line(1,0){.11969}}
\multiput(78.31,103.72)(.11969,.03066){8}{\line(1,0){.11969}}
\multiput(80.23,104.21)(.11969,.03066){8}{\line(1,0){.11969}}
\multiput(82.14,104.7)(.11969,.03066){8}{\line(1,0){.11969}}
\multiput(84.06,105.19)(.11969,.03066){8}{\line(1,0){.11969}}
\multiput(85.97,105.68)(.11969,.03066){8}{\line(1,0){.11969}}
%\end
%\dashline{1}(87,106)(89,93)
\multiput(86.93,105.93)(.0333,-.2167){4}{\line(0,-1){.2167}}
\multiput(87.2,104.2)(.0333,-.2167){4}{\line(0,-1){.2167}}
\multiput(87.46,102.46)(.0333,-.2167){4}{\line(0,-1){.2167}}
\multiput(87.73,100.73)(.0333,-.2167){4}{\line(0,-1){.2167}}
\multiput(88,99)(.0333,-.2167){4}{\line(0,-1){.2167}}
\multiput(88.26,97.26)(.0333,-.2167){4}{\line(0,-1){.2167}}
\multiput(88.53,95.53)(.0333,-.2167){4}{\line(0,-1){.2167}}
\multiput(88.8,93.8)(.0333,-.2167){4}{\line(0,-1){.2167}}
%\end
%\dashline{1}(89,93)(103.75,106)
\multiput(88.93,92.93)(.036967,.032581){19}{\line(1,0){.036967}}
\multiput(90.33,94.17)(.036967,.032581){19}{\line(1,0){.036967}}
\multiput(91.74,95.41)(.036967,.032581){19}{\line(1,0){.036967}}
\multiput(93.14,96.64)(.036967,.032581){19}{\line(1,0){.036967}}
\multiput(94.55,97.88)(.036967,.032581){19}{\line(1,0){.036967}}
\multiput(95.95,99.12)(.036967,.032581){19}{\line(1,0){.036967}}
\multiput(97.36,100.36)(.036967,.032581){19}{\line(1,0){.036967}}
\multiput(98.76,101.6)(.036967,.032581){19}{\line(1,0){.036967}}
\multiput(100.17,102.83)(.036967,.032581){19}{\line(1,0){.036967}}
\multiput(101.57,104.07)(.036967,.032581){19}{\line(1,0){.036967}}
\multiput(102.98,105.31)(.036967,.032581){19}{\line(1,0){.036967}}
%\end
%\dashline{1}(103.75,106)(107,93.25)
\multiput(103.68,105.93)(.03316,-.1301){7}{\line(0,-1){.1301}}
\multiput(104.14,104.11)(.03316,-.1301){7}{\line(0,-1){.1301}}
\multiput(104.61,102.29)(.03316,-.1301){7}{\line(0,-1){.1301}}
\multiput(105.07,100.47)(.03316,-.1301){7}{\line(0,-1){.1301}}
\multiput(105.54,98.64)(.03316,-.1301){7}{\line(0,-1){.1301}}
\multiput(106,96.82)(.03316,-.1301){7}{\line(0,-1){.1301}}
\multiput(106.47,95)(.03316,-.1301){7}{\line(0,-1){.1301}}
%\end
%\dashline{1}(119.5,93)(103.75,106.25)
\multiput(119.43,92.93)(-.039773,.03346){18}{\line(-1,0){.039773}}
\multiput(118,94.13)(-.039773,.03346){18}{\line(-1,0){.039773}}
\multiput(116.57,95.34)(-.039773,.03346){18}{\line(-1,0){.039773}}
\multiput(115.13,96.54)(-.039773,.03346){18}{\line(-1,0){.039773}}
\multiput(113.7,97.75)(-.039773,.03346){18}{\line(-1,0){.039773}}
\multiput(112.27,98.95)(-.039773,.03346){18}{\line(-1,0){.039773}}
\multiput(110.84,100.16)(-.039773,.03346){18}{\line(-1,0){.039773}}
\multiput(109.41,101.36)(-.039773,.03346){18}{\line(-1,0){.039773}}
\multiput(107.98,102.57)(-.039773,.03346){18}{\line(-1,0){.039773}}
\multiput(106.54,103.77)(-.039773,.03346){18}{\line(-1,0){.039773}}
\multiput(105.11,104.98)(-.039773,.03346){18}{\line(-1,0){.039773}}
%\end
%\dashline{1}(103.75,106.25)(136,93.25)
\multiput(103.68,106.18)(.08144,-.03283){11}{\line(1,0){.08144}}
\multiput(105.47,105.46)(.08144,-.03283){11}{\line(1,0){.08144}}
\multiput(107.26,104.74)(.08144,-.03283){11}{\line(1,0){.08144}}
\multiput(109.05,104.01)(.08144,-.03283){11}{\line(1,0){.08144}}
\multiput(110.85,103.29)(.08144,-.03283){11}{\line(1,0){.08144}}
\multiput(112.64,102.57)(.08144,-.03283){11}{\line(1,0){.08144}}
\multiput(114.43,101.85)(.08144,-.03283){11}{\line(1,0){.08144}}
\multiput(116.22,101.12)(.08144,-.03283){11}{\line(1,0){.08144}}
\multiput(118.01,100.4)(.08144,-.03283){11}{\line(1,0){.08144}}
\multiput(119.8,99.68)(.08144,-.03283){11}{\line(1,0){.08144}}
\multiput(121.6,98.96)(.08144,-.03283){11}{\line(1,0){.08144}}
\multiput(123.39,98.24)(.08144,-.03283){11}{\line(1,0){.08144}}
\multiput(125.18,97.51)(.08144,-.03283){11}{\line(1,0){.08144}}
\multiput(126.97,96.79)(.08144,-.03283){11}{\line(1,0){.08144}}
\multiput(128.76,96.07)(.08144,-.03283){11}{\line(1,0){.08144}}
\multiput(130.55,95.35)(.08144,-.03283){11}{\line(1,0){.08144}}
\multiput(132.35,94.62)(.08144,-.03283){11}{\line(1,0){.08144}}
\multiput(134.14,93.9)(.08144,-.03283){11}{\line(1,0){.08144}}
%\end
%\dashline{1}(136,93.25)(139,106.25)
\multiput(135.93,93.18)(.03333,.14444){6}{\line(0,1){.14444}}
\multiput(136.33,94.91)(.03333,.14444){6}{\line(0,1){.14444}}
\multiput(136.73,96.65)(.03333,.14444){6}{\line(0,1){.14444}}
\multiput(137.13,98.38)(.03333,.14444){6}{\line(0,1){.14444}}
\multiput(137.53,100.11)(.03333,.14444){6}{\line(0,1){.14444}}
\multiput(137.93,101.85)(.03333,.14444){6}{\line(0,1){.14444}}
\multiput(138.33,103.58)(.03333,.14444){6}{\line(0,1){.14444}}
\multiput(138.73,105.31)(.03333,.14444){6}{\line(0,1){.14444}}
%\end
%\dashline{1}(139,106.25)(153.75,93)
\multiput(138.93,106.18)(.036967,-.033208){19}{\line(1,0){.036967}}
\multiput(140.33,104.92)(.036967,-.033208){19}{\line(1,0){.036967}}
\multiput(141.74,103.66)(.036967,-.033208){19}{\line(1,0){.036967}}
\multiput(143.14,102.39)(.036967,-.033208){19}{\line(1,0){.036967}}
\multiput(144.55,101.13)(.036967,-.033208){19}{\line(1,0){.036967}}
\multiput(145.95,99.87)(.036967,-.033208){19}{\line(1,0){.036967}}
\multiput(147.36,98.61)(.036967,-.033208){19}{\line(1,0){.036967}}
\multiput(148.76,97.35)(.036967,-.033208){19}{\line(1,0){.036967}}
\multiput(150.17,96.08)(.036967,-.033208){19}{\line(1,0){.036967}}
\multiput(151.57,94.82)(.036967,-.033208){19}{\line(1,0){.036967}}
\multiput(152.98,93.56)(.036967,-.033208){19}{\line(1,0){.036967}}
%\end
%\dashline{1}(153.75,93)(157.75,96.5)
\multiput(153.68,92.93)(.038095,.033333){15}{\line(1,0){.038095}}
\multiput(154.82,93.93)(.038095,.033333){15}{\line(1,0){.038095}}
\multiput(155.97,94.93)(.038095,.033333){15}{\line(1,0){.038095}}
\multiput(157.11,95.93)(.038095,.033333){15}{\line(1,0){.038095}}
%\end
%\dashline{1}(153.75,93.25)(157.75,94.25)
\multiput(153.68,93.18)(.13333,.03333){5}{\line(1,0){.13333}}
\multiput(155.01,93.51)(.13333,.03333){5}{\line(1,0){.13333}}
\multiput(156.35,93.85)(.13333,.03333){5}{\line(1,0){.13333}}
%\end
%\dashline{1}(153.75,93.5)(157,102.25)
\multiput(153.68,93.43)(.0325,.0875){10}{\line(0,1){.0875}}
\multiput(154.33,95.18)(.0325,.0875){10}{\line(0,1){.0875}}
\multiput(154.98,96.93)(.0325,.0875){10}{\line(0,1){.0875}}
\multiput(155.63,98.68)(.0325,.0875){10}{\line(0,1){.0875}}
\multiput(156.28,100.43)(.0325,.0875){10}{\line(0,1){.0875}}
%\end
\put(13.25,87.75){\vector(1,0){5.75}}
\put(17,112.5){\vector(1,0){4.75}} 
\put(29,118){\vector(1,0){5.75}}
\put(30.5,135){\vector(1,0){5.75}} 
\thinlines
%\dashline{1}(36.25,131.75)(47.25,120.5)
\multiput(36.18,131.68)(.032164,-.032895){19}{\line(0,-1){.032895}}
\multiput(37.4,130.43)(.032164,-.032895){19}{\line(0,-1){.032895}}
\multiput(38.62,129.18)(.032164,-.032895){19}{\line(0,-1){.032895}}
\multiput(39.85,127.93)(.032164,-.032895){19}{\line(0,-1){.032895}}
\multiput(41.07,126.68)(.032164,-.032895){19}{\line(0,-1){.032895}}
\multiput(42.29,125.43)(.032164,-.032895){19}{\line(0,-1){.032895}}
\multiput(43.51,124.18)(.032164,-.032895){19}{\line(0,-1){.032895}}
\multiput(44.74,122.93)(.032164,-.032895){19}{\line(0,-1){.032895}}
\multiput(45.96,121.68)(.032164,-.032895){19}{\line(0,-1){.032895}}
%\end
\put(43.75,128.75){\makebox(0,0)[cc]{$\frac{3}{11}\pi$}}
\put(25.75,100.5){\makebox(0,0)[cc]{$\frac{3}{11}\pi$}}
\put(32.75,103.75){\makebox(0,0)[cc]{$\frac{3}{11}\pi$}}
\put(35,98){\makebox(0,0)[cc]{$\frac{4}{11}\pi$}}
\put(44.25,100.5){\makebox(0,0)[cc]{$\frac{2}{11}\pi$}}
\put(52,98.75){\makebox(0,0)[cc]{$\frac{1}{11}\pi$}}
\put(70,103.75){\makebox(0,0)[cc]{$\frac{1}{11}\pi$}}
\put(57,95){\makebox(0,0)[cc]{$\frac{1}{11}\pi$}}
\put(83.75,102.5){\makebox(0,0)[cc]{$\frac{5}{11}\pi$}}
\put(83,95.5){\makebox(0,0)[cc]{$\frac{5}{11}\pi$}}
\put(93,103.5){\makebox(0,0)[cc]{$\frac{5}{11}\pi$}}
\put(91.75,98.75){\makebox(0,0)[cc]{$\frac{3}{11}\pi$}}
\put(98,103.5){\makebox(0,0)[cc]{$\frac{3}{11}\pi$}}
\put(102.5,100){\makebox(0,0)[cc]{$\frac{4}{11}\pi$}}
\put(96.5,95.5){\makebox(0,0)[cc]{$\frac{3}{11}\pi$}}
\put(108.75,98){\makebox(0,0)[cc]{$\frac{2}{11}\pi$}}
\put(118.25,97.75){\makebox(0,0)[cc]{$\frac{1}{11}\pi$}}
\put(119.25,103){\makebox(0,0)[cc]{$\frac{1}{11}\pi$}}
\put(133.5,97){\makebox(0,0)[cc]{$\frac{5}{11}\pi$}}
\put(135,104.25){\makebox(0,0)[cc]{$\frac{5}{11}\pi$}}
\put(140.75,95.5){\makebox(0,0)[cc]{$\frac{5}{11}\pi$}}
\put(140.25,101.75){\makebox(0,0)[cc]{$\frac{3}{11}\pi$}}
\put(149.5,103.25){\makebox(0,0)[cc]{$\frac{3}{11}\pi$}}
\put(25.75,108.75){\makebox(0,0)[cc]{$A_{0}$}}
\put(44.5,108.75){\makebox(0,0)[cc]{$A_{2}$}}
\put(59.75,108.75){\makebox(0,0)[cc]{$A_{3}$}}
\put(87.25,108.75){\makebox(0,0)[cc]{$A_{4}$}}
\put(103.75,108.75){\makebox(0,0)[cc]{$A_{6}$}}
\put(139,108.25){\makebox(0,0)[cc]{$A_{10}$}}
\put(135.5,90){\makebox(0,0)[cc]{$A_{9}$}}
\put(153.5,90){\makebox(0,0)[cc]{$A_{11}$}}
\put(119,90){\makebox(0,0)[cc]{$A_{8}$}}
\put(106.25,90){\makebox(0,0)[cc]{$A_{7}$}}
\put(89.25,90){\makebox(0,0)[cc]{$A_{5}$}}
\put(35.75,90){\makebox(0,0)[cc]{$A_{1}$}}
\put(38.75,135.25){\makebox(0,0)[cc]{$A_{0}$}}
\put(82,117.25){\makebox(0,0)[cc]{(a)}}
\put(82,86.25){\makebox(0,0)[cc]{(b)}}
\put(82,76.25){\makebox(0,0)[cc]{Figure 2.9}}
\end{picture}
\end{center}

The proof of convergence for this folding procedure is similar to that for the primary folding procedure of period $2$. We leave the details to the reader. Let us explore here what we can do with this $\{ 1,3,1 \}$-tape. First, note that starting with the putative angle $\left( \frac{3}{11} \right) \pi$ at the top of the tape, we produce a putative angle of $\left( \frac{l}{11} \right) \pi$ at the bottom of the tape, then a putative angle $\left( \frac{5}{11} \right) \pi$ at the top of the tape, then a putative angle of $\left( \frac{3}{11} \right) \pi$ at the \emph{bottom} of the tape, and so on. Hence we see that we could use this tape to fold a star $\left\{ \frac{11}{3} \right\}$-gon, a convex $11$-gon, and a star $\left\{ \frac{11}{5} \right\}$-gon. More still is true; for, as we see, if there are crease lines enabling us to fold a star $\left\{ \frac{11}{a} \right\}$-gon, there will be crease lines enabling us to fold star $\left\{ \frac{11}{2^{k}a} \right\}$-gons, where $k \geq 0$ takes all values such that $2^{k+1}a < 11$. 

These features, as described for $b = 11$, apply for any odd number $b$. However, this tape has a special symmetry as a consequence of its \emph{odd} period; namely that if it is ``flipped'' about the horizontal line halfway between its parallel edges, the result is a translate of the original tape. As a practical matter, this special symmetry of the tape means that we can use either the top edge or the bottom edge of the tape to construct our polygons. On tapes with an \emph{even} period, the top edge and the bottom edge of the tape are not translates of each other (after the horizontal flip), which simply means that care must be taken in choosing the edge of the tape used to construct a specific polygon. Figures 2.l0(a,b) show the completed $\left\{ \frac{11}{3} \right\}, \left\{ \frac{11}{4} \right\}$-gons, respectively. (How would you guess we prepared these illustrations?!)

\begin{center}
%TeXCAD Options
%\grade{\on}
%\emlines{\off}
%\epic{\off}
%\beziermacro{\on}
%\reduce{\on}
%\snapping{\off}
%\quality{8.00}
%\graddiff{0.01}
%\snapasp{1}
%\zoom{4.0000}
\unitlength .6mm % = 2.85pt
\linethickness{0.4pt}
\ifx\plotpoint\undefined\newsavebox{\plotpoint}\fi % GNUPLOT compatibility
\begin{picture}(190,100)(0,25)
%\emline(31.5,115.25)(15.25,116.75)
\multiput(31.5,115.25)(-.3611111,.0333333){45}{\line(-1,0){.3611111}}
%\end
\put(15.25,116.75){\line(0,-1){13}}
%\emline(15.25,103.75)(.5,97.75)
\multiput(15.25,103.75)(-.08286517,-.03370787){178}{\line(-1,0){.08286517}}
%\end
%\emline(.5,97.75)(8.25,85)
\multiput(.5,97.75)(.03369565,-.05543478){230}{\line(0,-1){.05543478}}
%\end
%\emline(8.25,85)(0,73.5)
\multiput(8.25,85)(-.033673469,-.046938776){245}{\line(0,-1){.046938776}}
%\end
%\emline(15.25,67.5)(14.75,54.5)
\multiput(15.25,67.5)(-.033333,-.866667){15}{\line(0,-1){.866667}}
%\end
%\emline(14.75,54.5)(30.75,56.25)
\multiput(14.75,54.5)(.3076923,.0336538){52}{\line(1,0){.3076923}}
%\end
%\emline(30.75,56.25)(40,44)
\multiput(30.75,56.25)(.033636364,-.044545455){275}{\line(0,-1){.044545455}}
%\end
%\emline(40,44)(52.75,53.5)
\multiput(40,44)(.045212766,.033687943){282}{\line(1,0){.045212766}}
%\end
%\emline(52.75,53.5)(68,48)
\multiput(52.75,53.5)(.0929878,-.03353659){164}{\line(1,0){.0929878}}
%\end
%\emline(68,48)(72.5,61.25)
\multiput(68,48)(.03358209,.0988806){134}{\line(0,1){.0988806}}
%\end
%\emline(72.5,61.25)(88.75,63.75)
\multiput(72.5,61.25)(.2166667,.0333333){75}{\line(1,0){.2166667}}
%\end
%\emline(88.75,63.75)(84,77)
\multiput(88.75,63.75)(-.03368794,.09397163){141}{\line(0,1){.09397163}}
%\end
%\emline(84,77)(97.25,86.25)
\multiput(84,77)(.048181818,.033636364){275}{\line(1,0){.048181818}}
%\end
%\emline(97.25,86.25)(84.75,96)
\multiput(97.25,86.25)(-.043252595,.033737024){289}{\line(-1,0){.043252595}}
%\end
%\emline(84.75,96)(89.75,109.25)
\multiput(84.75,96)(.03355705,.08892617){149}{\line(0,1){.08892617}}
%\end
%\emline(89.75,109.25)(72,111)
\multiput(89.75,109.25)(-.3413462,.0336538){52}{\line(-1,0){.3413462}}
%\end
%\emline(72,111)(67.75,122.75)
\multiput(72,111)(-.03373016,.09325397){126}{\line(0,1){.09325397}}
%\end
%\emline(67.75,122.75)(52.5,118)
\multiput(67.75,122.75)(-.10815603,-.03368794){141}{\line(-1,0){.10815603}}
%\end
%\emline(41.5,115.25)(68.5,94.75)
\multiput(41.5,115.25)(.044407895,-.033717105){608}{\line(1,0){.044407895}}
%\end
%\emline(68.5,94.75)(61.5,114.5)
\multiput(68.5,94.75)(-.03365385,.09495192){208}{\line(0,1){.09495192}}
%\end
%\emline(61.5,114.5)(67.75,122.75)
\multiput(61.5,114.5)(.03360215,.04435484){186}{\line(0,1){.04435484}}
%\end
\put(61.75,114.25){\line(1,0){8.75}}
%\emline(89.5,109)(79.5,101.75)
\multiput(89.5,109)(-.04651163,-.03372093){215}{\line(-1,0){.04651163}}
%\end
\put(79.5,101.75){\line(5,-2){6.25}}
%\emline(79.75,102)(69.25,73)
\multiput(79.75,102)(-.033653846,-.092948718){312}{\line(0,-1){.092948718}}
%\end
%\emline(69.25,73)(86.25,85)
\multiput(69.25,73)(.047752809,.033707865){356}{\line(1,0){.047752809}}
%\end
%\emline(86.25,85)(97.5,86)
\multiput(86.25,85)(.375,.033333){30}{\line(1,0){.375}}
%\end
%\emline(86.5,85)(89.25,80.75)
\multiput(86.5,85)(.0335366,-.0518293){82}{\line(0,-1){.0518293}}
%\end
%\emline(69.25,73)(56.75,64.5)
\multiput(69.25,73)(-.049603175,-.033730159){252}{\line(-1,0){.049603175}}
%\end
%\emline(56.75,64.5)(81.5,68.5)
\multiput(56.75,64.5)(.20798319,.03361345){119}{\line(1,0){.20798319}}
%\end
%\emline(81.5,68.5)(88.75,63.75)
\multiput(81.5,68.5)(.05141844,-.03368794){141}{\line(1,0){.05141844}}
%\end
\put(81.5,68.5){\line(0,-1){5.75}}
%\emline(56.75,64.5)(44.75,63)
\multiput(56.75,64.5)(-.2666667,-.0333333){45}{\line(-1,0){.2666667}}
%\end
%\emline(44.75,63)(31.25,67)
\multiput(44.75,63)(-.11344538,.03361345){119}{\line(-1,0){.11344538}}
%\end
%\emline(31.25,67)(42.5,53)
\multiput(31.25,67)(.033682635,-.041916168){334}{\line(0,-1){.041916168}}
%\end
%\emline(42.5,53)(40.25,44.5)
\multiput(42.5,53)(-.0335821,-.1268657){67}{\line(0,-1){.1268657}}
%\end
%\emline(42.5,53)(36,49.5)
\multiput(42.5,53)(-.0625,-.03365385){104}{\line(-1,0){.0625}}
%\end
%\emline(45.5,62.75)(66,56.5)
\multiput(45.5,62.75)(.11021505,-.03360215){186}{\line(1,0){.11021505}}
%\end
%\emline(66,56.5)(68,48)
\multiput(66,56.5)(.0333333,-.1416667){60}{\line(0,-1){.1416667}}
%\end
%\emline(66.25,56.5)(61.75,50.25)
\multiput(66.25,56.5)(-.03358209,-.04664179){134}{\line(0,-1){.04664179}}
%\end
%\emline(59,58.5)(72.25,61.25)
\multiput(59,58.5)(.1615854,.0335366){82}{\line(1,0){.1615854}}
%\end
%\emline(64,59.25)(63.25,57.75)
\multiput(64,59.25)(-.032609,-.065217){23}{\line(0,-1){.065217}}
%\end
%\emline(61.5,59)(61,58)
\multiput(61.5,59)(-.033333,-.066667){15}{\line(0,-1){.066667}}
%\end
%\emline(84.5,95.75)(79,80)
\multiput(84.5,95.75)(-.03353659,-.09603659){164}{\line(0,-1){.09603659}}
%\end
%\emline(68.75,94.75)(72.75,83)
\multiput(68.75,94.75)(.03361345,-.0987395){119}{\line(0,-1){.0987395}}
%\end
%\emline(72,111)(77.25,96.25)
\multiput(72,111)(.03365385,-.09455128){156}{\line(0,-1){.09455128}}
%\end
%\emline(52.5,118)(63.75,108.5)
\multiput(52.5,118)(.039893617,-.033687943){282}{\line(1,0){.039893617}}
%\end
%\emline(60.25,112)(62,112.5)
\multiput(60.25,112)(.116667,.033333){15}{\line(1,0){.116667}}
%\end
%\emline(60,111.5)(62.25,111.25)
\multiput(60,111.5)(.28125,-.03125){8}{\line(1,0){.28125}}
%\end
%\emline(31.5,115)(43.5,113.75)
\multiput(31.5,115)(.3157895,-.0328947){38}{\line(1,0){.3157895}}
%\end
%\emline(42,115.25)(41.25,114.5)
\multiput(42,115.25)(-.032609,-.032609){23}{\line(0,-1){.032609}}
%\end
%\emline(41.5,114.75)(40.75,114.25)
\multiput(41.5,114.75)(-.05,-.033333){15}{\line(-1,0){.05}}
%\end
%\emline(23.75,115.5)(22,108)
\multiput(23.75,115.5)(-.0336538,-.1442308){52}{\line(0,-1){.1442308}}
%\end
%\emline(22,108)(15.5,116.75)
\multiput(22,108)(-.03367876,.04533679){193}{\line(0,1){.04533679}}
%\end
%\emline(22.5,108)(56,104.25)
\multiput(22.5,108)(.29910714,-.03348214){112}{\line(1,0){.29910714}}
%\end
\put(76,101.75){\line(1,0){3.25}}
%\emline(15.5,104)(25,107.75)
\multiput(15.5,104)(.08482143,.03348214){112}{\line(1,0){.08482143}}
%\end
%\emline(43.5,105.5)(12.5,92.5)
\multiput(43.5,105.5)(-.080310881,-.033678756){386}{\line(-1,0){.080310881}}
%\end
%\emline(12.5,92.5)(.75,97.5)
\multiput(12.5,92.5)(-.07885906,.03355705){149}{\line(-1,0){.07885906}}
%\end
%\emline(12.25,92.75)(9.25,101.25)
\multiput(12.25,92.75)(-.0337079,.0955056){89}{\line(0,1){.0955056}}
%\end
%\emline(17.25,94.5)(8.25,84.75)
\multiput(17.25,94.5)(-.033707865,-.036516854){267}{\line(0,-1){.036516854}}
%\end
%\emline(31,100.25)(12.75,76.5)
\multiput(31,100.25)(-.033733826,-.043900185){541}{\line(0,-1){.043900185}}
%\end
%\emline(12.75,76.5)(.25,73.5)
\multiput(12.75,76.5)(-.1404494,-.0337079){89}{\line(-1,0){.1404494}}
%\end
\put(12.75,76.5){\line(-3,2){7.5}}
%\emline(16.5,80.75)(15.5,66.5)
\multiput(16.5,80.75)(-.033333,-.475){30}{\line(0,-1){.475}}
%\end
%\emline(25.25,93.25)(23.25,61)
\multiput(25.25,93.25)(-.0333333,-.5375){60}{\line(0,-1){.5375}}
%\end
%\emline(23.25,61)(15,54.5)
\multiput(23.25,61)(-.04274611,-.03367876){193}{\line(-1,0){.04274611}}
%\end
\put(23,60.75){\line(-4,1){8}} 
\put(23.75,65.75){\line(3,-4){8.25}}
%\emline(23.5,61.5)(26.5,61.75)
\multiput(23.5,61.5)(.375,.03125){8}{\line(1,0){.375}}
%\end
%\emline(.25,73.75)(15.75,69.75)
\multiput(.25,73.75)(.1302521,-.03361345){119}{\line(1,0){.1302521}}
%\end
%\emline(31.25,67.25)(24,76.75)
\multiput(31.25,67.25)(-.03372093,.04418605){215}{\line(0,1){.04418605}}
%\end
\put(38.25,58.25){\line(3,-1){16.5}}
%\emline(41,54.75)(43.5,56.5)
\multiput(41,54.75)(.0480769,.0336538){52}{\line(1,0){.0480769}}
%\end
%\emline(83.75,76.75)(68.75,66.25)
\multiput(83.75,76.75)(-.048076923,-.033653846){312}{\line(-1,0){.048076923}}
%\end
%\emline(77,72.25)(74.25,67.5)
\multiput(77,72.25)(-.0335366,-.0579268){82}{\line(0,-1){.0579268}}
%\end
%\emline(76.5,71)(77.25,68)
\multiput(76.5,71)(.032609,-.130435){23}{\line(0,-1){.130435}}
%\end
%\emline(81,85)(82,82)
\multiput(81,85)(.033333,-.1){30}{\line(0,-1){.1}}
%\end
%\emline(52.75,117.75)(43.25,125.5)
\multiput(52.75,117.75)(-.04130435,.03369565){230}{\line(-1,0){.04130435}}
%\end
%\emline(41.25,115.25)(43.25,125.75)
\multiput(41.25,115.25)(.0333333,.175){60}{\line(0,1){.175}}
%\end
%\emline(41.5,115.75)(49.25,121)
\multiput(41.5,115.75)(.04967949,.03365385){156}{\line(1,0){.04967949}}
%\end
%\emline(43,125.25)(33.25,114.75)
\multiput(43,125.25)(-.033737024,-.03633218){289}{\line(0,-1){.03633218}}
%\end
%\emline(12.75,92.5)(13.75,91.25)
\multiput(12.75,92.5)(.033333,-.041667){30}{\line(0,-1){.041667}}
%\end
%\emline(13.5,91.5)(13.75,91)
\multiput(13.5,91.5)(.03125,-.0625){8}{\line(0,-1){.0625}}
%\end
%\emline(81.5,86.5)(84.75,84)
\multiput(81.5,86.5)(.0433333,-.0333333){75}{\line(1,0){.0433333}}
%\end
\thicklines
%\emline(137,118.25)(133.5,104)
\multiput(137,118.25)(-.03365385,-.13701923){104}{\line(0,-1){.13701923}}
%\end
%\emline(133.5,104)(118.25,109.25)
\multiput(133.5,104)(-.09775641,.03365385){156}{\line(-1,0){.09775641}}
%\end
%\emline(118.25,109.25)(126.25,94.75)
\multiput(118.25,109.25)(.03361345,-.06092437){238}{\line(0,-1){.06092437}}
%\end
%\emline(126.25,94.75)(108.75,89.75)
\multiput(126.25,94.75)(-.11744966,-.03355705){149}{\line(-1,0){.11744966}}
%\end
%\emline(108.75,89.75)(124,80.75)
\multiput(108.75,89.75)(.057116105,-.033707865){267}{\line(1,0){.057116105}}
%\end
%\emline(124,80.75)(113.75,69.5)
\multiput(124,80.75)(-.033717105,-.037006579){304}{\line(0,-1){.037006579}}
%\end
\put(113.75,69.5){\line(1,0){14.75}}
%\emline(128.5,69.5)(126.5,54)
\multiput(128.5,69.5)(-.0333333,-.2583333){60}{\line(0,-1){.2583333}}
%\end
%\emline(126.5,54)(142.5,63.25)
\multiput(126.5,54)(.058181818,.033636364){275}{\line(1,0){.058181818}}
%\end
%\emline(142.5,63.25)(149.25,48.75)
\multiput(142.5,63.25)(.03358209,-.0721393){201}{\line(0,-1){.0721393}}
%\end
%\emline(149.25,48.75)(156.75,63.25)
\multiput(149.25,48.75)(.03363229,.06502242){223}{\line(0,1){.06502242}}
%\end
%\emline(156.75,63.25)(173.5,53.75)
\multiput(156.75,63.25)(.059397163,-.033687943){282}{\line(1,0){.059397163}}
%\end
\put(173.5,53.75){\line(-1,5){3.25}}
%\emline(170.25,70)(187.5,69.75)
\multiput(170.25,70)(2.15625,-.03125){8}{\line(1,0){2.15625}}
%\end
%\emline(187.5,69.75)(175.5,80.25)
\multiput(187.5,69.75)(-.038461538,.033653846){312}{\line(-1,0){.038461538}}
%\end
%\emline(175.5,80.25)(190,88)
\multiput(175.5,80.25)(.06304348,.03369565){230}{\line(1,0){.06304348}}
%\end
%\emline(190,88)(171.25,93.25)
\multiput(190,88)(-.12019231,.03365385){156}{\line(-1,0){.12019231}}
%\end
%\emline(171.25,93.25)(179.25,107.75)
\multiput(171.25,93.25)(.03361345,.06092437){238}{\line(0,1){.06092437}}
%\end
%\emline(179.25,107.75)(162.5,103.25)
\multiput(179.25,107.75)(-.125,-.03358209){134}{\line(-1,0){.125}}
%\end
%\emline(162.5,103.25)(159,119.25)
\multiput(162.5,103.25)(-.03365385,.15384615){104}{\line(0,1){.15384615}}
%\end
%\emline(159,119.25)(148,107)
\multiput(159,119.25)(-.033639144,-.037461774){327}{\line(0,-1){.037461774}}
%\end
\put(148,107){\line(-1,1){10.75}} \put(143.25,111){\line(-1,0){8}}
%\emline(135.25,111)(157.25,87.5)
\multiput(135.25,111)(.033690658,-.035987749){653}{\line(0,-1){.035987749}}
%\end
%\emline(157.25,87.5)(152,111.75)
\multiput(157.25,87.5)(-.03365385,.15544872){156}{\line(0,1){.15544872}}
%\end
%\emline(154,112.5)(161,109.25)
\multiput(154,112.5)(.07216495,-.03350515){97}{\line(1,0){.07216495}}
%\end
%\emline(133.75,103.75)(145.75,99.5)
\multiput(133.75,103.75)(.0952381,-.03373016){126}{\line(1,0){.0952381}}
%\end
%\emline(116.5,84.75)(120,93.25)
\multiput(116.5,84.75)(.03365385,.08173077){104}{\line(0,1){.08173077}}
%\end
%\emline(130.25,105.25)(122.75,101.25)
\multiput(130.25,105.25)(-.06302521,-.03361345){119}{\line(-1,0){.06302521}}
%\end
%\emline(124.25,99)(154.5,90.5)
\multiput(124.25,99)(.120039683,-.033730159){252}{\line(1,0){.120039683}}
%\end
%\emline(154.5,90.5)(155.5,89.5)
\multiput(154.5,90.5)(.033333,-.033333){30}{\line(0,-1){.033333}}
%\end
%\emline(157.25,87.25)(158,83.5)
\multiput(157.25,87.25)(.032609,-.163043){23}{\line(0,-1){.163043}}
%\end
%\emline(158,83.5)(168,104.75)
\multiput(158,83.5)(.033670034,.071548822){297}{\line(0,1){.071548822}}
%\end
%\emline(171,105.5)(175.75,103)
\multiput(171,105.5)(.0633333,-.0333333){75}{\line(1,0){.0633333}}
%\end
%\emline(170.75,105.25)(174.5,99.75)
\multiput(170.75,105.25)(.03348214,-.04910714){112}{\line(0,-1){.04910714}}
%\end
%\emline(162.5,103.25)(164,96.5)
\multiput(162.5,103.25)(.0333333,-.15){45}{\line(0,-1){.15}}
%\end
%\emline(171.25,93.25)(167,84.25)
\multiput(171.25,93.25)(-.03373016,-.07142857){126}{\line(0,-1){.07142857}}
%\end
%\emline(167,84.25)(179.5,90.75)
\multiput(167,84.25)(.06476684,.03367876){193}{\line(1,0){.06476684}}
%\end
%\emline(167,84.25)(154,77.25)
\multiput(167,84.25)(-.0625,-.03365385){208}{\line(-1,0){.0625}}
%\end
%\emline(156.25,78.5)(158,84)
\multiput(156.25,78.5)(.0336538,.1057692){52}{\line(0,1){.1057692}}
%\end
%\emline(153.75,77.5)(179.75,76.5)
\multiput(153.75,77.5)(.866667,-.033333){30}{\line(1,0){.866667}}
%\end
%\emline(181.5,75)(176.75,69.5)
\multiput(181.5,75)(-.03368794,-.03900709){141}{\line(0,-1){.03900709}}
%\end
\put(153.5,77.5){\line(-1,0){6}}
%\emline(147.5,77.5)(171.75,63)
\multiput(147.5,77.5)(.056395349,-.03372093){430}{\line(1,0){.056395349}}
%\end
\put(172,61){\line(-1,-1){4}} 
\put(147.5,77.25){\line(-5,3){6.25}}
%\emline(141.25,81)(153.5,57.75)
\multiput(141.25,81)(.033653846,-.063873626){364}{\line(0,-1){.063873626}}
%\end
%\emline(153.5,57.75)(152.75,57)
\multiput(153.5,57.75)(-.032609,-.032609){23}{\line(0,-1){.032609}}
%\end
%\emline(151.5,53.5)(146.75,54.25)
\multiput(151.5,53.5)(-.206522,.032609){23}{\line(-1,0){.206522}}
%\end
%\emline(132,57.25)(129.5,59.5)
\multiput(132,57.25)(-.0373134,.0335821){67}{\line(-1,0){.0373134}}
%\end
%\emline(129.5,59.5)(127.25,61.25)
\multiput(129.5,59.5)(-.0432692,.0336538){52}{\line(-1,0){.0432692}}
%\end
%\emline(119.25,69.5)(117.5,73.5)
\multiput(119.25,69.5)(-.0336538,.0769231){52}{\line(0,1){.0769231}}
%\end
%\emline(139.75,83.75)(135.75,59.5)
\multiput(139.75,83.75)(-.03361345,-.20378151){119}{\line(0,-1){.20378151}}
%\end
%\emline(141.25,81)(139.75,84)
\multiput(141.25,81)(-.0333333,.0666667){45}{\line(0,1){.0666667}}
%\end
%\emline(140,83.75)(140.25,86.75)
\multiput(140,83.75)(.03125,.375){8}{\line(0,1){.375}}
%\end
%\emline(140.25,86.5)(124.25,69.5)
\multiput(140.25,86.5)(-.033684211,-.035789474){475}{\line(0,-1){.035789474}}
%\end
%\emline(140.5,86.5)(143,89.25)
\multiput(140.5,86.5)(.0333333,.0366667){75}{\line(0,1){.0366667}}
%\end
%\emline(143,89.25)(119.25,83.75)
\multiput(143,89.25)(-.14481707,-.03353659){164}{\line(-1,0){.14481707}}
%\end
%\emline(143.25,89.5)(151.5,91.5)
\multiput(143.25,89.5)(.1375,.0333333){60}{\line(1,0){.1375}}
%\end
\put(143.75,111){\line(-1,0){.5}}
\put(153.75,112.5){\line(-1,0){.5}} 
\put(175.5,103){\line(1,0){.75}}
\put(40.75,37){\makebox(0,0)[cc]{(a)}}
\put(149,37){\makebox(0,0)[cc]{(b)}}
\put(98.5,30){\makebox(0,0)[cc]{Figure 2.10}}
\end{picture}
\end{center}

Now, to set the scene for the number theory, let us look at the patterns in the arithmetic of the computations when $a = 3$ and $b = 11$. Referring to Figure 2.9, we observe that

\begin{center}
\begin{tabular}{lclclc}
\small{the angle to the}           &   & \small{is of the form}              &   & \small{and the number of}     &   \\
\small{right of $A_n$ where $n =$} & 0 & \small{$(a_i/11)\pi$ where $a_i =$} & 3 & \small{bisections at $A_n =$} & 3 \\
                                   & 1 &                                     & 1 &                               & 1 \\
                                   & 4 &                                     & 5 &                               & 1 \\
                                   & 5 &                                     & 3 &                               & 3 \\
                                   & 6 &                                     & 1 &                               & 1 \\
                                   & 9 &                                     & 5 &                               & 1 \\
\end{tabular}
\end{center}

We could write this in shorthand form as follows:
\begin{equation}
(b=) 11 
\begin{vmatrix}
(a=) & 3 & 1 & 5 \\
     & 3 & 1 & 1 
\end{vmatrix}  \tag{2.7}
\end{equation}

Observe that had we started with the putative angle of $\frac{\pi}{11}$, then the \emph{symbol} (2.7) would have taken the form
\begin{equation}
(b=) 11 
\begin{vmatrix}
(a=) & 1 & 5 & 3 \\
     & 1 & 1 & 3 
\end{vmatrix}  \tag{2.7'}
\end{equation}

In fact, it should be clear that we can \emph{start anywhere} (with $a = 1$, $3$, or $5$), and the resulting symbol, analogous to (2.7), will be the cyclic permutation of the interior that places our choice of $a$ in the first position along the top row.

The generalization of this symbol appears at the beginning of Section 4. But before leaving this section, let us comment that given any two odd numbers $a$ and $b$, with $a < \left( \frac{b}{2} \right)$, there is always a completely determined unique symbol like the one above (we do not need $a,b$ relatively prime). Appropriately interpreted, we can use this symbol to read off the folding procedure that produces the angle of $\left( \frac{a}{b} \right) \pi$ along the top edge of the tape, so that a symbol such as (2.7) encodes the folding procedure for producing a regular star $\left\{ \frac{b}{a} \right\}$-gon, and also tells us what other star polygons we can obtain from the same tape (of course, for each symbol, a diagram similar to Figure 2.9 can be drawn to illustrate the relative positions of the angles $\left( \frac{a_{i}}{b} \right) \pi$ ).

The answer to question 4 is that the folding process described above is the \emph{most efficient} one possible. That is, there could not be any folding procedure of this type that would produce the required star polygons with fewer folds. It is also optimal from the point of view of ``difficulty of execution'', for it keeps the number of bisections at each vertex to a minimum. 

These last comments are explained as follows. If a folding procedure $\{ k_{1}, k_{2},$ $\cdots, k_{r} \}$ produces an angle $\left( \frac{a}{b} \right) \pi$, then (see (4.3), (4.4) ) $b\ |\ 2^k \pm 1$, where $\displaystyle k = \sum^{r}_{i=1} k_{i}$. If we adopt the procedures described in this section, we will have a procedure $\{ \ell_{1}, \ell_{2}, \cdots, \ell_{s} \}$, such that $\displaystyle \ell = \sum^{s}_{j=1} \ell_{j}$ is the \emph{smallest} number $m$ such that $b\ |\ 2^{m} \pm 1$; that is, $\ell$ is the \emph{quasi-order} of $2 \mod{b}$. In fact, $\ell\ |\ k$ and $r$ will be a multiple of $s$. See Section 4 for details.




\newpage

\section{The Number Theory Of Period 2 Folding}
\markboth{3. Theory Of Period 2 Folding}{3. Theory Of Period 2 Folding}


We have seen that the $\{ D^{m} U^{n} \}$-procedure, or more briefly the $(m,n)$-procedure, enables us to fold star $\left\{ \frac{b}{a} \right\}$-gons, where $b = 2^{m+n}-1$, $a = 2^{n}-1$. We will generalize to the extent of replacing the base $2$ by an arbitrary positive integer $t \geq 2$, so that we are concerned with rational numbers
\begin{equation}
\frac{t^{u}-1}{t^{v}-1},\ u > v \geq 1  \tag{3.1}
\end{equation}

Our first result establishes how to reduce such fractions (3.1).

\begin{description}

  \item[Theorem 3.1:] Let $d = \gcd(u,v)$. Then $\gcd (t^{u}-1, t^{v}-1) = t^{d}-1$.

  \item[Proof:] Let us use base $t$, and let $\imath$ = t - 1. Then
    \begin{align*}
    & t^{u}-1 = \imath\ \imath\ .\ .\ .\ .\ \imath & \text{ ($u$ times)} \\
    & t^{v}-1 = \imath\ \imath\ .\ .\ \imath       & \text{ ($v$ times)}
    \end{align*}

    It is now obvious that if we divide $(t^{u}-1)$ by $(t^{v}-1)$, the remainder is $(t^{r}-1)$, where $r$ is the remainder when we divide $u$ by $v$. Thus, an application of the Euclidean algorithm establishes the theorem.

  \item[Remarks:] $\ $
    \begin{enumerate}[\hspace{1cm}(i)]
      \item This theorem is propaganda for the Euclidean algorithm. We defy the reader to prove it by factoring $(t^{u}-1), (t^{v}-1)$ into prime factors.
      \item The statement becomes false if we replace $\gcd$ by $\text{lcm}$. Thus if $\ell = \text{lcm}(u,v)$, then it is obvious that $\text{lcm} (t^{u}-1, t^{v}-1)$ divides $(t^{\ell}-1)$, but it may divide $(t^{\ell}-1)$ properly. For example, $\text{lcm} (2^{2}-1, 2^{3}-1) = 21$, and $21\ |\ (2^{6}-1) = 63$.
    \end{enumerate}

  \item[Corollary 3.2:] $\frac{(t^{u}-1)}{(t^{v}-1)}$ is an integer if and only if $v\ |\ u$.

\end{description}

We call such integers \emph{$t$-folding numbers}, and write $F_{t}$ for the set of such numbers. It is sometimes convenient to include $1$ as a \emph{degenerate} $t$-folding number; we write $\bar{F}_{t}$ for $F_{t} \cap \{1\}$.

\begin{description}

  \item[Corollary 3.3:] The rational number (3.1) reduces to the irreducible fraction $\frac{\beta}{\alpha}$, where $\alpha \in \bar{F}_{t}$, $\beta \in F_{t}$, and
    \begin{equation}
    \beta = \frac{t^{u}-1}{t^{d}-1}, \quad  \alpha = \frac{t^{v}-1}{t^{d}-1}, \quad  d = \gcd(u,v)  \tag{3.2}
    \end{equation}

    Let us change notation, setting $d = x$, $u = xy$, $v = xy'$. Thus
    \begin{equation}
    \beta = \frac{t^{xy}-1}{t^{x}-1},\ x \geq 1,\ y \geq 2  \tag{3.3}
    \end{equation}

    Then the question naturally arises of how we should express the possible values of $\alpha$ in terms of $x$ and $y$. To answer this, let us write $\beta$ in base $t$, so that
    \begin{equation}
    \beta = 10\ldots0\ \ 10\ldots0\ \ \cdots\ \ 10\ldots01  \tag{3.4}
    \end{equation}
    where the \emph{repeating} part $10\ldots0$ consists of $1$ followed by $(x-1)$ zeros, and there are $y$ $1$s on the right-hand side of equation (3.4). Then $\alpha$ is a \emph{(right) section} of $\beta$, meaning that the $t$-numeral for $\alpha$ has the same repeating part but fewer $1$s. In fact, it has $y'$ $1$s, because
    \begin{equation}
    \alpha = \frac{t^{xy'}-1}{t^{x}-1}  \tag{3.5}
    \end{equation}

    Since $\gcd (xy,xy') = x$, it follows that $y'$ is restricted to those integers $< y$ which are prime to $y$. Thus we obtain $\phi(y)$ possible values of $\alpha$, which correspond to the \emph{relatively prime sections} of $\beta$.

  \item[Example:] Consider the $2$-folding number $21$. Now, $21 = \frac{(2^{6}-1)}{(2^{2}-1)}$, so that $x = 2$, $y = 3$. Thus the possible values of $\alpha$ are given by (3.5) with $t = 2$, $x = 2$, and $y' =$ $1$ or $2$. Of course, $y' = 1$ gives $\alpha = 1$; and $y' = 2$ gives $\alpha = \frac{(2^{4}-1)}{(2^{2}-1)} = 5$. Thus our period $2$ folding procedure enables us to fold the convex $2l$-gon or the star $\{ \frac{2l}{5} \}$-gon. Obviously, the convex $2l$-gon is obtained by the $(4,2)$-procedure, while the star $\{ \frac{2l}{5} \}$-gon is obtained by the $(2,4)$-procedure. (Alternatively, we can say that it is obtained by the $(4,2)$-procedure, but using the FAT-algorithm on the \emph{bottom} edge of the tape.)

\end{description}

We now prove two basic theorems on $t$-folding numbers. For $\beta$ of the form (3.3), the base $t$ representation (3.4) shows that $x,y$ are determined by $\beta$. We write $\tau(\beta) = xy$.

\begin{description}

  \item[Theorem 3.4:] Let $a > 1$ be an integer prime to $t$. Then for any $x$, there exists $\beta$ of the form (3.3) with $a\ |\ \beta$. Moreover, there exists a positive integer $h(a)$, called the \emph{height} of $a$, such that
    \begin{enumerate}[\hspace{1cm}(i)]
      \item there exists a $\beta_{0} \in F_{t}$, with $a\ |\ \beta_{0}$ and $\tau(\beta_{0}) = h(a)$;
      \item if $\beta \in F_{t}$ and $a\ |\ \beta$, then $h(a)\ |\ \tau(\beta)$.
    \end{enumerate}

  \item[Proof:] Now $t$ is prime to $a$ and to $t^{x}-1$, so that $t$ is prime to $a(t^{x}-1)$. Let $z$ be the order of $t \mod{a(t^{x}-1)}$. Then $t^{z} \equiv 1 \mod{a(t^{x}-1)}$. Since $t^{x}-1\ |\ t^{z}-1$, it follows by Corollary 3.2 that $x\ |\ z$, so we may write $z = xy$. Also $y \geq 2$, since $a > 1$. Thus $\frac{(t^{xy}-1)}{(t^{x}-1)} \in F$, and $a\ |\ \frac{(t^{xy}-1)}{(t^{x}-1)}$. We use the above argument with $x = 1$ to obtain $\beta_{0} = \frac{(t^{y}-1)}{(t-1)}$, and set $y_{0} = h(a)$, establishing (i). Now, if $a\ |\ \frac{(t^{xy}-1)}{(t^{x}-1)}$, then $a\ |\ \frac{(t^{xy}-1)}{(t-1)}$, so that $t^{xy} \equiv 1 \mod{a(t-1)}$. But $h(a)$ is the order of $t \mod{a(t-1)}$, so that $h(a)\ |\ xy = \tau(\beta)$, establishing (ii).

  \item[Theorem 3.5:] Let $a > 1$ be an integer prime to $t$, and let $\beta$ be the smallest $t$-folding number such that $a\ |\ \beta$. Then $h(a) = \tau(\beta)$.

  \item[Proof:] Let $\beta = \frac{(t^{xy}-1)}{(t^{x}-1)}$ so that $\tau(\beta) = xy$, and let $\beta_{0} = \frac{(t^{h(a)}-1)}{(t-1)}$ as in the proof of Theorem 3.4. We show that if $h(a) \neq \tau(\beta)$, then $\beta > \beta_{0}$, contradicting the minimality of $\beta$. Now $h(a)\ |\ \tau(\beta)$, so that if $h(a) \neq \tau(\beta)$, then $h(a) \leq \left( \frac{1}{2} \right) \tau(\beta)$. Thus, since $y \geq 2$, we have
    \begin{equation*}
    \beta > \frac{t^{xy}}{t^{x}} = t^{x(y-1)} \geq t^{xy/2} \geq t^{h(a)} > \beta_{0}
    \end{equation*}

  \item[Remarks:] $\ $
    \begin{enumerate}[\hspace{1cm}(i)]
      \item Theorem 3.5 provides us with an algorithm for finding the smallest $\beta$ such that $a\ |\ \beta$. We first compute $h(a)$, and then test the $t$-folding numbers
        \begin{equation*}
        \frac{t^{h(a)}-1}{t^{x_{1}}-1},\ \frac{t^{h(a)}-1}{t^{x_{2}}-1},\ \cdots
        \end{equation*}
        as $x_{1}, x_{2}, \cdots$ descend through the proper divisors of $h(a)$, to find the first of the sequence which is divisible by $a$. Notice that if $a$ is prime, this must be $\frac{(t^{h(a)}-1)}{(t^{x_{1}}-1)}$. For
        \begin{equation*}
        a\ |\ \frac{t^{h(a)}-1}{t^{x_1}-1} \cdot \frac{t^{x_1}-1}{t-1},\ \text{ but } a \not|\ \frac{t^{x_1}-1}{t-1}
        \end{equation*}
        the latter by the definition of $h(a)$. For example, if $a = 7$ and $t = 3$, then $h(a) = 6$, so the minimum $\beta$ is $\frac{(3^{6}-1)}{(3^{3}-1)} = 28$.
      \item Theorem 3.5 is very useful to the paper-folder. For it affirms that in choosing the smallest folding number $\beta$ having a given odd number $a$ as a factor, we automatically involve ourselves in the smallest number of fold lines -- recall that with $t = 2$, $\tau(\beta) = xy = m+n$.
    \end{enumerate}

\end{description}

Our basic theorems have the following surprising consequence.

\begin{description}

  \item[Corollary 3.6:] If for some $t$,
    \begin{equation*}
    \frac{t^{xy}-1}{t^{x}-1}\ \left|\ \frac{t^{x'y'}-1}{t^{x'}-1} \right.
    \end{equation*}
    with $x,x' \geq 1$, $y,y' \geq 2$, then $xy\ |\ x' y'$.

  \item[Proof:] Set $a = \frac{(t^{xy}-1)}{(t^{x}-1)}$. Since $a \in F_{t}$, it is obviously the smallest $\beta \in F_{t}$ such that $a\ |\ \beta$. Thus by Theorem 3.5, $h(a) = \tau(a) = xy$. But by Theorem 3.4(ii), $h(a)\ |\ \tau(\beta') = x' y'$, where $\beta' = \frac{(t^{x'y'}-1)}{(t^{x}-1)}$.

\end{description}

This result is the natural starting point for a systematic study of divisibility questions. We consider $t$-folding numbers $\beta_{t}\ |\ \beta'_{t}$ as in Corollary 3.6, and ask: For what values of $t$ do we have $\beta_{t}\ |\ \beta'_{t}$? Corollary 3.6 tells us that we may as well assume that $xy\ |\ x' y'$; otherwise, the divisibility condition holds for \emph{no} values of $t$. Thus we assume $xy\ |\ x' y'$. We now give a complete answer to our question, which we enunciate as a Divisibility Theorem.

Set $d = \gcd(xy,x')$, $h = \gcd(x,x') = \gcd(x,d)$. Now $xy = ad$, $x' = bd$, where $\gcd(a,b) = 1$. Thus $ad\ |\ bdy'$, so $a\ | by'$, and hence $a\ |\ y'$. We can now formulate the promised answer.

\begin{description}

  \item[Theorem 3.7 (Divisibility Theorem):] Let $\beta = \frac{(t^{xy}-1)}{(t^{x}-1)}$, $\beta' = \frac{(t^{x'y'}-1)}{(t^{x'}-1)}$ be $t$-folding numbers. Then, assuming $xy\ |\ x'y'$,
    \begin{equation*}
    \beta_{t}\ |\ \beta'_{t}\ \Longleftrightarrow\ \frac{t^{d}-1}{t^{h}-1}\ \left|\ \frac{y'}{a} \right.
    \end{equation*}

\end{description}

We base our proof on 3 easy lemmas. For the first, we allow ourselves to write $\left( \frac{m}{n} \right)\ \left| k \right.$ if $m\ |\ nk$, even if $n \not|\ m$.

\begin{description}

  \item[Lemma 3.8:] $\left( \frac{m}{n} \right)\ \left|\ k\ \Longleftrightarrow\ \left( \frac{m}{g} \right)\ \right|\ k$, where $g = \gcd(m,n)$.

  \item[Proof:] If $m\ |\ gk$, then $m\ |\ nk$. Conversely, let $m = m'g$, $n = n'g$. Then if $m\ |\ nk$, $m'\ |\ n'k$, $m'\ |\ k$ or $\left( \frac{m}{g} \right)\ \left|\ k \right.$.

  \item[Lemma 3.9:] If $g = \gcd(u,v)$, then $\text{lcm} (gu,gv) = uv$.

  \item[Proof:] Let $\ell = \text{lcm}(u,v)$. Then $g\ell = uv$. But $\text{lcm} (gu,gv) = g \text{lcm}(u,v) = g\ell$.

  \item[Lemma 3.10:] Let $d\ |\ q$. Then
    \begin{equation*}
    \frac{t^{d}-1}{t^{x}-1}\ \left|\ \frac{t^{qr}-1}{t^{q}-1}\ \Longleftrightarrow\ \frac{t^{d}-1}{t^{x}-1}\ \right|\ r
    \end{equation*}

  \item[Proof:] Note first that by Lemma 3.8 and Theorem 3.1, we may replace $x$ by $h$ in this statement. To prove it, we set $t^{d} = s$, $q = dk$, and compute the remainder on dividing the polynomial $\frac{(s^{kr}-1)}{(s^{k}-1)}$ by $s-1$. By L'Hopital's Rule, this remainder is
    \begin{equation*} 
    \lim_{s \to 1} \frac{s^{kr}-1}{s^{k}-1} = r
    \end{equation*}

    Thus $\left( \frac{(t^{qr}-1)}{(t^{q}-1)} \right) = Q (t^{d}) (t^{d}-1) + r$, where $Q$ is a polynomial with integer coefficients.

    It follows that
    \begin{equation*} 
    \frac{t^{qr}-1}{t^{q}-1} = (t^{h}-1) Q(t^{d}) \frac{t^{d}-1}{t^{h}-1} + r
    \end{equation*}
    from which the lemma immediately follows.

  \item[Proof of Theorem 3.7:] Set $u = t^{xy}-1$, $v = t^{x'}-1$. Then $g = \gcd(u,v) = t^{d}-1$, by Theorem 3.1. If $c = (t^{x}-1)(t^{x'y'}-1)$, then
    \begin{align*}
    \beta_{t}\ |\ \beta_{t}'\ \Longleftrightarrow\ uv\ |\ c\ \Longleftrightarrow\ (t^{d}-1)(t^{xy}-1)\ |\ c & \text{ and} \\
    (t^{d}-1)(t^{x'}-1)\ |\ c, & \text{ by Lemma 3.9}
    \end{align*}

    \begin{align*}
    & \Longleftrightarrow\ \frac{t^{d}-1}{t^{x}-1}\ \left|\ \frac{t^{x'y'}-1}{t^{xy}-1} \right. && \text{ and }\ \left. \frac{t^{d}-1}{t^{x}-1}\ \right|\ \frac{t^{x'y'}-1}{t^{x'}-1} \\
    & \Longleftrightarrow\ \frac{t^{d}-1}{t^{x}-1}\ \left|\ \frac{x'y'}{xy}             \right. && \text{ and }\ \left. \frac{t^{d}-1}{t^{x}-1}\ \right|\ y',\ \text{ by Lemma 3.10} \\
    & \Longleftrightarrow\ \frac{t^{d}-1}{t^{h}-1}\ \left|\ \frac{x'y}{xy}              \right. && \text{ and }\ \left. \frac{t^{d}-1}{t^{h}-1}\ \right|\ y',\ \text{ by Thm 3.1 / Lem. 3.8} \\
    & \Longleftrightarrow\ \frac{t^{d}-1}{t^{h}-1}\ \left|\ \gcd \left( \frac{x'y}{xy}, y' \right) \right. 
    \end{align*}

    Now $x' y' = bdy'$, $xy = ad$, so $\frac{x'y'}{xy} = \frac{by'}{a}$. Thus $\gcd \left( \left( \frac{x'y'}{xy} \right), y' \right)  =  \gcd \left( b \left( \frac{y'}{a} \right), a \left( \frac{y'}{a} \right) \right) = \frac{y'}{a}$, since $a,b$ are mutually prime. Thus, finally,
    \begin{equation*}
    \beta_{t}\ |\ \beta'_{t} \Longleftrightarrow \frac{t^{d}-1}{t^{h}-1}\ \left|\ \frac{y'}{a} \right.
    \end{equation*}

\end{description}

We draw some consequences from this theorem; we preserve all the hypotheses and notation of the theorem.

\begin{description}

  \item[Corollary 3.11:] Either $d\ |\ x$ and $\beta_{t}\ |\ \beta'_{t}$ for all $t$, or $d \not|\ x$ and $\beta_{t}\ |\ \beta'_{t}$ for only finitely many $t$.

  \item[Corollary 3.12:] $\left[ \frac{(t^{xy}-1)}{(t^{x}-1)} \right]\ \left|\ \left[ \frac{(t^{x'y'}-1)}{(t^{x}-1)} \right]\ \Longleftrightarrow\ \left[ \frac{(t^{xy}-1)}{(t^{h}-1)} \right]\ \right|\ \left[ \frac{(t^{x'y'}-1)}{(t^{x}-1)} \right]$.

\end{description}

We also have the following interesting result.

\begin{description}

  \item[Theorem 3.13:] The following conditions are equivalent:
    \begin{enumerate}[\hspace{1cm}(i)]
      \item $\beta_{t}\ |\ \beta'_{t}$ for all $t$;
      \item $\beta_{t}\ |\ \beta'_{t}$ for infinitely many $t$;
      \item $\beta_{t}\ |\ \beta'_{t}$ as polynomials in $Z[t]$;
      \item $xy\ |\ x' y'$ and $d\ |\ x$.
    \end{enumerate}

  \item[Proof:] We already know the equivalence of (i), (ii), and (iv). Since plainly (iii) $\Rightarrow$ (i), it remains to show that (i) $\Rightarrow$ (iii). This follows from the following lemma.

  \item[Lemma 3.14:] Let $f(t), g(t)$ be polynomials in $Z[t]$ with $g(t)$ monic, and suppose that $g(m)\ |\ f(m)$ for infinitely many integers $m$. Then $g(t)\ |\ f(t)$.

  \item[Proof:] Since $g(t)$ is monic, we may execute the division algorithm in $Z[t]$, obtaining
    \begin{equation*}
    f(t) = q(t)g(t) + r(t), \text{ where } r(t) = 0 \text{ or } \deg r < \deg g
    \end{equation*}

    We will obtain a contradiction from assuming $r(t) \neq 0$. Now 
    \begin{equation*}
    g(m)\ |\ f(m)\ \Longleftrightarrow\ g(m)\ |\ r(m)
    \end{equation*}

    Moreover, since $r(t)$ or $g(t)$ vanishes for only finitely many values of $t$, we may suppose that $g(m)\ |\ r(m)$ for infinitely many $m$, for none of which either $r(t)$ or $g(t)$ vanishes. We consider the function $\frac{r(t)}{g(t)}$. Then $\left| \frac{r(t)}{g(t)} \right|\ \geq 1$ for infinitely many integral values of $t$. On the other hand, since $\deg r < \deg g$, $\frac{r(t)}{g(t)} \to 0$ as $t \to \infty$ and as $t \to -\infty$. Thus $\left| \frac{r(t)}{g(t)} \right|\ \leq \frac{1}{2}$ for $|t|$ sufficiently large. This is the contradiction we seek.

  \item[Examples:] $\ $
    \begin{enumerate}[\hspace{1cm}(1)]
      \item Let $\beta_{t} = \left[ \frac{(t^{3}-1)}{(t-1)} \right]$, $\beta'_{t} = \left[ \frac{(t^{6}-1)}{(t^{2}-1)} \right]$. Here $d = 1$, $x = 1$, so we are in the situation covered by Theorem 3.13. Indeed,
        \begin{align*}
        \beta_{t} & = t^{2} + t + 1, \\
        \beta'_{t} & = t^{4} + t^{2} + 1, \text{ and } \\
        t^{4} + t^{2} + 1 & = (t^{2} + t + 1) (t^{2} - t + 1)
        \end{align*}

        Writing in base $t$, we have $\beta_{t} = 111$, $\beta'_{t} = 10101$, and we know that
        \begin{equation*}
        111\ |\ 10101, \text{ for all } t
        \end{equation*}

        With $t = 2$, we have $7\ |\ 21$; with $t = 3$, we have $13\ |\ 91$; with $t = 4$, we have $21\ |\ 273$; $\cdots$.

      \item Let $\beta_{t} = \left[ \frac{(t^{2}-1)}{(t-1)} \right]$, $\beta'_{t} = \left[ \frac{(t^{6}-1)}{(t^{2}-1)} \right]$. Here $d = 2$, $x = 1$, so $h = 1$; also, $a = b = 1$, $y' = 3$, so
        \begin{equation*}
        \beta_{t}\ \left|\ \beta'_{t}\ \Longleftrightarrow\ \frac{t^{2}-1}{t-1}\ \right|\ 3\ \Longleftrightarrow\ (t+1)\ |\ 3
        \end{equation*}

        Since $t \geq 2$ is an integer, this means that $\beta_{t} + \beta'_{t}\ \Longleftrightarrow\ t = 2$. Thus the assertion $11\ |\ 10101$ is only true in base $2$. There it asserts that $3\ |\ 21$; we verify that $4 \not|\ 91$, $5 \not|\ 273$, $6 \not|\ 651$, $\cdots$.

        These two examples might be (lightheartedly!) summed up by saying that $7\ |\ 21$ always, while $3\ |\ 21$ almost never!

      \item Let $\beta_{t} = \left[ \frac{(t^{24}-1)}{(t^{6}-1)} \right]$, $\beta'_{t} = \left[ \frac{(t^{60y}-1)}{(t^{60}-1)} \right]$. Since we need $24\ |\ 60y'$, we must have $y'$ even, if $\beta_{t}\ |\ \beta'_{t}$ for any $t$. Now $d = 12$, $h = 6$, so $a = 2$ and
        \begin{equation*}
        \beta_{t}\ \left|\ \beta'_{t}\ \Longleftrightarrow\ \frac{t^{12}-1}{t^{6}-1}\ \right|\ \frac{y'}{2}
        \end{equation*}

        We thus see that $\beta_{t} \not|\ \beta_{t}'$ for all $t$, if $y' < 130$. If $y' = 130$, then $\beta_{t}\ |\ \beta_{t}'$ just for $t = 2$. The smallest $y'$ such that $\beta_{3}\ |\ \beta_{3}'$ is given by $y' = 1460$. Arguing the other way, we see that if $y' \leq 2 \cdot 10^{6}$, then the only possible values of $t$ such that $\beta_{t}\ |\ \beta_{t}'$ are in the range $2 \leq t \leq 9$.
    \end{enumerate}

\end{description}

We close this section with a divisibility property, relating to $t$-folding numbers, of a very different kind.

\begin{description}

  \item[Theorem 3.15:] Let $\beta \in F_{t}$. Then
    \begin{enumerate}[\hspace{1cm}(i)]
      \item the order of $t \mod{\beta}$ is $\tau(\beta)$,
      \item $\tau(\beta) + \Phi(\beta)$, where $\Phi$ is the Euler totient function.
    \end{enumerate}

  \item[Proof:] $\ $
    \begin{enumerate}[\hspace{1cm}(i)]
      \item If $\beta = \left[ \frac{(t^{xy}-1)}{(t^{x}-1)} \right]$, then $xy = \tau(\beta)$ and $t^{xy} \equiv 1 \mod{\beta}$. Thus if $h$ is the order of $t$ modulo $\beta$, then $h\ |\ xy$. Suppose $h \sim xy$. Then $h \leq \frac{1}{2} xy$, so (compare the proof of Theorem 3.5):
        \begin{equation*}
        \beta_{t} > \frac{t^{xy}}{t^{x}} = t^{x(y-1)} \geq t^{xy/2} \geq t^{h} > t^{h}-1
        \end{equation*}
        so it is not possible that $t^{h} \equiv 1 \mod{\beta}$. Thus $h = xy$.
      \item Since $\Phi(\beta)$ is the order of the multiplicative group of residues prime to $\beta$, it follows that $\tau^{\Phi(\beta)} = 1 \mod{\beta}$. Thus from (i), $\tau(\beta)\ |\ \Phi(\beta)$.
    \end{enumerate}

\end{description}

It is again interesting to remark that, while $\Phi(\beta)$ certainly depends on $t$, $\tau(\beta)$ does not. Thus we may re-express Theorem 3.15 (ii) as follows. Let $q$ be a fixed positive integer. Consider all the integers $\beta$ of the form $\frac{(t^{q}-1)}{(t^{r}-1)}$, where $t$ is an arbitrary integer $\geq 2$, and $r$ is an arbitrary proper factor of $q$. Then $q\ |\ \Phi(\beta)$ for any such integer $\beta$.

Finally, we note that an even stronger result than Theorem 3.15 is available. For it follows easily from Theorem 4.10 below, that if $\beta = \frac{(t^{xy}-1)}{(t^{x}-1)}$, so that $\tau(\beta) = xy$, then $\tau(\beta)$ is the \emph{quasi-order} of $t \mod{\beta}$ unless $y = 2$, when the quasi-order is $x = \left( \frac{1}{2} \right) \tau(\beta)$ (see Theorem 4.10 for the precise definition of quasi-order). Since we know that $t^{\tau(\beta)} \equiv 1 \mod{\beta}$, this is a finer result than Theorem 3.l5 (i). As to Theorem 3.l5 (ii), we can improve this too by replacing $\Phi(\beta)$ by the \emph{exponent} $\psi(\beta)$, of the multiplicative group of residues prime to $\beta$. For example, if $\beta = 91 = \frac{(3^{6}-1)}{(3^{2}-1)}$, then $\tau(\beta) = 6$, $\Phi(\beta) = 72$, $\psi(\beta) = 12$.




\newpage

\section{The Number Theory Of The General Folding Procedure}
\markboth{4. Theory Of General Folding}{4. Theory Of General Folding}


Let
\begin{equation}
b \begin{vmatrix}
a_{1} & a_{2} & \cdots & a_{r} \\
k_{1} & k_{2} & \cdots & k_{r}
\end{vmatrix} \tag{4.1}
\end{equation}
be a ($2$-)\emph{symbol} as presented in Section 2, where only special cases (e.g. (2.7)) were described explicitly. Thus $b,a_{i}$ ($a_{1} = a$) are odd, $a_{1} < \frac{b}{2}$, and
\begin{equation}
b = a_{i} + 2^{k_i} a_{i+1}, \quad i = 1,2,\cdots,r, \quad a_{r+1} = a_{1}  \tag{4.2}
\end{equation}

At this stage, we do not assume that $\gcd(b,a) = 1$, nor do we assume that the list $a_{1}, a_{2}, \cdots, a_{r}$ is without repeats. Indeed, if $\gcd(b,a) = 1$, we say that the symbol (4.1) is \emph{reduced}; and if there are no repeats, we say that the symbol (4.1) is \emph{contracted}. Thus the folding instructions for regular star polygons of Section 2 are encoded in reduced contracted symbols.

We now study the number theory embodied in the symbols (4.1). Our first observation is a trivial consequence of (4.2).

\begin{description}

  \item[Lemma 4.1:] $\ $
    \begin{enumerate}[\hspace{1cm}(i)]
      \item Given the symbol (4.1), then $\gcd(b, a_{i})$ is independent of $i$ and $a_{i} < \frac{b}{2}$, $1 \leq i \leq r$.
      \item (4.1) is a symbol if and only if
        \begin{equation*}
        qb \begin{vmatrix}
        qa_{1} & qa_{2} & \cdots & qa_{r} \\
        k_{1} & k_{2} & \cdots & k_{r}
        \end{vmatrix} 
        \end{equation*}
        is a symbol, where $q$ is an arbitrary odd number. (This symbol is said to arise by \emph{inflating} (4.1).)
    \end{enumerate}

\end{description}

If we regard (4.2) as a set of $r$ simultaneous linear equations in the ``unknowns'' $a_{1}, a_{2}, \cdots, a_{r}$, then their solution is given by
\begin{equation}
Ba_{i} = bA_{i}, \quad i = 1,2,\cdots,r  \tag{4.3}
\end{equation}
where 
\begin{equation}
B = 2^{k} - (-1)^{r}, \quad k = \Sigma k_{i}  \tag{4.4}
\end{equation}
and 
\begin{equation}
A_{i} = 2^{k-k_{i-1}} - 2^{k-k_{i-1}-k_{i-2}} + \cdots + (-1)^{r_{2}^{k_{i}}} - (-1)^{r}, \quad i = 1,2,\cdots,r  \tag{4.5}
\end{equation}

Here, and subsequently, we find it convenient to think of the suffices $i$ as running over the residue classes $\mod{r}$, so that $k_{0} = k_{r}$, etc. We now announce a portmanteau lemma.

\begin{description}

  \item[Lemma 4.2:] $\ $
    \begin{enumerate}[\hspace{1cm}(i)]
      \item The symbol (4.1) is determined up to contradiction by $b$ and $a_{1}$.
      \item The symbol (4.1) is determined by $b$ and $k_{1}, k_{2}, \cdots, k_{r}$.
      \item $\ $
        \begin{equation}
        B \begin{vmatrix}
        A_{1} & A_{2} & \cdots & A_{r} \\
        k_{1} & k_{2} & \cdots & k_{r} 
        \end{vmatrix}  \tag{4.6}
        \end{equation}
        is a symbol.
      \item If (4.1) is reduced, then it is obtained by reducing (4.6).
      \item $A_{i}$ is independent of $k_{i-1}$.
    \end{enumerate}

  \item[Proof:]
    \begin{enumerate}[\hspace{1cm}(i)]
      \item follows from the definitions. 
      \item is clear from (4.3, 4.4, 4.5). 
      \item follows by observing that
        \begin{equation*}
        b' \begin{vmatrix}
        a'_{1} & a'_{2} & \cdots & a'_{r} \\
        k_{1} & k_{2} & \cdots & k_{r}
        \end{vmatrix}
        \end{equation*}
        is a symbol provided that $b',a'_{i}$ are odd integers satisfying $a'_{i} < \left( \frac{b'}{2} \right)$, and (4.3) with $b',a'_{i}$ replacing $b,a_{i}$. But clearly $A_{i} < (B/2)$, and the equations $Ba_{i}' = b' A_{i}$ are obviously satisfied by $b' = B$, $a_{i}' = A_{i}$. 
      \item If $\gcd(b,a_{i}) = 1$, then from (4.3), there exists a $q$ (odd) such that $B = qb$, $A_{i} = qa_{i}$. 
      \item is a consequence of (4.5).
    \end{enumerate}

\end{description}

We call (4.6) a \emph{canonical} symbol, or, more precisely, the \emph{canonical symbol associated with} $(k_{1}, k_{2}, \cdots, k_{r})$. As we see, every symbol is obtained by reducing a canonical symbol and (perhaps) then inflating the resulting symbol. Thus while symbols do in a sense generalize the notion of a fraction, we see that there are two kinds of standard symbol: the reduced contracted symbols, and the canonical symbols.

Of course, canonical symbols need not be reduced. The following lemma deals with the question of whether they are contracted.

\begin{description}

  \item[Lemma 4.3:] $\ $
    \begin{enumerate}[\hspace{1cm}(i)]
      \item If the sequence $k_{1}, k_{2}, \cdots, k_{r}$ is obtained by repeating a subsequence $k_{1}, k_{2}, \cdots, k_{s}$, then the sequence $a_{1}, a_{2}, \cdots, a_{r}$ is obtained by repeating the subsequence $a_{1}, a_{2}, \cdots, a_{s}$.
      \item The canonical symbol (4.6) is contracted if and only if the symbol (4.1) is contracted.
    \end{enumerate}

  \item[Proof:] $\ $
    \begin{enumerate}[\hspace{1cm}(i)]
      \item We suppose that there exists an $s$ such that $s\ |\ r$, $s \neq r$, and $k_{s+i} = k_{i}$ for all $i$. It is then plain from (4.5) that $A_{s+i} = A_{i}$ for all $i$, so that by (4.3), $a_{s+i} = a_{i}$ for all $i$.
      \item Using (i), we observe that the symbol (4.1) (and the symbol (4.6)) involves a repeat if and only if the sequence $k_{1}, k_{2}, \cdots, k_{r}$ involves a repeat.
    \end{enumerate}

  \item[Remarks:] $\ $
    \begin{enumerate}[\hspace{1cm}(i)]
      \item When we talk of the symbol (4.1) \emph{repeating}, we require only that for some $s,t$ with $1 \leq s$, $t \leq r$, and $s \neq t$, $a_{s} = a_{t}$. It is easy to see that if this happens, then in fact $a_{1}$ must repeat, say with $a_{n} = a_{1}$, and then the entire symbol repeats with ``period'' $(n-1)$. However, in Lemma 4.3, we do not merely require that some $k_{i}$ is repeated, but that the sequence repeats in the sense described.
      \item We particularly note that Lemma 4.2 (iii),(iv) answers our question as to what regular star $\left\{ \frac{b}{a} \right\}$-gons are obtained by the folding instruction $(k_{1}, k_{2}, \cdots, k_{r})$, which is of course assumed to be non-repeating. Thus we calculate $B$, $A_{1}, \cdots, A_{r}$ from (4.4), (4.5) and then reduce by dividing by $d = \gcd(B,A_{i})$, which is known to be independent of $i$. If $B = bd$, $A_{i} = a_{i} d$, then we can construct by the given folding instruction, precisely the $\left\{ \frac{b}{a_{1}} \right\}$-gons. This, of course, assumes we are prepared to use the FAT-algorithm on either the top or the bottom of the tape.
    \end{enumerate}

\end{description}

We now move towards our principal number-theoretical result. (All integers referred to are positive.)

\begin{description}

  \item[Lemma 4.4:] Let $b$ be odd $\geq 3$, and let $S$ be the set of odd integers $a$ with $a < \frac{b}{2}$. Then the function $\psi$ on $S$, given by $\psi(a) = a'$, where $b = a + 2^{k} a'$, $a'$ odd, is a permutation of $S$.

  \item[Proof:] Since $k \geq 1$, it is obvious that $a' < \frac{b}{2}$. We show that $\psi$ is a permutation of $S$ by defining an inverse $\varphi$. Thus, given $a' \in S$, let $k$ be maximal with $2^{k} a' < b$. Then $k \geq 1$, so that if $a = b - 2^{k} a'$, then $a$ is odd and $a < \frac{b}{2}$. Of course, we define $\varphi(a') = a$.

\end{description}

This lemma shows that, as claimed, a symbol always exists for given odd numbers $b,a$ with $a < \frac{b}{2}$; and it is, of course, unique up to contraction.

\begin{description}

  \item[Theorem 4.5 (Quasi-Order Theorem in base 2):] Let (4.1) be a reduced contracted symbol, and let $k = \Sigma k_{i}$. Then $2^{k} \equiv (-1)^{r} \mod{b}$. Indeed, $k$ is the quasi-order of $2 \mod{b}$; that is, the smallest positive integer $n$ such that $2^{n} \equiv \pm 1 \mod{b}$.

  \item[Proof:] It is obvious from (4.3), since $\gcd(b,a_{i}) = 1$, that $b\ |\ 2^{k}-(-l)^{r}$. To prove that $k$ is the quasi-order, we effectively employ the $\varphi$ function introduced in the proof of Lemma 4.4. Thus consider the sequence of $(k+1)$ positive integers $< \frac{b}{2}$ (see (4.2)).
    \begin{align*}
    & a_{r}, 2a_{r}, \cdots, 2^{k_{r-1}-1} a_{r}, a_{r-1}, 2a_{r-1}, \cdots, a_{i+1}, 2a_{i+1}, \cdots, 2^{k_{i}-1} a_{i+1}, \\
    & a_{i}, 2a_{i}, \cdots, a_{1}, 2a_{1}, \cdots, a^{k_{r}-1} a_{1}, a_{r} 
    \end{align*}

    Call this sequence $s_{1}, s_{2}, \cdots, s_{k+1}$. Then plainly,
    \begin{equation*}
    s_{j+1} \equiv \pm 2s_{j} \mod{b}
    \end{equation*}
    the minus sign obtaining precisely when $s_{j+1}$ is some $a_{i}$. (This provides an alternative proof that $2^{k} \equiv (-1)^{r} \mod{b}$.) 

    Thus, if $2^{\ell} \equiv \pm 1 \mod{b}$ for some positive $\ell < k$, we would have $s_{\ell+1} \equiv \pm a_{r} \mod{b}$. We show that this is impossible. For if $s_{\ell+1}$ $\equiv$ a$_{r} \mod{b}$, then since $s_{\ell+1},a_{r} < \frac{b}{2}$, we have $s_{\ell+1} = a_{r}$. This is impossible if $s_{\ell+1} = 2^{m} a_{i}$, $m \geq 1$, since $a_{r}$ is odd; but it is also impossible if $s_{\ell+1} = a_{i}$, since $i \not|\ r$ and the symbol is contracted. On the other hand, if $s_{\ell+1} \equiv -a_{r} \mod{b}$, then $b\ |\ (s_{\ell+1} + a_{r})$, which is impossible since $0 < s_{\ell+1} + a_{r} < b$. This proves the theorem.

\end{description}

There are many interesting consequences of this theorem. We give one of significance for paper-folding.

\begin{description}

  \item[Corollary 4.6:] $\ $
    \begin{enumerate}[\hspace{1cm}(i)]
      \item Let $k \geq 3$. Then every contracted reduced symbol (4.1) with $b = 2^{k}-1$ is canonical.
      \item Let $k \geq 1$. Then every contracted reduced symbol (4.1) with $b = 2^{k}+1$ is canonical.
    \end{enumerate}

  \item[Proof:] $\ $
    \begin{enumerate}[\hspace{1cm}(i)]
      \item Obviously, the quasi-order of $2 \mod{(2^{k}-1)}$ is $k$ if $k \geq 3$. Thus if $\gcd (2^{k}-1, a_{1}) = 1$, the contracted reduced symbol we obtain, namely
        \begin{equation}
        2^{k}-1 \begin{vmatrix}
        a_{1} & a_{2} & \cdots & a_{r} \\
        k_{1} & k_{2} & \cdots & k_{r}
        \end{vmatrix}  \tag{4.7}
        \end{equation}
        has by Theorem 4.5, $\Sigma k_{i} = k$ (and $r$ even). By Lemma 4.2 (ii,iii), it follows that this symbol is canonical. 
      \item is proved similarly.
    \end{enumerate}

  \item[Remark:] We need $k \geq 3$ in (i), since $b \geq 3$, and if $k = 2$, then the quasi-order of $2 \mod{(2^{k}-1)}$ is $1$, not $2$.

  \item[Proposition 4.7:] Let $k = \Sigma k_{i}$.
    \begin{enumerate}[\hspace{1cm}(i)]
      \item Let $k \geq 2$. Then if (4.7) is a symbol, there is a symbol
        \begin{equation}
        2^{k+\ell}-1 \begin{vmatrix}
        a_{1} & a_{2}^{(\ell)} & \cdots & a_{r-1}^{(\ell)} & a_{r}^{(\ell)} \\
        k_{1} & k_{2}          & \cdots & k_{r-1}          & k_{r}+\ell 
        \end{vmatrix}  \tag{4.8-}
        \end{equation}
        for some $a_{2}^{(\ell)}, \cdots, ar^{(\ell)}$, for any $\ell \geq 0$.
      \item Let $k \geq 1$. Then if
        \begin{equation*}
        2^{k}+1 \begin{vmatrix}
        a_{1} & a_{2} & \cdots & a_{r} \\
        k_{1} & k_{2} & \cdots & k_{r}
        \end{vmatrix} 
        \end{equation*}
        is a symbol, there is a symbol
        \begin{equation}
        2^{k+\ell}+1 \begin{vmatrix}
        a_{1} & a_{2}^{(\ell)} & \cdots & a_{r-1}^{(\ell)} & a_{r}^{(\ell)} \\
        k_{1} & k_{2}          & \cdots & k_{r-1}          & k_{r}+\ell 
        \end{vmatrix}  \tag{4.8+}
        \end{equation}
        for some $a_{2}^{(\ell)}, \cdots, a_{r}^{(\ell)}$, for any $\ell \geq 0$.
    \end{enumerate}

  \item[Proof:] $\ $
    \begin{enumerate}[\hspace{1cm}(i)]
      \item If (4.7) is a symbol, then by (4.3),
        \begin{equation}
        (2^{k}-(-1)^{r}) a_{1} = (2^{k}-1) A_{1}
        \end{equation}

        The assumption that $r$ is odd leads to the conclusion $2^{k}-1\ |\ a_{1}$ contradicting $a_{1} < \left( \frac{1}{2} \right) (2^{k}-1)$. Thus $r$ is even, $a_{1} = A_{1}$, and (4.7) is canonical. Recalling Lemma 4.2 (v), it follows that the canonical symbol associated with $k_{1}, k_{2}, \cdots, k_{r-1}, k_{r}+\ell$ has the same $A_{1}$ as the canonical symbol associated with $k_{1}, k_{2}, \cdots, k_{r-1}, k_{r}$. This proves (i).

      \item is proved similarly.
    \end{enumerate}

  \item[Corollary 4.8:] $\ $
    \begin{enumerate}[\hspace{1cm}(i)]
      \item Let $\gcd(2^{k}-1,a) = 1$, $k \geq 3$. If the regular $\left\{ \frac{(2^{k}-1)}{a} \right\}$-star polygon is folded by the procedure $(k_{1}, k_{2}, \cdots, k_{r})$, then the regular $\left\{ \frac{(2^{k+\ell}-1)}{a} \right\}$-star polygon is folded by the procedure $(k_{1}, k_{2}, \cdots, k_{r-1}, k_{r}+\ell)$.
      \item Let $\gcd(2^{k}+1,a) = 1$, $k \geq 1$. If the regular $\left\{ \frac{(2^{k}+1)}{a} \right\}$-star polygon is folded by the procedure $(k_{1}, k_{2}, \cdots, k_{r})$, then the regular $\left\{ \frac{(2^{k+\ell}+1)}{a} \right\}$-star polygon is folded by the procedure $(k_{1}, k_{2}, \cdots, k_{r-1}, k_{r}+\ell)$.
    \end{enumerate}

  \item[Proof:] This is an immediate consequence of Proposition 4.7 and Corollary 4.6.

  \item[Remark:] For some $\ell$, $\frac{(2^{k+\ell} \pm 1)}{a}$ may fail to be a reduced fraction. Moreover, it may happen (for at most one value of $\ell$ -- see [5]) that the sequence $k_{1}, k_{2}, \cdots, k_{r-1}, k_{r}+\ell$ repeats. In this latter very special case, we will certainly know from Theorem 4.5 that $\frac{(2^{k+\ell} \pm 1)}{a}$ is not reduced, and we will be able to contract the folding instructions.

\end{description}

We now consider how to generalize the notion of a $2$-symbol to that of a $t$-symbol to obtain a generalization of Theorem 4.5. It turns out to be easier to generalize the $\varphi$-algorithm than the $\psi$-algorithm in Lemma 4.4.

In order to accommodate both the $\psi$-algorithm and the $\varphi$-algorithm, it is preferable to revise the format of a general $2$-symbol, so that it appears as
\begin{equation}
b \begin{vmatrix}
a_{1} &       & a_{2} & \cdot & \cdot & \cdot & a_{r} &       & a_{1} \\
      & k_{1} &       & k_{2} & \cdot & \cdot & \cdot & k_{r} &       
\end{vmatrix}  \tag{4.9}
\end{equation}

Then, if we are emphasizing the $\varphi$-algorithm, we would rewrite (4.9) as
\begin{equation}
b \begin{vmatrix}
c_{1} &          & c_{r} & \cdot & \cdot & \cdot    & c_{2} &          & c_{1} \\
      & \ell_{r} &       &       &       & \ell_{2} &       & \ell_{1} &       
\end{vmatrix}  \tag{4.10}
\end{equation}
so that $b = c_{i+1} + 2^{\ell_i} c_{i}$, and $\varphi (c_{i}) = c_{i+1}$, $i = 1,2,\cdots,r$ ($c_{r+l} = c_{1}$). Thus the data for the $\varphi$-algorithm with $t = 2$ are the following: we take $b$ odd, and consider the set $S$ of odd positive integers $c < \frac{b}{2}$. For such $c \in S$, let $\ell$ be maximal subject to $2^{\ell} c < b$, and define $c' = \varphi(c)$ by the rule
\begin{equation}
b = c' + 2^{\ell} c  \tag{4.11}
\end{equation}

To generalize from $t = 2$ to the case of a general positive integer $t \geq 2$, we find we should proceed as follows:
\begin{enumerate}[\hspace{1cm}(1)]
  \item The condition that $b$ is odd is generalized to the condition that $b$ \emph{is prime to} $t$.
  \item The condition that $c$ is odd is generalized to the condition that $t \not|\ c$.
  \item The condition $e < \frac{b}{2}$ is generalized to the condition $c < \frac{b}{2}$ (this was the most surprising).
\end{enumerate}

Having generalized the set $S$, it remains to generalize the $\varphi$-algorithm. We first generalize the characterization of $\ell$ above.

\begin{description}
  \item{(4)} The condition that $\ell$ is \emph{maximal} subject to $2^{\ell} c < b$ is equivalent to the condition that $\ell$ is \emph{minimal} subject to $2^{\ell} c > \frac{b}{2}$, and may therefore be generalized to the condition that $\ell$ is \emph{minimal subject to} $t^{\ell} c > \frac{b}{2}$; note that, as with $t = 2$, it follows that $\ell \geq 1$.
\end{description}

We are now ready to generalize $\varphi$. We observe that in (4.11), the multiple of $b$ nearest to $2^{\ell} c$ is $b$ itself; in fact, $b > 2^{\ell} c$, and $c' = \varphi(c)$ is the distance from $2^{\ell} c$ to the nearest multiple of $b$, namely to $b$ itself. In general, we define $\varphi(c)$ to be the distance from $t^{\ell} c$ to the nearest multiple of $b$, where $\ell$ is defined as in (4) above. We immediately note that this definition makes sense since we cannot have $t^{\ell} c = \frac{nb}{2}$, for this would imply $b\ |\ 2c$, whereas $2c < b$. Thus if $\varphi(c) = c'$, then
\begin{equation}
qb + (-1)^{\epsilon} c' = t^{\ell} c, \quad \epsilon = 0 \text{ or } 1  \tag{4.12}
\end{equation}
with $q \geq 1$. (Note that if $t = 2$, then $q = 1$ and $\epsilon = 1$.)

We remark that the minimality of $\ell$ implies that $t^{\ell} c < \frac{tb}{2}$, so that in (4.12), $q \leq \frac{t}{2}$ (and $q < \frac{t}{2}$ if $\epsilon = 0$). In fact, the restriction on $q$ in (4.12) is precisely
\begin{equation}
\begin{cases}
& 1 \geq q \geq \frac{t-1}{2} \text{ if $t$ is odd} \\
& 1 \geq q \geq \frac{t}{2} \text{ if $t$ is even and $\epsilon = 1$} \\
& 1 \geq q \geq \frac{t-2}{2} \text{ if $t$ is even and $\epsilon = 0$}
\end{cases}  \tag{4.13}
\end{equation}

The definition of $\varphi$ makes it obvious that $0 < c' < \frac{b}{2}$. Since $q \leq \frac{t}{2}$, it follows that $t \not|\ q$, so $t \not|\ qb$, and hence from (4.12), $t \not|\ c'$. Thus $c' \in S$, and $\varphi$ is a function
\begin{equation*}
\varphi\ :\ S \to S
\end{equation*}

\begin{description}

  \item[Theorem 4.9:] The function $\varphi$ is a permutation of $S$.

  \item[Proof:] First let $t$ be odd, and consider for a given $c' \in S$, the set $Q = \{ qb + (-l)^{\epsilon} c$, $1 \leq q \leq \frac{(t-1)}{2}$, $\epsilon = 0,1 \}$. We claim that exactly one of the $(t-1)$ numbers in $Q$ is divisible by $t$. Certainly, exactly one of the $t$ numbers $sb + c'$, $-\frac{(t-1)}{2} \leq s \leq \frac{(t-1)}{2}$, is divisible by $t$; and this number cannot be given by $s = 0$, since $t \not|\ c'$. If this number has $s$ positive, set $s = q$, so that $t\ |\ (qb + c')$, $1 \leq q \leq \frac{(t-1)}{2}$; if it has $s$ negative, set $s = -q$, so that $t\ |\ (qb - c')$, $1 \leq q \leq \frac{(t-1)}{2}$. Thus we have found $q,\epsilon$ satisfying $1 \leq q \leq \frac{(t-1)}{2}$, $\epsilon = 0,1$, such that
    \begin{equation*}
    t^{\ell} c = qb + (-1)^{\epsilon} c'
    \end{equation*}
    with $\ell \geq 1$ and $t \not|\ c$. It remains to show that $t^{\ell-1} c < \frac{b}{2}$, since this guarantees both the minimality of $\ell$ and the fact that $c \in S$. But $t^{\ell} c < \left( \frac{(t-1)}{2} \right) b + \frac{b}{2} = \frac{tb}{2}$, establishing the required inequality. Thus, $c \in S$ and $\varphi(c) = c'$.

    A small modification is required if $t$ is even. We now consider the set
    \begin{equation*}
    Q' = \left\{ qb+c',\ 1 \leq q \leq \frac{t}{2}-1;\ \ qb-c',\ 1 \leq q \leq \frac{t}{2} \right\}
    \end{equation*}

    Once again, we argue that exactly one of the $(t-1)$ numbers in $Q'$ is divisible by $t$, and we choose it to get the equation
    \begin{equation*}
    t^{\ell} c = qb + (-1)^{\epsilon} c'
    \end{equation*}

    where $\ell > 1$, $t \not|\ c$, and $1 \leq q \leq \frac{t}{2}$ if $\epsilon = 1$, while $1 \leq q \leq \frac{t}{2}-1$ if $\epsilon = 0$. We again argue that $t^{\ell-1} c < \frac{b}{2}$. For if $\epsilon = 0$, then $t^{\ell} c < \left( \frac{t}{2}-1 \right) b + \frac{b}{2} < \frac{tb}{2}$; and if $\epsilon = 1$, then $t^{\ell} c < \left( \frac{t}{2} \right) b$. This completes the proof of the theorem, for we have shown $\varphi\ :\ S \rightarrow S$ to be a surjective function from a finite set to itself, and hence a permutation.

  \item[Remarks:] $\ $
    \begin{enumerate}[\hspace{1cm}(i)]
      \item The proof of the theorem implicitly contains the definition of the inverse algorithm $\psi\ :\ S \rightarrow S$. Notice that this $\psi$ generalizes the algorithm $\psi$ of Lemma 4.4, since with $t = 2$ we must take $\epsilon = 1$, $q = 1$, so that $\psi(a) = a'$, where $2^{\ell} a' = b-a$.
      \item It will be convenient to enlarge $S$ to allow $a \leq \frac{b}{2}$; and when $b$ is even, to define $\varphi \left( \frac{b}{2} \right) = \frac{b}{2}$, $\psi \left( \frac{b}{2} \right) = \frac{b}{2}$. Thus the choice for (4.12) which we make in this singular case is
        \begin{equation*}
        \frac{t-1}{2} b + \frac{b}{2} = t \frac{b}{2}
        \end{equation*}
        which accords with the restriction $1 \leq q \leq \frac{(t-1)}{2}$ in the case that $t$ is odd. We write $\bar{S}$ for $S \cup \left\{ \frac{b}{2} \right\}$, for $b$ even; and may write $\bar{S}$ for $S$ if $b$ is odd.
    \end{enumerate}

\end{description}

Our $t$-symbol is now more elaborate than in the case $t = 2$, and contains 4 rows, though only the second and third rows are relevant to the quasi-order theorem. Thus let $b,t$ be mutually prime, and let $a_{1} \in S$. We iterate the $\psi$-algorithm, getting
\begin{equation}
t^{k_i} a_{i+1} = q_{i} b + (-1)^{\epsilon_i} a_{i}, \quad i = 1,2,\cdots,r, \quad a_{r+1} = a_{1}  \tag{4.14}
\end{equation}
where the positive integers $q_{i}$ are restricted as in the proof of Theorem 4.9, and record (4.14) in the symbol (in its new format)

\begin{tabular}{l|ccccccccc|r}
$b$ & $a_{1}$ &                & $a_{2}$ &     $\cdot$    & $\cdot$ & $\cdot$ & $a_{r}$ &                & $a_{1}$ &     \\
    &         & $k_{1}$        &         & $k_{2}$        & $\cdot$ & $\cdot$ & $\cdot$ & $k_{r}$        &         &     \\
    &         & $\epsilon_{1}$ &         & $\epsilon_{2}$ & $\cdot$ & $\cdot$ & $\cdot$ & $\epsilon_{r}$ &         &     \\
    &         & $q_{1}$        &         & $q_{2}$        & $\cdot$ & $\cdot$ & $\cdot$ & $q_{r}$        &         & $t$ 
\end{tabular} 
\begin{equation}
k = \Sigma k_{i},\ \epsilon = \Sigma \epsilon_{i}  \tag{4.15}
\end{equation}

Note again that if $t = 2$, then all $\epsilon_{i} = 1$ and all $q_{i} = 1$, so that the last two rows of (4.15) may be omitted, and $\epsilon = r$.

We may clearly introduce the notions \emph{contracted, reduced} for $t$-symbols. We will not copy all the earlier material on $2$-symbols in this more general context, but will proceed immediately to enunciate the general quasi-order theorem.

\begin{description}

  \item[Theorem 4.10:] Let the symbol (4.15) be contracted and reduced. Then $t^{k} \equiv (-1)^{\epsilon} \mod{b}$, and $k$ is the quasi-order of $t \mod{b}$.

  \item[Proof:] We argue just as in the case $t = 2$ from the $\varphi$-algorithm. We leave the details to the reader.

  \item[Remark:] By including the case $b = 2$, $a_{1} = 1$, we make our quasi-order theorem completely general. The only contracted reduced symbol for odd $t$, with $a = \frac{b}{2}$, is
    \begin{equation*}
    \begin{array}{c}
    2 \\
    \ \\
    \ \\
    \ 
    \end{array}
    \begin{vmatrix}
    1 &   & 1 \\
      & 1 &   \\
      & 0 &   \\
      & \frac{t-1}{2} & \\
    \end{vmatrix}
    \begin{array}{c}
    \ \\
    \ \\
    \ \\
    t
    \end{array}
    \end{equation*}
    yielding the true (if unsensational) fact that if $t$ is odd, its quasi-order $\mod{2}$ is $1$. Let us close with a more interesting numerical example.

  \item[Example:] Let us take $t = 3$. The $\psi$-algorithm then proceeds by deciding, at each stage, whether $b-a$ or $b+a$ is divisible by $3$ (remember that both $a$ and $b$ are prime to $3$), and recording $\epsilon = 1$ in the first case, $\epsilon = 0$ in the second. It is unnecessary to record the value of $q$, since with $t = 3$, it is always $1$. Thus, for example, we have symbols
\end{description}
{\tiny
\begin{tabular}{l|lllllll|lllllll|l}
37&1 &4 &11 &16 &7 &10 &1 &2 &13 &8 &5 &14 &17 &2 &   \\
  &\quad 2&\quad 1&\quad 1&\quad 1&\quad 1&\quad 3& &\quad 1&\quad 1&\quad 2&\quad 1&\quad 1&\quad 3& &  \\
  &\quad 1&\quad 1&\quad 0&\quad 1&\quad 1&\quad 1& &\quad 0&\quad 1&\quad 0&\quad 0&\quad 0&\quad 0& & 3
\end{tabular}
}
\begin{description}
  \item[] $\ $

    each symbol shows that the quasi-order of $3 \mod{37}$ is $9$, and that $3^{9} \equiv -1 \mod{37}$. Despite some valid remarks in [2] on the relative complexity of the $\psi$-algorithm, we unrepentantly claim that it is more attractive to intelligent human beings than the $\varphi$-algorithm -- though of course the $\varphi$-algorithm is certainly the more computer-friendly. We are not surprised at this dichotomy.
\end{description}




\newpage

\section{Folding All Star Polygons}
\markboth{5. Folding All Star Polygons}{5. Folding All Star Polygons}


In this section, we give an algorithm for introducing \emph{secondary fold lines} (if necessary) into our tape, in order to be able by means of the FAT-algorithm, to construct any given regular $\left\{ \frac{b}{a} \right\}$-gon.

We assume that we can produce tape suitable for constructing (by the FAT-algorithm) any regular $\left\{ \frac{b}{a} \right\}$-gon \emph{with $a,b$ both odd}, and $a < \frac{b}{2}$. We first consider the case when we wish to construct a regular star $\left\{ \frac{b}{a} \right\}$-gon, where $b$ is odd and $a$ is even. Then $a = 2^{k} a'$ with $a'$ odd, and we assert

\begin{description}

  \item[Proposition 5.1:] The tape prepared for folding the regular star $\left\{ \frac{b}{a'} \right\}$-gon has a crease line making an angle of $\left( \frac{a}{b} \right) \pi$ with the forward direction of the top of the tape, and a second crease line under it making an angle of $\left( \frac{a}{b} \right) \pi$ with it.

\end{description}
\begin{center}
%TeXCAD Picture [fig5.1.pic]. Options:
%\grade{\on}
%\emlines{\off}
%\epic{\off}
%\beziermacro{\on}
%\reduce{\on}
%\snapping{\off}
%\quality{8.00}
%\graddiff{0.01}
%\snapasp{1}
%\zoom{4.0000}
\unitlength .95mm % = 2.85pt
\linethickness{0.4pt}
\ifx\plotpoint\undefined\newsavebox{\plotpoint}\fi % GNUPLOT compatibility
\begin{picture}(126.25,34)(0,0)
\put(3.25,34){\line(1,0){119.75}}
\put(126.25,16.25){\line(-1,0){124.25}}
%\dashline{1}(5.75,16.25)(67,33.75)
\multiput(5.68,16.18)(.11779,.03365){8}{\line(1,0){.11779}}
\multiput(7.56,16.72)(.11779,.03365){8}{\line(1,0){.11779}}
\multiput(9.45,17.26)(.11779,.03365){8}{\line(1,0){.11779}}
\multiput(11.33,17.8)(.11779,.03365){8}{\line(1,0){.11779}}
\multiput(13.22,18.33)(.11779,.03365){8}{\line(1,0){.11779}}
\multiput(15.1,18.87)(.11779,.03365){8}{\line(1,0){.11779}}
\multiput(16.99,19.41)(.11779,.03365){8}{\line(1,0){.11779}}
\multiput(18.87,19.95)(.11779,.03365){8}{\line(1,0){.11779}}
\multiput(20.76,20.49)(.11779,.03365){8}{\line(1,0){.11779}}
\multiput(22.64,21.03)(.11779,.03365){8}{\line(1,0){.11779}}
\multiput(24.53,21.56)(.11779,.03365){8}{\line(1,0){.11779}}
\multiput(26.41,22.1)(.11779,.03365){8}{\line(1,0){.11779}}
\multiput(28.3,22.64)(.11779,.03365){8}{\line(1,0){.11779}}
\multiput(30.18,23.18)(.11779,.03365){8}{\line(1,0){.11779}}
\multiput(32.06,23.72)(.11779,.03365){8}{\line(1,0){.11779}}
\multiput(33.95,24.26)(.11779,.03365){8}{\line(1,0){.11779}}
\multiput(35.83,24.8)(.11779,.03365){8}{\line(1,0){.11779}}
\multiput(37.72,25.33)(.11779,.03365){8}{\line(1,0){.11779}}
\multiput(39.6,25.87)(.11779,.03365){8}{\line(1,0){.11779}}
\multiput(41.49,26.41)(.11779,.03365){8}{\line(1,0){.11779}}
\multiput(43.37,26.95)(.11779,.03365){8}{\line(1,0){.11779}}
\multiput(45.26,27.49)(.11779,.03365){8}{\line(1,0){.11779}}
\multiput(47.14,28.03)(.11779,.03365){8}{\line(1,0){.11779}}
\multiput(49.03,28.56)(.11779,.03365){8}{\line(1,0){.11779}}
\multiput(50.91,29.1)(.11779,.03365){8}{\line(1,0){.11779}}
\multiput(52.8,29.64)(.11779,.03365){8}{\line(1,0){.11779}}
\multiput(54.68,30.18)(.11779,.03365){8}{\line(1,0){.11779}}
\multiput(56.56,30.72)(.11779,.03365){8}{\line(1,0){.11779}}
\multiput(58.45,31.26)(.11779,.03365){8}{\line(1,0){.11779}}
\multiput(60.33,31.8)(.11779,.03365){8}{\line(1,0){.11779}}
\multiput(62.22,32.33)(.11779,.03365){8}{\line(1,0){.11779}}
\multiput(64.1,32.87)(.11779,.03365){8}{\line(1,0){.11779}}
\multiput(65.99,33.41)(.11779,.03365){8}{\line(1,0){.11779}}
%\end
%\dashline{1}(67,33.75)(95.25,16.75)
\multiput(66.93,33.68)(.055392,-.033333){15}{\line(1,0){.055392}}
\multiput(68.59,32.68)(.055392,-.033333){15}{\line(1,0){.055392}}
\multiput(70.25,31.68)(.055392,-.033333){15}{\line(1,0){.055392}}
\multiput(71.92,30.68)(.055392,-.033333){15}{\line(1,0){.055392}}
\multiput(73.58,29.68)(.055392,-.033333){15}{\line(1,0){.055392}}
\multiput(75.24,28.68)(.055392,-.033333){15}{\line(1,0){.055392}}
\multiput(76.9,27.68)(.055392,-.033333){15}{\line(1,0){.055392}}
\multiput(78.56,26.68)(.055392,-.033333){15}{\line(1,0){.055392}}
\multiput(80.22,25.68)(.055392,-.033333){15}{\line(1,0){.055392}}
\multiput(81.89,24.68)(.055392,-.033333){15}{\line(1,0){.055392}}
\multiput(83.55,23.68)(.055392,-.033333){15}{\line(1,0){.055392}}
\multiput(85.21,22.68)(.055392,-.033333){15}{\line(1,0){.055392}}
\multiput(86.87,21.68)(.055392,-.033333){15}{\line(1,0){.055392}}
\multiput(88.53,20.68)(.055392,-.033333){15}{\line(1,0){.055392}}
\multiput(90.19,19.68)(.055392,-.033333){15}{\line(1,0){.055392}}
\multiput(91.86,18.68)(.055392,-.033333){15}{\line(1,0){.055392}}
\multiput(93.52,17.68)(.055392,-.033333){15}{\line(1,0){.055392}}
%\end
%\dashline{1}(94.5,17.25)(96.25,16.25)
\multiput(94.43,17.18)(.05833,-.03333){10}{\line(1,0){.05833}}
\multiput(95.6,16.51)(.05833,-.03333){10}{\line(1,0){.05833}}
%\end
%\dashline{1}(67.25,32.75)(71.25,16.5)
\multiput(67.18,32.68)(.03175,-.12897){7}{\line(0,-1){.12897}}
\multiput(67.62,30.87)(.03175,-.12897){7}{\line(0,-1){.12897}}
\multiput(68.07,29.07)(.03175,-.12897){7}{\line(0,-1){.12897}}
\multiput(68.51,27.26)(.03175,-.12897){7}{\line(0,-1){.12897}}
\multiput(68.96,25.46)(.03175,-.12897){7}{\line(0,-1){.12897}}
\multiput(69.4,23.65)(.03175,-.12897){7}{\line(0,-1){.12897}}
\multiput(69.85,21.85)(.03175,-.12897){7}{\line(0,-1){.12897}}
\multiput(70.29,20.04)(.03175,-.12897){7}{\line(0,-1){.12897}}
\multiput(70.74,18.24)(.03175,-.12897){7}{\line(0,-1){.12897}}
%\end
\put(61.5,28.25){\makebox(0,0)[cc]{$6 \frac{\pi}{13}$}}
\put(73,25.25){\makebox(0,0)[cc]{$3 \frac{\pi}{13}$}}
\put(81.5,30){\makebox(0,0)[cc]{$3 \frac{\pi}{13}$}}
\put(60,11){\makebox(0,0)[cc]{Illustration of Proposition 5.1}}
\put(67.25,3.5){\makebox(0,0)[cc]{Figure 5.1}}
\end{picture}
\end{center}
    \begin{equation*}
    a=6, \quad a'=3, \quad k=1, \quad b=13
    \end{equation*}
\begin{description}
  \item[Proof:] Our folding algorithm always starts with an \emph{obtuse} angle, and proceeds by repeated bisection. Thus the angle of $\pi \frac{a'}{b}$ on the tape was created by successive bisection, by folding, of an angle $\frac{2^{\ell} \pi a'}{b}$ where
    \begin{equation*}
    \frac{2^{\ell} \pi a'}{b} > \frac{\pi}{2}
    \end{equation*}

Since $\left( \frac{2^{k} \pi a'}{b} \right) < \frac{\pi}{2}$, it follows that $k < \ell$, so that the bisection process must have produced the stated configuration.

\end{description}

Thus, in this case, we can simply use the FAT-algorithm, on the appropriate crease lines of the $\left\{ \frac{b}{a'} \right\}$-tape, to construct the desired $\left\{ \frac{b}{a} \right\}$-gon; no secondary folds are needed.

We now consider the problem of constructing a regular star $\left\{ \frac{b}{a} \right\}$-gon, where $b$ is even and $a$ is odd. Then $b = 2^{k} b'$, with $b'$ odd. Let $a$ be written in base $2$ as
\begin{equation}
a = \epsilon_{0} \epsilon_{1} \cdots \epsilon_{n},  \quad  \epsilon_{0} = 1,  \quad  \epsilon_{i} = 0 \text{ or } 1,  \quad  1 \leq i \leq n \text{ (in fact, $\epsilon_{n} = 1$ )}  \tag{5.1}
\end{equation}

We distinguish two possibilities. First, let $b' = 1$, so that $b = 2^{k}$. Since $2^{n} \leq a \leq 2^{k-1}$, we have $n \leq k-2$. We initiate our algorithm with tape folded to produce a crease line making an angle of $\frac{\pi}{2^{k-n}}$ with the forward direction of the top of the tape, with a second crease line under it making an equal angle with it. Second, let $b' \neq 1$. If $k > n$, start with the $b'$-tape and introduce, by folding $(k-n)$ secondary crease lines, a crease line making an angle of $\frac{\pi}{2^{k-n} b'}$ with the forward direction of the top of the tape, with a second crease line under it making an equal angle with it. If $k \leq n$, then the $b'$-tape will already have on it a crease line making an angle of $\frac{\pi}{2^{k-n} b'}$ ($= \frac{2^{n-k} \pi}{b'}$) with the forward direction of the top of the tape, with a second crease line under it making an equal angle with it (see the proof of Proposition 5.1).

We have now described the \emph{initial configuration}, consisting of two crease lines, the first making an angle of $\alpha_{0}$ with the top of the tape, the second making an angle $\sigma_{0}$ with the first. We call $\alpha_{0}$ the $0$th \emph{proximand} and $\sigma_{0}$ the $0$th \emph{support}, and proceed to define inductively the $i$th \emph{proximand} $\alpha_{i}$ and the $i$th \emph{support} $\sigma_{i}$, $0 \leq i \leq n$. Of course, $\alpha_{0} = \sigma_{0} = \frac{\pi}{2^{k-n} b}$.

Suppose $\alpha_{i},\sigma_{i}$ already defined, $i < n$. Bisect $\sigma_{i}$ by a new crease line, into two equal angles $\lambda_{i},\rho_{i}$, with $\lambda_{i}$ to the left of $\rho_{i}$. If $\epsilon_{i+1} = 0$, define $\alpha_{i+1} = \alpha_{i}$, $\sigma_{i+1} = \rho_{i}$. If $\epsilon_{i+1} = 1$, define $\alpha_{i+1} = \alpha_{i} \cup \rho_{i}$, $\sigma_{i+1} = \lambda_{i}$. See Figure 5.2.

\begin{center}
%TeXCAD Picture [fig5.2.pic]. Options:
%\grade{\on}
%\emlines{\off}
%\epic{\off}
%\beziermacro{\on}
%\reduce{\on}
%\snapping{\off}
%\quality{8.00}
%\graddiff{0.01}
%\snapasp{1}
%\zoom{4.0000}
\unitlength 1mm % = 2.85pt
\linethickness{0.4pt}
\ifx\plotpoint\undefined\newsavebox{\plotpoint}\fi % GNUPLOT compatibility
\begin{picture}(117.5,39.75)(0,0)
\put(3,39.75){\line(1,0){35.5}}
\put(2.75,19.75){\line(1,0){34.75}}
%\dashline{1}(13,39.5)(15.5,19.75)
\multiput(12.93,39.43)(.0313,-.2469){4}{\line(0,-1){.2469}}
\multiput(13.18,37.45)(.0313,-.2469){4}{\line(0,-1){.2469}}
\multiput(13.43,35.48)(.0313,-.2469){4}{\line(0,-1){.2469}}
\multiput(13.68,33.5)(.0313,-.2469){4}{\line(0,-1){.2469}}
\multiput(13.93,31.53)(.0313,-.2469){4}{\line(0,-1){.2469}}
\multiput(14.18,29.55)(.0313,-.2469){4}{\line(0,-1){.2469}}
\multiput(14.43,27.58)(.0313,-.2469){4}{\line(0,-1){.2469}}
\multiput(14.68,25.6)(.0313,-.2469){4}{\line(0,-1){.2469}}
\multiput(14.93,23.63)(.0313,-.2469){4}{\line(0,-1){.2469}}
\multiput(15.18,21.65)(.0313,-.2469){4}{\line(0,-1){.2469}}
%\end
%\dashline{1}(13,39.25)(24,19.75)
\multiput(12.93,39.18)(.032738,-.058036){14}{\line(0,-1){.058036}}
\multiput(13.85,37.55)(.032738,-.058036){14}{\line(0,-1){.058036}}
\multiput(14.76,35.93)(.032738,-.058036){14}{\line(0,-1){.058036}}
\multiput(15.68,34.3)(.032738,-.058036){14}{\line(0,-1){.058036}}
\multiput(16.6,32.68)(.032738,-.058036){14}{\line(0,-1){.058036}}
\multiput(17.51,31.05)(.032738,-.058036){14}{\line(0,-1){.058036}}
\multiput(18.43,29.43)(.032738,-.058036){14}{\line(0,-1){.058036}}
\multiput(19.35,27.8)(.032738,-.058036){14}{\line(0,-1){.058036}}
\multiput(20.26,26.18)(.032738,-.058036){14}{\line(0,-1){.058036}}
\multiput(21.18,24.55)(.032738,-.058036){14}{\line(0,-1){.058036}}
\multiput(22.1,22.93)(.032738,-.058036){14}{\line(0,-1){.058036}}
\multiput(23.01,21.3)(.032738,-.058036){14}{\line(0,-1){.058036}}
%\end
%\dashline{1}(13.25,39.25)(33,20)
\multiput(13.18,39.18)(.034052,-.03319){20}{\line(1,0){.034052}}
\multiput(14.54,37.85)(.034052,-.03319){20}{\line(1,0){.034052}}
\multiput(15.9,36.52)(.034052,-.03319){20}{\line(1,0){.034052}}
\multiput(17.27,35.2)(.034052,-.03319){20}{\line(1,0){.034052}}
\multiput(18.63,33.87)(.034052,-.03319){20}{\line(1,0){.034052}}
\multiput(19.99,32.54)(.034052,-.03319){20}{\line(1,0){.034052}}
\multiput(21.35,31.21)(.034052,-.03319){20}{\line(1,0){.034052}}
\multiput(22.71,29.89)(.034052,-.03319){20}{\line(1,0){.034052}}
\multiput(24.08,28.56)(.034052,-.03319){20}{\line(1,0){.034052}}
\multiput(25.44,27.23)(.034052,-.03319){20}{\line(1,0){.034052}}
\multiput(26.8,25.9)(.034052,-.03319){20}{\line(1,0){.034052}}
\multiput(28.16,24.58)(.034052,-.03319){20}{\line(1,0){.034052}}
\multiput(29.52,23.25)(.034052,-.03319){20}{\line(1,0){.034052}}
\multiput(30.89,21.92)(.034052,-.03319){20}{\line(1,0){.034052}}
\multiput(32.25,20.59)(.034052,-.03319){20}{\line(1,0){.034052}}
%\end
\put(41.75,39.75){\line(1,0){34.5}}
\put(41.75,19.75){\line(1,0){34.75}}
\put(83.5,39.5){\line(1,0){34}}
\put(82.75,20){\line(1,0){34.5}}
%\dashline{1}(51,39.5)(55,19.75)
\multiput(50.93,39.43)(.03175,-.15675){6}{\line(0,-1){.15675}}
\multiput(51.31,37.55)(.03175,-.15675){6}{\line(0,-1){.15675}}
\multiput(51.69,35.67)(.03175,-.15675){6}{\line(0,-1){.15675}}
\multiput(52.07,33.79)(.03175,-.15675){6}{\line(0,-1){.15675}}
\multiput(52.45,31.91)(.03175,-.15675){6}{\line(0,-1){.15675}}
\multiput(52.83,30.02)(.03175,-.15675){6}{\line(0,-1){.15675}}
\multiput(53.22,28.14)(.03175,-.15675){6}{\line(0,-1){.15675}}
\multiput(53.6,26.26)(.03175,-.15675){6}{\line(0,-1){.15675}}
\multiput(53.98,24.38)(.03175,-.15675){6}{\line(0,-1){.15675}}
\multiput(54.36,22.5)(.03175,-.15675){6}{\line(0,-1){.15675}}
\multiput(54.74,20.62)(.03175,-.15675){6}{\line(0,-1){.15675}}
%\end
%\dashline{1}(51,39.25)(63.25,19.75)
\multiput(50.93,39.18)(.031901,-.050781){16}{\line(0,-1){.050781}}
\multiput(51.95,37.55)(.031901,-.050781){16}{\line(0,-1){.050781}}
\multiput(52.97,35.93)(.031901,-.050781){16}{\line(0,-1){.050781}}
\multiput(53.99,34.3)(.031901,-.050781){16}{\line(0,-1){.050781}}
\multiput(55.01,32.68)(.031901,-.050781){16}{\line(0,-1){.050781}}
\multiput(56.03,31.05)(.031901,-.050781){16}{\line(0,-1){.050781}}
\multiput(57.05,29.43)(.031901,-.050781){16}{\line(0,-1){.050781}}
\multiput(58.08,27.8)(.031901,-.050781){16}{\line(0,-1){.050781}}
\multiput(59.1,26.18)(.031901,-.050781){16}{\line(0,-1){.050781}}
\multiput(60.12,24.55)(.031901,-.050781){16}{\line(0,-1){.050781}}
\multiput(61.14,22.93)(.031901,-.050781){16}{\line(0,-1){.050781}}
\multiput(62.16,21.3)(.031901,-.050781){16}{\line(0,-1){.050781}}
%\end
%\dashline{1}(51.25,39.5)(72,19.75)
\multiput(51.18,39.43)(.034583,-.032917){20}{\line(1,0){.034583}}
\multiput(52.56,38.11)(.034583,-.032917){20}{\line(1,0){.034583}}
\multiput(53.95,36.8)(.034583,-.032917){20}{\line(1,0){.034583}}
\multiput(55.33,35.48)(.034583,-.032917){20}{\line(1,0){.034583}}
\multiput(56.71,34.16)(.034583,-.032917){20}{\line(1,0){.034583}}
\multiput(58.1,32.85)(.034583,-.032917){20}{\line(1,0){.034583}}
\multiput(59.48,31.53)(.034583,-.032917){20}{\line(1,0){.034583}}
\multiput(60.86,30.21)(.034583,-.032917){20}{\line(1,0){.034583}}
\multiput(62.25,28.9)(.034583,-.032917){20}{\line(1,0){.034583}}
\multiput(63.63,27.58)(.034583,-.032917){20}{\line(1,0){.034583}}
\multiput(65.01,26.26)(.034583,-.032917){20}{\line(1,0){.034583}}
\multiput(66.4,24.95)(.034583,-.032917){20}{\line(1,0){.034583}}
\multiput(67.78,23.63)(.034583,-.032917){20}{\line(1,0){.034583}}
\multiput(69.16,22.31)(.034583,-.032917){20}{\line(1,0){.034583}}
\multiput(70.55,21)(.034583,-.032917){20}{\line(1,0){.034583}}
%\end
%\dashline{1}(92.75,39.25)(97,20.25)
\multiput(92.68,39.18)(.03373,-.15079){6}{\line(0,-1){.15079}}
\multiput(93.08,37.37)(.03373,-.15079){6}{\line(0,-1){.15079}}
\multiput(93.49,35.56)(.03373,-.15079){6}{\line(0,-1){.15079}}
\multiput(93.89,33.75)(.03373,-.15079){6}{\line(0,-1){.15079}}
\multiput(94.3,31.94)(.03373,-.15079){6}{\line(0,-1){.15079}}
\multiput(94.7,30.13)(.03373,-.15079){6}{\line(0,-1){.15079}}
\multiput(95.11,28.32)(.03373,-.15079){6}{\line(0,-1){.15079}}
\multiput(95.51,26.51)(.03373,-.15079){6}{\line(0,-1){.15079}}
\multiput(95.92,24.7)(.03373,-.15079){6}{\line(0,-1){.15079}}
\multiput(96.32,22.89)(.03373,-.15079){6}{\line(0,-1){.15079}}
\multiput(96.73,21.08)(.03373,-.15079){6}{\line(0,-1){.15079}}
%\end
%\dashline{1}(93,39)(105,20)
\multiput(92.93,38.93)(.033333,-.052778){15}{\line(0,-1){.052778}}
\multiput(93.93,37.35)(.033333,-.052778){15}{\line(0,-1){.052778}}
\multiput(94.93,35.76)(.033333,-.052778){15}{\line(0,-1){.052778}}
\multiput(95.93,34.18)(.033333,-.052778){15}{\line(0,-1){.052778}}
\multiput(96.93,32.6)(.033333,-.052778){15}{\line(0,-1){.052778}}
\multiput(97.93,31.01)(.033333,-.052778){15}{\line(0,-1){.052778}}
\multiput(98.93,29.43)(.033333,-.052778){15}{\line(0,-1){.052778}}
\multiput(99.93,27.85)(.033333,-.052778){15}{\line(0,-1){.052778}}
\multiput(100.93,26.26)(.033333,-.052778){15}{\line(0,-1){.052778}}
\multiput(101.93,24.68)(.033333,-.052778){15}{\line(0,-1){.052778}}
\multiput(102.93,23.1)(.033333,-.052778){15}{\line(0,-1){.052778}}
\multiput(103.93,21.51)(.033333,-.052778){15}{\line(0,-1){.052778}}
%\end
%\dashline{1}(92.75,39.5)(113.5,20)
\multiput(92.68,39.43)(.034583,-.0325){20}{\line(1,0){.034583}}
\multiput(94.06,38.13)(.034583,-.0325){20}{\line(1,0){.034583}}
\multiput(95.45,36.83)(.034583,-.0325){20}{\line(1,0){.034583}}
\multiput(96.83,35.53)(.034583,-.0325){20}{\line(1,0){.034583}}
\multiput(98.21,34.23)(.034583,-.0325){20}{\line(1,0){.034583}}
\multiput(99.6,32.93)(.034583,-.0325){20}{\line(1,0){.034583}}
\multiput(100.98,31.63)(.034583,-.0325){20}{\line(1,0){.034583}}
\multiput(102.36,30.33)(.034583,-.0325){20}{\line(1,0){.034583}}
\multiput(103.75,29.03)(.034583,-.0325){20}{\line(1,0){.034583}}
\multiput(105.13,27.73)(.034583,-.0325){20}{\line(1,0){.034583}}
\multiput(106.51,26.43)(.034583,-.0325){20}{\line(1,0){.034583}}
\multiput(107.9,25.13)(.034583,-.0325){20}{\line(1,0){.034583}}
\multiput(109.28,23.83)(.034583,-.0325){20}{\line(1,0){.034583}}
\multiput(110.66,22.53)(.034583,-.0325){20}{\line(1,0){.034583}}
\multiput(112.05,21.23)(.034583,-.0325){20}{\line(1,0){.034583}}
%\end
\put(17,28.5){\vector(-1,0){2.25}}
%\vector(20.5,29)(22.25,30)
\put(22.25,30){\vector(2,1){.07}}\multiput(20.5,29)(.058333,.033333){30}{\line(1,0){.058333}}
%\end
%\qbezvec[both](97.46,39.43)(99.99,34.16)(95.38,35.12)
\put(95.38,35.12){\vector(-1,0){.07}}\put(97.46,39.43){\vector(-1,2){.07}}\qbezier(97.46,39.43)(99.99,34.16)(95.38,35.12)
%\end
\put(18.5,28.75){\makebox(0,0)[cc]{$\sigma_{i}$}}
\put(18,23.75){\makebox(0,0)[cc]{$\lambda_{i}$}}
\put(26.5,22.75){\makebox(0,0)[cc]{$\rho_{i}$}}
\put(19.75,37.25){\makebox(0,0)[cc]{$\alpha_{i}$}}
\put(61.25,36){\makebox(0,0)[cc]{$\alpha_{i+1}$}}
\put(61.85,25.77){\makebox(0,0)[cc]{$\sigma_{i+1}$}}
\put(97.45,25.76){\makebox(0,0)[cc]{$\sigma_{i+1}$}}
\put(102,35){\makebox(0,0)[cc]{$\alpha_{i+1}$}}
\put(20,13.25){\makebox(0,0)[cc]{$i$th stage}}
\put(60,13.25){\makebox(0,0)[cc]{$(i+1)$st stage, $\epsilon_{i+1} = 0$}}
\put(99.75,13.25){\makebox(0,0)[cc]{$(i+1)$st stage, $\epsilon_{i+1} = 1$}}
\put(60,2.75){\makebox(0,0)[cc]{Figure 5.2}}
\end{picture}
\end{center}

\begin{description}

  \item[Proposition 5.2:] The $n$th proximand $\alpha_{n}$ yields the required angle of $\left( \frac{a}{b} \right) \pi$.

  \item[Proof:] We argue by induction on $i$ that
    \begin{equation}
    \alpha_{i} = \frac{\pi a_{i}}{2^{k}b'}, \qquad \sigma_{i} = \frac{\pi}{2^{k-n+i}b'}  \tag{5.2}
    \end{equation}
    where
    \begin{equation}
    a_{i} = \epsilon_{0} \epsilon_{1} \cdots \epsilon_{i}0 \cdots 0 \text{ (length $(n+1)$)}   \tag{5.3}
    \end{equation}

    Then (5.2) clearly holds for $i = 0$, since $a_{0} = 2^{n}$. Given (5.2) and $\epsilon_{i+1} = 0$, then $\alpha_{i+1} = \alpha_{i} = \frac{\pi a_{i+1}}{2^{k} b'}$, since $a_{i+1} = a_{i}$, and $\sigma_{i+1} = \left( \frac{1}{2} \right) \sigma_{i} = \frac{\pi}{2^{k-n+i+1} b'}$. Given (5.2) and $\epsilon_{i+1} = 1$, then
    \begin{align*}
    \alpha_{i+1} = \alpha_{i} + \frac{1}{2} \sigma_{i} & = \frac{\pi a_{i}}{2^{k}b'} + \frac{\pi}{2^{k-n+i+1}b'} = \frac{\pi}{2^{k}b'} (a_{i} + 2^{n-i-1}) \\
    & = \frac{\pi}{2^{k}b'} (\epsilon_{0} \epsilon_{1} \cdots \epsilon_{i}0 \cdots 0 + 00 \cdots 010 \cdots 0)
    \end{align*}
    where the digit $1$ is in place $(n-i)$ counting from the right, and thus in place $(i+1)$ counting from the left. Thus
    \begin{equation*}
    \alpha_{i+1} = \frac{\pi}{2^{k}b'} (\epsilon_{0} \epsilon_{1} \cdots \epsilon_{i} 10 \cdots 0) = \frac{\pi a_{i+1}}{2^{k}b'}; \qquad \sigma_{i+1} = \frac{1}{2} \sigma_{i} = \frac{\pi}{2^{k-n+i+1}b'}
    \end{equation*}

\end{description}

This establishes (5.2), and we see immediately that $a_{n} = a$, so that $\alpha_{n} = \frac{\pi a}{2^{k} b'} = \frac{\pi a}{b}$. Given the crease line on the tape making an angle of $\frac{\pi a}{b}$, we may plainly complete the FAT-algorithm to construct the regular star $\left\{ \frac{b}{a} \right\}$-gon. Figure 5.3 illustrates this process for the $\left\{ \frac{10}{3} \right\}$-gon. Figure 5.4 illustrates a $\left\{ \frac{17}{6} \right\}$-gon constructed from $\{ 1,1,2 \}$ tape in accordance with Proposition 5.1.

\begin{center}
%TeXCAD Picture [fig5.3.pic]. Options:
%\grade{\on}
%\emlines{\off}
%\epic{\off}
%\beziermacro{\on}
%\reduce{\on}
%\snapping{\off}
%\quality{8.00}
%\graddiff{0.01}
%\snapasp{1}
%\zoom{4.0000}
\unitlength 1mm % = 2.85pt
\linethickness{0.4pt}
\ifx\plotpoint\undefined\newsavebox{\plotpoint}\fi % GNUPLOT compatibility
\begin{picture}(100.5,134.34)(0,0)
\put(1.5,131.59){\line(1,0){99}}
\put(6.25,122.59){\line(1,0){94}}
%\dashline{1}(9.5,131.34)(7.75,122.84)
\multiput(9.43,131.27)(-.03241,-.15741){6}{\line(0,-1){.15741}}
\multiput(9.04,129.39)(-.03241,-.15741){6}{\line(0,-1){.15741}}
\multiput(8.65,127.5)(-.03241,-.15741){6}{\line(0,-1){.15741}}
\multiput(8.26,125.61)(-.03241,-.15741){6}{\line(0,-1){.15741}}
\multiput(7.87,123.72)(-.03241,-.15741){6}{\line(0,-1){.15741}}
%\end
%\dashline{1}(7.75,122.84)(22,131.09)
\multiput(7.68,122.77)(.056548,.032738){14}{\line(1,0){.056548}}
\multiput(9.26,123.69)(.056548,.032738){14}{\line(1,0){.056548}}
\multiput(10.85,124.61)(.056548,.032738){14}{\line(1,0){.056548}}
\multiput(12.43,125.52)(.056548,.032738){14}{\line(1,0){.056548}}
\multiput(14.01,126.44)(.056548,.032738){14}{\line(1,0){.056548}}
\multiput(15.6,127.36)(.056548,.032738){14}{\line(1,0){.056548}}
\multiput(17.18,128.27)(.056548,.032738){14}{\line(1,0){.056548}}
\multiput(18.76,129.19)(.056548,.032738){14}{\line(1,0){.056548}}
\multiput(20.35,130.11)(.056548,.032738){14}{\line(1,0){.056548}}
%\end
%\dashline{1}(22,131.09)(36.75,122.59)
\multiput(21.93,131.02)(.058532,-.03373){14}{\line(1,0){.058532}}
\multiput(23.57,130.08)(.058532,-.03373){14}{\line(1,0){.058532}}
\multiput(25.21,129.14)(.058532,-.03373){14}{\line(1,0){.058532}}
\multiput(26.85,128.19)(.058532,-.03373){14}{\line(1,0){.058532}}
\multiput(28.49,127.25)(.058532,-.03373){14}{\line(1,0){.058532}}
\multiput(30.12,126.3)(.058532,-.03373){14}{\line(1,0){.058532}}
\multiput(31.76,125.36)(.058532,-.03373){14}{\line(1,0){.058532}}
\multiput(33.4,124.41)(.058532,-.03373){14}{\line(1,0){.058532}}
\multiput(35.04,123.47)(.058532,-.03373){14}{\line(1,0){.058532}}
%\end
%\dashline{1}(36.75,122.59)(49.75,131.34)
\multiput(36.68,122.52)(.047794,.032169){16}{\line(1,0){.047794}}
\multiput(38.21,123.55)(.047794,.032169){16}{\line(1,0){.047794}}
\multiput(39.74,124.58)(.047794,.032169){16}{\line(1,0){.047794}}
\multiput(41.27,125.61)(.047794,.032169){16}{\line(1,0){.047794}}
\multiput(42.8,126.64)(.047794,.032169){16}{\line(1,0){.047794}}
\multiput(44.33,127.67)(.047794,.032169){16}{\line(1,0){.047794}}
\multiput(45.86,128.7)(.047794,.032169){16}{\line(1,0){.047794}}
\multiput(47.39,129.73)(.047794,.032169){16}{\line(1,0){.047794}}
\multiput(48.92,130.76)(.047794,.032169){16}{\line(1,0){.047794}}
%\end
%\dashline{1}(49.75,131.34)(64.5,122.84)
\multiput(49.68,131.27)(.058532,-.03373){14}{\line(1,0){.058532}}
\multiput(51.32,130.33)(.058532,-.03373){14}{\line(1,0){.058532}}
\multiput(52.96,129.39)(.058532,-.03373){14}{\line(1,0){.058532}}
\multiput(54.6,128.44)(.058532,-.03373){14}{\line(1,0){.058532}}
\multiput(56.24,127.5)(.058532,-.03373){14}{\line(1,0){.058532}}
\multiput(57.87,126.55)(.058532,-.03373){14}{\line(1,0){.058532}}
\multiput(59.51,125.61)(.058532,-.03373){14}{\line(1,0){.058532}}
\multiput(61.15,124.66)(.058532,-.03373){14}{\line(1,0){.058532}}
\multiput(62.79,123.72)(.058532,-.03373){14}{\line(1,0){.058532}}
%\end
%\dashline{1}(64.5,122.84)(79.25,131.59)
\multiput(64.43,122.77)(.05463,.032407){15}{\line(1,0){.05463}}
\multiput(66.07,123.75)(.05463,.032407){15}{\line(1,0){.05463}}
\multiput(67.71,124.72)(.05463,.032407){15}{\line(1,0){.05463}}
\multiput(69.35,125.69)(.05463,.032407){15}{\line(1,0){.05463}}
\multiput(70.99,126.66)(.05463,.032407){15}{\line(1,0){.05463}}
\multiput(72.62,127.64)(.05463,.032407){15}{\line(1,0){.05463}}
\multiput(74.26,128.61)(.05463,.032407){15}{\line(1,0){.05463}}
\multiput(75.9,129.58)(.05463,.032407){15}{\line(1,0){.05463}}
\multiput(77.54,130.55)(.05463,.032407){15}{\line(1,0){.05463}}
%\end
%\dashline{1}(79.25,131.59)(92.25,122.59)
\multiput(79.18,131.52)(.048148,-.033333){15}{\line(1,0){.048148}}
\multiput(80.62,130.52)(.048148,-.033333){15}{\line(1,0){.048148}}
\multiput(82.07,129.52)(.048148,-.033333){15}{\line(1,0){.048148}}
\multiput(83.51,128.52)(.048148,-.033333){15}{\line(1,0){.048148}}
\multiput(84.96,127.52)(.048148,-.033333){15}{\line(1,0){.048148}}
\multiput(86.4,126.52)(.048148,-.033333){15}{\line(1,0){.048148}}
\multiput(87.85,125.52)(.048148,-.033333){15}{\line(1,0){.048148}}
\multiput(89.29,124.52)(.048148,-.033333){15}{\line(1,0){.048148}}
\multiput(90.74,123.52)(.048148,-.033333){15}{\line(1,0){.048148}}
%\end
%\dashline{1}(22.25,130.59)(24.75,122.84)
\multiput(22.18,130.52)(.03086,-.09568){9}{\line(0,-1){.09568}}
\multiput(22.74,128.8)(.03086,-.09568){9}{\line(0,-1){.09568}}
\multiput(23.29,127.08)(.03086,-.09568){9}{\line(0,-1){.09568}}
\multiput(23.85,125.36)(.03086,-.09568){9}{\line(0,-1){.09568}}
\multiput(24.4,123.64)(.03086,-.09568){9}{\line(0,-1){.09568}}
%\end
%\dashline{1}(23,130.34)(29.5,122.59)
\multiput(22.93,130.27)(.032828,-.039141){18}{\line(0,-1){.039141}}
\multiput(24.11,128.87)(.032828,-.039141){18}{\line(0,-1){.039141}}
\multiput(25.29,127.46)(.032828,-.039141){18}{\line(0,-1){.039141}}
\multiput(26.48,126.05)(.032828,-.039141){18}{\line(0,-1){.039141}}
\multiput(27.66,124.64)(.032828,-.039141){18}{\line(0,-1){.039141}}
\multiput(28.84,123.23)(.032828,-.039141){18}{\line(0,-1){.039141}}
%\end
%\dashline{1}(27.5,124.84)(34,117.09)
\multiput(27.43,124.77)(.032828,-.039141){18}{\line(0,-1){.039141}}
\multiput(28.61,123.37)(.032828,-.039141){18}{\line(0,-1){.039141}}
\multiput(29.79,121.96)(.032828,-.039141){18}{\line(0,-1){.039141}}
\multiput(30.98,120.55)(.032828,-.039141){18}{\line(0,-1){.039141}}
\multiput(32.16,119.14)(.032828,-.039141){18}{\line(0,-1){.039141}}
\multiput(33.34,117.73)(.032828,-.039141){18}{\line(0,-1){.039141}}
%\end
%\dashline{1}(36.5,122.84)(38.75,131.34)
\multiput(36.43,122.77)(.03214,.12143){7}{\line(0,1){.12143}}
\multiput(36.88,124.47)(.03214,.12143){7}{\line(0,1){.12143}}
\multiput(37.33,126.17)(.03214,.12143){7}{\line(0,1){.12143}}
\multiput(37.78,127.87)(.03214,.12143){7}{\line(0,1){.12143}}
\multiput(38.23,129.57)(.03214,.12143){7}{\line(0,1){.12143}}
%\end
%\dashline{1}(49.75,130.84)(53.5,122.59)
\multiput(49.68,130.77)(.03125,-.06875){12}{\line(0,-1){.06875}}
\multiput(50.43,129.12)(.03125,-.06875){12}{\line(0,-1){.06875}}
\multiput(51.18,127.47)(.03125,-.06875){12}{\line(0,-1){.06875}}
\multiput(51.93,125.82)(.03125,-.06875){12}{\line(0,-1){.06875}}
\multiput(52.68,124.17)(.03125,-.06875){12}{\line(0,-1){.06875}}
%\end
%\dashline{1}(64.75,122.59)(67.25,131.84)
\multiput(64.68,122.52)(.03247,.12013){7}{\line(0,1){.12013}}
\multiput(65.13,124.21)(.03247,.12013){7}{\line(0,1){.12013}}
\multiput(65.59,125.89)(.03247,.12013){7}{\line(0,1){.12013}}
\multiput(66.04,127.57)(.03247,.12013){7}{\line(0,1){.12013}}
\multiput(66.5,129.25)(.03247,.12013){7}{\line(0,1){.12013}}
\multiput(66.95,130.93)(.03247,.12013){7}{\line(0,1){.12013}}
%\end
%\dashline{1}(79.25,131.09)(83,123.09)
\multiput(79.18,131.02)(.03125,-.06667){12}{\line(0,-1){.06667}}
\multiput(79.93,129.42)(.03125,-.06667){12}{\line(0,-1){.06667}}
\multiput(80.68,127.82)(.03125,-.06667){12}{\line(0,-1){.06667}}
\multiput(81.43,126.22)(.03125,-.06667){12}{\line(0,-1){.06667}}
\multiput(82.18,124.62)(.03125,-.06667){12}{\line(0,-1){.06667}}
%\end
%\dashline{1}(83,123.59)(83.25,122.84)
\multiput(82.93,123.52)(.0313,-.0938){4}{\line(0,-1){.0938}}
%\end
\put(2.75,104.34){\line(1,0){22.25}}
\put(25,104.34){\line(1,-1){8.25}}
%\emline(33.25,96.09)(28.75,84.09)
\multiput(33.25,96.09)(-.03358209,-.08955224){134}{\line(0,-1){.08955224}}
%\end
%\emline(15.75,79.84)(25,104.09)
\multiput(15.75,79.84)(.033636364,.088181818){275}{\line(0,1){.088181818}}
%\end
\put(21.75,95.59){\line(-1,0){16.5}}
%\dashline{1}(5.25,98.59)(10.5,95.84)
\multiput(5.18,98.52)(.0625,-.03274){12}{\line(1,0){.0625}}
\multiput(6.68,97.74)(.0625,-.03274){12}{\line(1,0){.0625}}
\multiput(8.18,96.95)(.0625,-.03274){12}{\line(1,0){.0625}}
\multiput(9.68,96.17)(.0625,-.03274){12}{\line(1,0){.0625}}
%\end
%\dashline{1}(10.5,95.84)(13.75,104.59)
\multiput(10.43,95.77)(.0325,.0875){10}{\line(0,1){.0875}}
\multiput(11.08,97.52)(.0325,.0875){10}{\line(0,1){.0875}}
\multiput(11.73,99.27)(.0325,.0875){10}{\line(0,1){.0875}}
\multiput(12.38,101.02)(.0325,.0875){10}{\line(0,1){.0875}}
\multiput(13.03,102.77)(.0325,.0875){10}{\line(0,1){.0875}}
%\end
%\dashline{1}(10.5,95.59)(25,104.34)
\multiput(10.43,95.52)(.053704,.032407){15}{\line(1,0){.053704}}
\multiput(12.04,96.5)(.053704,.032407){15}{\line(1,0){.053704}}
\multiput(13.65,97.47)(.053704,.032407){15}{\line(1,0){.053704}}
\multiput(15.26,98.44)(.053704,.032407){15}{\line(1,0){.053704}}
\multiput(16.87,99.41)(.053704,.032407){15}{\line(1,0){.053704}}
\multiput(18.49,100.39)(.053704,.032407){15}{\line(1,0){.053704}}
\multiput(20.1,101.36)(.053704,.032407){15}{\line(1,0){.053704}}
\multiput(21.71,102.33)(.053704,.032407){15}{\line(1,0){.053704}}
\multiput(23.32,103.3)(.053704,.032407){15}{\line(1,0){.053704}}
%\end
%\dashline{1}(25.25,104.09)(30.75,90.09)
\multiput(25.18,104.02)(.03125,-.07955){11}{\line(0,-1){.07955}}
\multiput(25.87,102.27)(.03125,-.07955){11}{\line(0,-1){.07955}}
\multiput(26.55,100.52)(.03125,-.07955){11}{\line(0,-1){.07955}}
\multiput(27.24,98.77)(.03125,-.07955){11}{\line(0,-1){.07955}}
\multiput(27.93,97.02)(.03125,-.07955){11}{\line(0,-1){.07955}}
\multiput(28.62,95.27)(.03125,-.07955){11}{\line(0,-1){.07955}}
\multiput(29.3,93.52)(.03125,-.07955){11}{\line(0,-1){.07955}}
\multiput(29.99,91.77)(.03125,-.07955){11}{\line(0,-1){.07955}}
%\end
%\dashline{1}(30.75,90.09)(16.75,82.34)
\multiput(30.68,90.02)(-.058824,-.032563){14}{\line(-1,0){.058824}}
\multiput(29.03,89.11)(-.058824,-.032563){14}{\line(-1,0){.058824}}
\multiput(27.39,88.2)(-.058824,-.032563){14}{\line(-1,0){.058824}}
\multiput(25.74,87.29)(-.058824,-.032563){14}{\line(-1,0){.058824}}
\multiput(24.09,86.38)(-.058824,-.032563){14}{\line(-1,0){.058824}}
\multiput(22.44,85.47)(-.058824,-.032563){14}{\line(-1,0){.058824}}
\multiput(20.8,84.55)(-.058824,-.032563){14}{\line(-1,0){.058824}}
\multiput(19.15,83.64)(-.058824,-.032563){14}{\line(-1,0){.058824}}
\multiput(17.5,82.73)(-.058824,-.032563){14}{\line(-1,0){.058824}}
%\end
%\dashline{1}(30.75,90.09)(20,91.09)
\put(30.68,90.02){\line(-1,0){.977}}
\put(28.73,90.21){\line(-1,0){.977}}
\put(26.77,90.39){\line(-1,0){.977}}
\put(24.82,90.57){\line(-1,0){.977}}
\put(22.86,90.75){\line(-1,0){.977}}
\put(20.91,90.93){\line(-1,0){.977}}
%\end
%\dashline{1}(24.75,103.59)(26.25,107.59)
\multiput(24.68,103.52)(.03125,.08333){8}{\line(0,1){.08333}}
\multiput(25.18,104.86)(.03125,.08333){8}{\line(0,1){.08333}}
\multiput(25.68,106.19)(.03125,.08333){8}{\line(0,1){.08333}}
%\end
%\dashline{1}(25.75,106.34)(27.25,110.34)
\multiput(25.68,106.27)(.03125,.08333){8}{\line(0,1){.08333}}
\multiput(26.18,107.61)(.03125,.08333){8}{\line(0,1){.08333}}
\multiput(26.68,108.94)(.03125,.08333){8}{\line(0,1){.08333}}
%\end
\put(69,104.34){\line(1,0){23.5}}
%\emline(92.5,104.34)(82.75,79.34)
\multiput(92.5,104.34)(-.033737024,-.08650519){289}{\line(0,-1){.08650519}}
%\end
%\emline(71.25,82.34)(77.5,95.84)
\multiput(71.25,82.34)(.03360215,.07258065){186}{\line(0,1){.07258065}}
%\end
\put(77.5,95.84){\line(-1,0){9}}
%\emline(77.25,95.34)(79.5,100.59)
\multiput(77.25,95.34)(.0335821,.0783582){67}{\line(0,1){.0783582}}
%\end
%\emline(79.5,100.59)(92.25,104.09)
\multiput(79.5,100.59)(.12259615,.03365385){104}{\line(1,0){.12259615}}
%\end
%\emline(79.75,100.59)(88.75,95.59)
\multiput(79.75,100.59)(.06040268,-.03355705){149}{\line(1,0){.06040268}}
%\end
%\dashline{1}(79.5,100.84)(81,105.09)
\multiput(79.43,100.77)(.03125,.08854){8}{\line(0,1){.08854}}
\multiput(79.93,102.19)(.03125,.08854){8}{\line(0,1){.08854}}
\multiput(80.43,103.61)(.03125,.08854){8}{\line(0,1){.08854}}
%\end
%\dashline{1}(77.25,95.59)(92.25,103.84)
\multiput(77.18,95.52)(.060729,.033401){13}{\line(1,0){.060729}}
\multiput(78.76,96.39)(.060729,.033401){13}{\line(1,0){.060729}}
\multiput(80.34,97.26)(.060729,.033401){13}{\line(1,0){.060729}}
\multiput(81.92,98.13)(.060729,.033401){13}{\line(1,0){.060729}}
\multiput(83.5,99)(.060729,.033401){13}{\line(1,0){.060729}}
\multiput(85.07,99.87)(.060729,.033401){13}{\line(1,0){.060729}}
\multiput(86.65,100.74)(.060729,.033401){13}{\line(1,0){.060729}}
\multiput(88.23,101.6)(.060729,.033401){13}{\line(1,0){.060729}}
\multiput(89.81,102.47)(.060729,.033401){13}{\line(1,0){.060729}}
\multiput(91.39,103.34)(.060729,.033401){13}{\line(1,0){.060729}}
%\end
%\dashline{1}(73,99.09)(86.75,89.59)
\multiput(72.93,99.02)(.047743,-.032986){16}{\line(1,0){.047743}}
\multiput(74.46,97.97)(.047743,-.032986){16}{\line(1,0){.047743}}
\multiput(75.99,96.91)(.047743,-.032986){16}{\line(1,0){.047743}}
\multiput(77.51,95.86)(.047743,-.032986){16}{\line(1,0){.047743}}
\multiput(79.04,94.8)(.047743,-.032986){16}{\line(1,0){.047743}}
\multiput(80.57,93.75)(.047743,-.032986){16}{\line(1,0){.047743}}
\multiput(82.1,92.69)(.047743,-.032986){16}{\line(1,0){.047743}}
\multiput(83.62,91.64)(.047743,-.032986){16}{\line(1,0){.047743}}
\multiput(85.15,90.58)(.047743,-.032986){16}{\line(1,0){.047743}}
%\end
%\dashline{1}(77.5,95.59)(84,82.34)
\multiput(77.43,95.52)(.03125,-.063702){13}{\line(0,-1){.063702}}
\multiput(78.24,93.87)(.03125,-.063702){13}{\line(0,-1){.063702}}
\multiput(79.05,92.21)(.03125,-.063702){13}{\line(0,-1){.063702}}
\multiput(79.87,90.56)(.03125,-.063702){13}{\line(0,-1){.063702}}
\multiput(80.68,88.9)(.03125,-.063702){13}{\line(0,-1){.063702}}
\multiput(81.49,87.24)(.03125,-.063702){13}{\line(0,-1){.063702}}
\multiput(82.3,85.59)(.03125,-.063702){13}{\line(0,-1){.063702}}
\multiput(83.12,83.93)(.03125,-.063702){13}{\line(0,-1){.063702}}
%\end
%\dashline{1}(84,82.34)(71.5,82.34)
\put(83.93,82.27){\line(-1,0){.962}}
\put(82.01,82.27){\line(-1,0){.962}}
\put(80.08,82.27){\line(-1,0){.962}}
\put(78.16,82.27){\line(-1,0){.962}}
\put(76.24,82.27){\line(-1,0){.962}}
\put(74.31,82.27){\line(-1,0){.962}}
\put(72.39,82.27){\line(-1,0){.962}}
%\end
\thicklines
\put(28.68,64.09){\line(0,-1){10.5}}
\put(56.18,20.59){\line(5,-1){13.75}}
\put(69.93,17.84){\line(0,1){10.25}}
%\emline(55.93,20.84)(51.68,21.84)
\multiput(55.93,20.84)(-.141667,.033333){30}{\line(-1,0){.141667}}
%\end
\thinlines
%\emline(69.93,17.84)(65.68,26.59)
\multiput(69.93,17.84)(-.03373016,.06944444){126}{\line(0,1){.06944444}}
%\end
%\emline(48.63,11.99)(51.63,22.24)
\multiput(48.63,11.99)(.0337079,.1151685){89}{\line(0,1){.1151685}}
%\end
%\emline(32.28,55.09)(54.78,47.84)
\multiput(32.28,55.09)(.10465116,-.03372093){215}{\line(1,0){.10465116}}
%\end
%\emline(71.43,37.59)(48.68,30.84)
\multiput(71.43,37.59)(-.11318408,-.03358209){201}{\line(-1,0){.11318408}}
%\end
%\dashline{1}(55.43,60.34)(32.43,55.09)
\multiput(55.36,60.27)(-.13143,-.03){7}{\line(-1,0){.13143}}
\multiput(53.52,59.85)(-.13143,-.03){7}{\line(-1,0){.13143}}
\multiput(51.68,59.43)(-.13143,-.03){7}{\line(-1,0){.13143}}
\multiput(49.84,59.01)(-.13143,-.03){7}{\line(-1,0){.13143}}
\multiput(48,58.59)(-.13143,-.03){7}{\line(-1,0){.13143}}
\multiput(46.16,58.17)(-.13143,-.03){7}{\line(-1,0){.13143}}
\multiput(44.32,57.75)(-.13143,-.03){7}{\line(-1,0){.13143}}
\multiput(42.48,57.33)(-.13143,-.03){7}{\line(-1,0){.13143}}
\multiput(40.64,56.91)(-.13143,-.03){7}{\line(-1,0){.13143}}
\multiput(38.8,56.49)(-.13143,-.03){7}{\line(-1,0){.13143}}
\multiput(36.96,56.07)(-.13143,-.03){7}{\line(-1,0){.13143}}
\multiput(35.12,55.65)(-.13143,-.03){7}{\line(-1,0){.13143}}
\multiput(33.28,55.23)(-.13143,-.03){7}{\line(-1,0){.13143}}
%\end
%\dashline{1}(40.43,60.34)(25.68,43.84)
\multiput(40.36,60.27)(-.032065,-.03587){20}{\line(0,-1){.03587}}
\multiput(39.08,58.84)(-.032065,-.03587){20}{\line(0,-1){.03587}}
\multiput(37.8,57.4)(-.032065,-.03587){20}{\line(0,-1){.03587}}
\multiput(36.51,55.97)(-.032065,-.03587){20}{\line(0,-1){.03587}}
\multiput(35.23,54.54)(-.032065,-.03587){20}{\line(0,-1){.03587}}
\multiput(33.95,53.1)(-.032065,-.03587){20}{\line(0,-1){.03587}}
\multiput(32.67,51.67)(-.032065,-.03587){20}{\line(0,-1){.03587}}
\multiput(31.38,50.23)(-.032065,-.03587){20}{\line(0,-1){.03587}}
\multiput(30.1,48.8)(-.032065,-.03587){20}{\line(0,-1){.03587}}
\multiput(28.82,47.36)(-.032065,-.03587){20}{\line(0,-1){.03587}}
\multiput(27.53,45.93)(-.032065,-.03587){20}{\line(0,-1){.03587}}
\multiput(26.25,44.49)(-.032065,-.03587){20}{\line(0,-1){.03587}}
%\end
%\dashline{1}(23.43,40.84)(38.18,24.34)
\multiput(23.36,40.77)(.032065,-.03587){20}{\line(0,-1){.03587}}
\multiput(24.64,39.34)(.032065,-.03587){20}{\line(0,-1){.03587}}
\multiput(25.93,37.9)(.032065,-.03587){20}{\line(0,-1){.03587}}
\multiput(27.21,36.47)(.032065,-.03587){20}{\line(0,-1){.03587}}
\multiput(28.49,35.04)(.032065,-.03587){20}{\line(0,-1){.03587}}
\multiput(29.77,33.6)(.032065,-.03587){20}{\line(0,-1){.03587}}
\multiput(31.06,32.17)(.032065,-.03587){20}{\line(0,-1){.03587}}
\multiput(32.34,30.73)(.032065,-.03587){20}{\line(0,-1){.03587}}
\multiput(33.62,29.3)(.032065,-.03587){20}{\line(0,-1){.03587}}
\multiput(34.9,27.86)(.032065,-.03587){20}{\line(0,-1){.03587}}
\multiput(36.19,26.43)(.032065,-.03587){20}{\line(0,-1){.03587}}
\multiput(37.47,24.99)(.032065,-.03587){20}{\line(0,-1){.03587}}
%\end
%\dashline{1}(29.18,28.34)(52.43,21.59)
\multiput(29.11,28.27)(.10333,-.03){9}{\line(1,0){.10333}}
\multiput(30.97,27.73)(.10333,-.03){9}{\line(1,0){.10333}}
\multiput(32.83,27.19)(.10333,-.03){9}{\line(1,0){.10333}}
\multiput(34.69,26.65)(.10333,-.03){9}{\line(1,0){.10333}}
\multiput(36.55,26.11)(.10333,-.03){9}{\line(1,0){.10333}}
\multiput(38.41,25.57)(.10333,-.03){9}{\line(1,0){.10333}}
\multiput(40.27,25.03)(.10333,-.03){9}{\line(1,0){.10333}}
\multiput(42.13,24.49)(.10333,-.03){9}{\line(1,0){.10333}}
\multiput(43.99,23.95)(.10333,-.03){9}{\line(1,0){.10333}}
\multiput(45.85,23.41)(.10333,-.03){9}{\line(1,0){.10333}}
\multiput(47.71,22.87)(.10333,-.03){9}{\line(1,0){.10333}}
\multiput(49.57,22.33)(.10333,-.03){9}{\line(1,0){.10333}}
\multiput(51.43,21.79)(.10333,-.03){9}{\line(1,0){.10333}}
%\end
%\dashline{1}(56.43,20.59)(71.43,37.59)
\multiput(56.36,20.52)(.032895,.037281){19}{\line(0,1){.037281}}
\multiput(57.61,21.94)(.032895,.037281){19}{\line(0,1){.037281}}
\multiput(58.86,23.36)(.032895,.037281){19}{\line(0,1){.037281}}
\multiput(60.11,24.77)(.032895,.037281){19}{\line(0,1){.037281}}
\multiput(61.36,26.19)(.032895,.037281){19}{\line(0,1){.037281}}
\multiput(62.61,27.61)(.032895,.037281){19}{\line(0,1){.037281}}
\multiput(63.86,29.02)(.032895,.037281){19}{\line(0,1){.037281}}
\multiput(65.11,30.44)(.032895,.037281){19}{\line(0,1){.037281}}
\multiput(66.36,31.86)(.032895,.037281){19}{\line(0,1){.037281}}
\multiput(67.61,33.27)(.032895,.037281){19}{\line(0,1){.037281}}
\multiput(68.86,34.69)(.032895,.037281){19}{\line(0,1){.037281}}
\multiput(70.11,36.11)(.032895,.037281){19}{\line(0,1){.037281}}
%\end
%\dashline{1}(68.93,52.84)(44.43,59.59)
\multiput(68.86,52.77)(-.11779,.03245){8}{\line(-1,0){.11779}}
\multiput(66.98,53.29)(-.11779,.03245){8}{\line(-1,0){.11779}}
\multiput(65.09,53.81)(-.11779,.03245){8}{\line(-1,0){.11779}}
\multiput(63.21,54.33)(-.11779,.03245){8}{\line(-1,0){.11779}}
\multiput(61.32,54.85)(-.11779,.03245){8}{\line(-1,0){.11779}}
\multiput(59.44,55.37)(-.11779,.03245){8}{\line(-1,0){.11779}}
\multiput(57.55,55.89)(-.11779,.03245){8}{\line(-1,0){.11779}}
\multiput(55.67,56.41)(-.11779,.03245){8}{\line(-1,0){.11779}}
\multiput(53.78,56.93)(-.11779,.03245){8}{\line(-1,0){.11779}}
\multiput(51.9,57.45)(-.11779,.03245){8}{\line(-1,0){.11779}}
\multiput(50.01,57.97)(-.11779,.03245){8}{\line(-1,0){.11779}}
\multiput(48.13,58.49)(-.11779,.03245){8}{\line(-1,0){.11779}}
\multiput(46.25,59.01)(-.11779,.03245){8}{\line(-1,0){.11779}}
%\end
%\qbezvec[both](25,128.25)(27.13,129.13)(26.75,131.5)
\put(26.75,131.5){\vector(0,1){.07}}\put(25,128.25){\vector(-3,-2){.07}}\qbezier(25,128.25)(27.13,129.13)(26.75,131.5)
%\end
\put(28.75,16.75){\line(0,1){15.75}}
%\dashline{1}(28.75,53.5)(28.75,32.75)
\put(28.68,53.43){\line(0,-1){.988}}
\put(28.68,51.45){\line(0,-1){.988}}
\put(28.68,49.48){\line(0,-1){.988}}
\put(28.68,47.5){\line(0,-1){.988}}
\put(28.68,45.52){\line(0,-1){.988}}
\put(28.68,43.55){\line(0,-1){.988}}
\put(28.68,41.57){\line(0,-1){.988}}
\put(28.68,39.6){\line(0,-1){.988}}
\put(28.68,37.62){\line(0,-1){.988}}
\put(28.68,35.64){\line(0,-1){.988}}
\put(28.68,33.67){\line(0,-1){.988}}
%\end
%\dashline{1}(28,32.75)(28,32.75)
%\end
%\dashline{1}(69.75,28)(69.75,48.75)
\put(69.68,27.93){\line(0,1){.988}}
\put(69.68,29.91){\line(0,1){.988}}
\put(69.68,31.88){\line(0,1){.988}}
\put(69.68,33.86){\line(0,1){.988}}
\put(69.68,35.83){\line(0,1){.988}}
\put(69.68,37.81){\line(0,1){.988}}
\put(69.68,39.79){\line(0,1){.988}}
\put(69.68,41.76){\line(0,1){.988}}
\put(69.68,43.74){\line(0,1){.988}}
\put(69.68,45.72){\line(0,1){.988}}
\put(69.68,47.69){\line(0,1){.988}}
%\end
%\dashline{1}(69.75,48.75)(69.75,48.75)
%\end
%\dashline{1}(74.25,40.25)(59.75,58.25)
\multiput(74.18,40.18)(-.032222,.04){18}{\line(0,1){.04}}
\multiput(73.02,41.62)(-.032222,.04){18}{\line(0,1){.04}}
\multiput(71.86,43.06)(-.032222,.04){18}{\line(0,1){.04}}
\multiput(70.7,44.5)(-.032222,.04){18}{\line(0,1){.04}}
\multiput(69.54,45.94)(-.032222,.04){18}{\line(0,1){.04}}
\multiput(68.38,47.38)(-.032222,.04){18}{\line(0,1){.04}}
\multiput(67.22,48.82)(-.032222,.04){18}{\line(0,1){.04}}
\multiput(66.06,50.26)(-.032222,.04){18}{\line(0,1){.04}}
\multiput(64.9,51.7)(-.032222,.04){18}{\line(0,1){.04}}
\multiput(63.74,53.14)(-.032222,.04){18}{\line(0,1){.04}}
\multiput(62.58,54.58)(-.032222,.04){18}{\line(0,1){.04}}
\multiput(61.42,56.02)(-.032222,.04){18}{\line(0,1){.04}}
\multiput(60.26,57.46)(-.032222,.04){18}{\line(0,1){.04}}
%\end
\put(69.75,48.5){\line(0,1){16.25}}
%\emline(56,60.25)(69.75,64.5)
\multiput(56,60.25)(.10912698,.03373016){126}{\line(1,0){.10912698}}
%\end
\put(38.25,42.5){\line(0,-1){18.5}}
%\emline(38,24)(28.75,26.25)
\multiput(38,24)(-.1380597,.0335821){67}{\line(-1,0){.1380597}}
%\end
%\emline(25.75,44)(15.75,33)
\multiput(25.75,44)(-.033670034,-.037037037){297}{\line(0,-1){.037037037}}
%\end
%\emline(56.25,20.5)(48.75,11.75)
\multiput(56.25,20.5)(-.03363229,-.03923767){223}{\line(0,-1){.03923767}}
%\end
%\emline(60,58.25)(49.25,70.75)
\multiput(60,58.25)(-.03369906,.039184953){319}{\line(0,1){.039184953}}
%\end
%\emline(41,60.75)(49.25,70.75)
\multiput(41,60.75)(.033673469,.040816327){245}{\line(0,1){.040816327}}
%\end
%\emline(49.25,70.75)(44.5,59.25)
\multiput(49.25,70.75)(-.03368794,-.08156028){141}{\line(0,-1){.08156028}}
%\end
%\emline(44.71,59.7)(56.27,62.65)
\multiput(44.71,59.7)(.1313901,.0334447){88}{\line(1,0){.1313901}}
%\end
\put(59.46,58.8){\line(0,-1){20.81}}
%\emline(54.8,32.58)(69.73,49.82)
\multiput(54.8,32.58)(.033692802,.038912814){443}{\line(0,1){.038912814}}
%\end
%\emline(69.73,52.56)(81.92,49.61)
\multiput(69.73,52.56)(.1385568,-.0334447){88}{\line(1,0){.1385568}}
%\end
\put(81.92,49.61){\line(-1,0){12.19}}
%\emline(82.13,49.4)(71.62,37.63)
\multiput(82.13,49.4)(-.03368976,-.037732531){312}{\line(0,-1){.037732531}}
%\end
%\emline(69.73,49.61)(75.82,41.83)
\multiput(69.73,49.61)(.03368231,-.04297399){181}{\line(0,-1){.04297399}}
%\end
%\emline(59.65,58.69)(69.73,64.35)
\multiput(59.65,58.69)(.05997781,.03367175){168}{\line(1,0){.05997781}}
%\end
%\emline(59.65,58.69)(69.73,55.68)
\multiput(59.65,58.69)(.1119586,-.0333912){90}{\line(1,0){.1119586}}
%\end
%\emline(44.62,59.4)(59.3,42.25)
\multiput(44.62,59.4)(.033729806,-.039419171){435}{\line(0,-1){.039419171}}
%\end
%\emline(15.81,32.88)(28.71,28.46)
\multiput(15.81,32.88)(.09850915,-.03373601){131}{\line(1,0){.09850915}}
%\end
\put(15.63,32.7){\line(1,0){12.9}}
%\emline(28.54,32.7)(22,39.77)
\multiput(28.54,32.7)(-.03371514,.0364488){194}{\line(0,1){.0364488}}
%\end
%\emline(51.87,21.92)(38.26,37.3)
\multiput(51.87,21.92)(-.033692588,.038068249){404}{\line(0,1){.038068249}}
%\end
%\emline(32.25,54.98)(15.45,50.91)
\multiput(32.25,54.98)(-.13879162,-.03360218){121}{\line(-1,0){.13879162}}
%\end
%\emline(15.45,50.73)(23.23,41.19)
\multiput(15.45,50.73)(.03367175,-.04132442){231}{\line(0,-1){.04132442}}
%\end
%\emline(15.63,50.73)(25.88,44.02)
\multiput(15.63,50.73)(.05126524,-.03358757){200}{\line(1,0){.05126524}}
%\end
%\emline(25.88,44.02)(49.04,49.67)
\multiput(25.88,44.02)(.13784373,.03367175){168}{\line(1,0){.13784373}}
%\end
%\emline(28.71,32.53)(43.56,48.26)
\multiput(28.71,32.53)(.033671751,.035676022){441}{\line(0,1){.035676022}}
%\end
%\emline(32.6,54.98)(38.08,61.34)
\multiput(32.6,54.98)(.03362011,.03904271){163}{\line(0,1){.03904271}}
%\end
%\emline(32.45,54.95)(28.56,64.14)
\multiput(32.45,54.95)(-.03352661,.07924473){116}{\line(0,1){.07924473}}
%\end
%\emline(28.71,64.17)(44.62,59.57)
\multiput(28.71,64.17)(.11613068,-.03354886){137}{\line(1,0){.11613068}}
%\end
\put(25.88,44.19){\line(0,1){9.19}}
%\emline(38.41,24.01)(48.84,11.82)
\multiput(38.41,24.01)(.033644597,-.039347071){310}{\line(0,-1){.039347071}}
%\end
%\emline(65.36,26.49)(59.53,19.95)
\multiput(65.36,26.49)(-.03372041,-.03780773){173}{\line(0,-1){.03780773}}
%\end
%\dashline{1}(40.73,20.86)(65.48,26.52)
\multiput(40.66,20.79)(.13598,.03108){7}{\line(1,0){.13598}}
\multiput(42.57,21.22)(.13598,.03108){7}{\line(1,0){.13598}}
\multiput(44.47,21.66)(.13598,.03108){7}{\line(1,0){.13598}}
\multiput(46.37,22.09)(.13598,.03108){7}{\line(1,0){.13598}}
\multiput(48.28,22.53)(.13598,.03108){7}{\line(1,0){.13598}}
\multiput(50.18,22.97)(.13598,.03108){7}{\line(1,0){.13598}}
\multiput(52.09,23.4)(.13598,.03108){7}{\line(1,0){.13598}}
\multiput(53.99,23.84)(.13598,.03108){7}{\line(1,0){.13598}}
\multiput(55.89,24.27)(.13598,.03108){7}{\line(1,0){.13598}}
\multiput(57.8,24.71)(.13598,.03108){7}{\line(1,0){.13598}}
\multiput(59.7,25.14)(.13598,.03108){7}{\line(1,0){.13598}}
\multiput(61.6,25.58)(.13598,.03108){7}{\line(1,0){.13598}}
\multiput(63.51,26.01)(.13598,.03108){7}{\line(1,0){.13598}}
%\end
%\dashline{1}(65.48,26.52)(65.48,26.52)
%\end
%\emline(28.71,16.62)(38.26,23.86)
\multiput(28.71,16.62)(.04439973,.0337109){215}{\line(1,0){.04439973}}
%\end
%\emline(28.71,16.79)(40.91,20.68)
\multiput(28.71,16.79)(.10515165,.03352661){116}{\line(1,0){.10515165}}
%\end
\put(51.87,21.74){\line(-3,-1){9.02}}
%\emline(65.48,26.52)(81.75,31.11)
\multiput(65.48,26.52)(.11871136,.03354886){137}{\line(1,0){.11871136}}
%\end
%\emline(81.75,31.11)(74.14,40.31)
\multiput(81.75,31.11)(-.0336345,.04067428){226}{\line(0,1){.04067428}}
%\end
%\emline(71.57,37.61)(81.83,30.92)
\multiput(71.57,37.61)(.05154227,-.03361452){199}{\line(1,0){.05154227}}
%\end
\put(71.72,37.61){\line(0,-1){9.22}}
%\emline(65.62,26.46)(42.43,32.55)
\multiput(65.62,26.46)(-.128119,.0336723){181}{\line(-1,0){.128119}}
%\end
\put(22.7,106.5){\makebox(0,0)[cc]{$A_{4k}$}}
\put(94,105.59){\makebox(0,0)[cc]{$A_{4k}$}}
\put(21.5,134.34){\makebox(0,0)[cc]{$A_{4k}$}}
\put(36,89.84){\makebox(0,0)[cc]{$A_{4k+2}$}}
\put(23.25,119.84){\makebox(0,0)[cc]{$A_{4k+1}$}}
\put(39,119.84){\makebox(0,0)[cc]{$A_{4k+2}$}}
\put(84.25,94.84){\makebox(0,0)[cc]{$A_{4k+2}$}}
\put(36,96.09){\makebox(0,0)[cc]{$B$}}
\put(77,101.34){\makebox(0,0)[cc]{$B$}}
\put(32.5,120.84){\makebox(0,0)[cc]{$B$}}
\put(31.25,129.09){\makebox(0,0)[cc]{$\frac{2\pi}{10}$}}
\put(57.5,129.09){\makebox(0,0)[cc]{$\frac{\pi}{5}$}}
\put(54.5,126.34){\makebox(0,0)[cc]{$\frac{\pi}{5}$}}
\put(50.75,112.34){\makebox(0,0)[cc]{(b)}}
\put(15.75,75.84){\makebox(0,0)[cc]{(c)}}
\put(77.25,75.59){\makebox(0,0)[cc]{(d)}}
\put(34.75,114.84){\makebox(0,0)[cc]{Fold Line}}
\put(33.5,107){\makebox(0,0)[cc]{Fold Line}}
\put(48.43,3.34){\makebox(0,0)[cc]{Figure 5.3}}
\end{picture}
\end{center}

\begin{center}
%TeXCAD Picture [fig5.4.pic]. Options:
%\grade{\on}
%\emlines{\off}
%\epic{\off}
%\beziermacro{\on}
%\reduce{\on}
%\snapping{\off}
%\quality{8.00}
%\graddiff{0.01}
%\snapasp{1}
%\zoom{4.0000}
\unitlength 1mm % = 2.85pt
\linethickness{0.4pt}
\ifx\plotpoint\undefined\newsavebox{\plotpoint}\fi % GNUPLOT compatibility
\begin{picture}(62.25,87)(0,0)
\put(6,87){\line(1,0){52.25}}
\put(58.25,87){\line(0,1){0}}
\put(57.75,81.75){\line(-1,0){51.5}}
%\dottedline(8.75,81.75)(6,85.5)
\multiput(8.68,81.68)(-.55,.75){6}{{\rule{.4pt}{.4pt}}}
%\end
%\dottedline(6,85.5)(15.5,81.75)
\multiput(5.93,85.43)(.8636,-.3409){12}{{\rule{.4pt}{.4pt}}}
%\end
%\dottedline(15.5,81.75)(18,86.75)
\multiput(15.43,81.68)(.3571,.7143){8}{{\rule{.4pt}{.4pt}}}
%\end
%\dottedline(18,86.75)(22.75,82)
\multiput(17.93,86.68)(.5938,-.5938){9}{{\rule{.4pt}{.4pt}}}
%\end
%\dottedline(22.75,82)(31.75,86.75)
\multiput(22.68,81.93)(.8182,.4318){12}{{\rule{.4pt}{.4pt}}}
%\end
%\dottedline(31.75,86.75)(35.75,81.75)
\multiput(31.68,86.68)(.5,-.625){9}{{\rule{.4pt}{.4pt}}}
%\end
%\dottedline(35.75,81.75)(40,86.75)
\multiput(35.68,81.68)(.5313,.625){9}{{\rule{.4pt}{.4pt}}}
%\end
%\dottedline(40,86.75)(43.75,82)
\multiput(39.93,86.68)(.5357,-.6786){8}{{\rule{.4pt}{.4pt}}}
%\end
%\dottedline(40.5,86.75)(51.25,82)
\multiput(40.43,86.68)(.8958,-.3958){13}{{\rule{.4pt}{.4pt}}}
%\end
%\dottedline(51.25,82)(52.25,86.5)
\multiput(51.18,81.93)(.1667,.75){7}{{\rule{.4pt}{.4pt}}}
%\end
%\dottedline(52.25,86.5)(58.25,81.75)
\multiput(52.18,86.43)(.6667,-.5278){10}{{\rule{.4pt}{.4pt}}}
%\end
%\dottedline(23,82)(25.5,87)
\multiput(22.93,81.93)(.3571,.7143){8}{{\rule{.4pt}{.4pt}}}
%\end
\thicklines
%\emline(42.25,11.25)(43,22.75)
\multiput(42.25,11.25)(.032609,.5){23}{\line(0,1){.5}}
%\end
%\emline(61.5,52.25)(50,53)
\multiput(61.5,52.25)(-.5,.032609){23}{\line(-1,0){.5}}
%\end
\thinlines
%\emline(29.25,34.75)(49.5,26.75)
\multiput(29.25,34.75)(.08508403,-.03361345){238}{\line(1,0){.08508403}}
%\end
%\emline(38,39.25)(46,59.5)
\multiput(38,39.25)(.03361345,.08508403){238}{\line(0,1){.08508403}}
%\end
%\emline(29.5,34.5)(28.25,35)
\multiput(29.5,34.5)(-.083333,.033333){15}{\line(-1,0){.083333}}
%\end
%\emline(38.25,39.5)(37.75,38.25)
\multiput(38.25,39.5)(-.033333,-.083333){15}{\line(0,-1){.083333}}
%\end
\put(19,78){\makebox(0,0)[cc]{$D^{1} U^{1} D^{2}$ tape}}
%\emline(38,25)(52.5,19)
\multiput(38,25)(.08146067,-.03370787){178}{\line(1,0){.08146067}}
%\end
%\emline(47.75,48)(53.75,62.5)
\multiput(47.75,48)(.03370787,.08146067){178}{\line(0,1){.08146067}}
%\end
%\emline(52.5,19)(48,29.5)
\multiput(52.5,19)(-.03358209,.07835821){134}{\line(0,1){.07835821}}
%\end
%\emline(53.75,62.5)(43.25,58)
\multiput(53.75,62.5)(-.07835821,-.03358209){134}{\line(-1,0){.07835821}}
%\end
\put(49.75,25.5){\line(-1,-1){3.75}}
\put(47.25,59.75){\line(1,-1){3.75}}
\put(8.5,55.75){\line(1,1){3.75}}
\put(16,20.25){\line(-1,1){3.75}}
%\emline(59.5,28)(41,30.25)
\multiput(59.5,28)(-.2761194,.0335821){67}{\line(-1,0){.2761194}}
%\end
%\emline(44.75,69.5)(42.5,51)
\multiput(44.75,69.5)(-.0335821,-.2761194){67}{\line(0,-1){.2761194}}
%\end
%\emline(28.25,35.25)(55,33.75)
\multiput(28.25,35.25)(.5944444,-.0333333){45}{\line(1,0){.5944444}}
%\end
%\emline(59.5,28)(52.25,37)
\multiput(59.5,28)(-.03372093,.04186047){215}{\line(0,1){.04186047}}
%\end
%\emline(44.75,69.5)(35.75,62.25)
\multiput(44.75,69.5)(-.04186047,-.03372093){215}{\line(-1,0){.04186047}}
%\end
%\emline(56,32.75)(53.25,28.75)
\multiput(56,32.75)(-.0335366,-.0487805){82}{\line(0,-1){.0487805}}
%\end
%\emline(53.75,52.75)(34.25,35.25)
\multiput(53.75,52.75)(-.037572254,-.03371869){519}{\line(-1,0){.037572254}}
%\end
%\emline(24.5,27.25)(30.5,11.75)
\multiput(24.5,27.25)(.03370787,-.08707865){178}{\line(0,-1){.08707865}}
%\end
%\emline(46,34.75)(61.5,40.75)
\multiput(46,34.75)(.08707865,.03370787){178}{\line(1,0){.08707865}}
%\end
%\emline(30.5,11.75)(35.25,20)
\multiput(30.5,11.75)(.03368794,.05851064){141}{\line(0,1){.05851064}}
%\end
%\emline(61.5,40.75)(53.25,45.5)
\multiput(61.5,40.75)(-.05851064,.03368794){141}{\line(-1,0){.05851064}}
%\end
%\emline(56,44)(37,35)
\multiput(56,44)(-.071161049,-.033707865){267}{\line(-1,0){.071161049}}
%\end
%\emline(28.75,16.75)(33,16)
\multiput(28.75,16.75)(.184783,-.032609){23}{\line(1,0){.184783}}
%\end
%\emline(56.5,39)(57.25,43.25)
\multiput(56.5,39)(.032609,.184783){23}{\line(0,1){.184783}}
%\end
%\emline(29,27.25)(42.25,11)
\multiput(29,27.25)(.033715013,-.041348601){393}{\line(0,-1){.041348601}}
%\end
%\emline(46,39.25)(62.25,52.5)
\multiput(46,39.25)(.041348601,.033715013){393}{\line(1,0){.041348601}}
%\end
%\emline(34.75,71.25)(38,50.25)
\multiput(34.75,71.25)(.03350515,-.21649485){97}{\line(0,-1){.21649485}}
%\end
%\emline(34.75,71.25)(31.5,60.25)
\multiput(34.75,71.25)(-.03350515,-.11340206){97}{\line(0,-1){.11340206}}
%\end
%\emline(32.75,65)(35.75,65.25)
\multiput(32.75,65)(.375,.03125){8}{\line(1,0){.375}}
%\end
\put(38.5,64.25){\line(0,-1){22.5}}
%\emline(38.5,41.5)(32,62.25)
\multiput(38.5,41.5)(-.03367876,.10751295){193}{\line(0,1){.10751295}}
%\end
%\emline(34.25,55.25)(26.75,69.75)
\multiput(34.25,55.25)(-.03363229,.06502242){223}{\line(0,1){.06502242}}
%\end
%\emline(26.5,69.75)(25.25,58.75)
\multiput(26.5,69.75)(-.0328947,-.2894737){38}{\line(0,-1){.2894737}}
%\end
%\emline(25.5,62)(37.25,45.25)
\multiput(25.5,62)(.033667622,-.047994269){349}{\line(0,-1){.047994269}}
%\end
%\emline(31.25,54)(14.75,68.5)
\multiput(31.25,54)(-.038372093,.03372093){430}{\line(-1,0){.038372093}}
%\end
%\emline(14.75,68.5)(15.5,58.25)
\multiput(14.75,68.5)(.032609,-.445652){23}{\line(0,-1){.445652}}
%\end
%\emline(15.5,61.5)(35.75,47.5)
\multiput(15.5,61.5)(.048795181,-.03373494){415}{\line(1,0){.048795181}}
%\end
%\emline(31.5,50.25)(9.5,54)
\multiput(31.5,50.25)(-.19642857,.03348214){112}{\line(-1,0){.19642857}}
%\end
%\emline(23.75,56)(4,61)
\multiput(23.75,56)(-.13255034,.03355705){149}{\line(-1,0){.13255034}}
%\end
%\emline(4,61)(11.25,51.5)
\multiput(4,61)(.03372093,-.04418605){215}{\line(0,-1){.04418605}}
%\end
%\emline(1.25,50.25)(12,43.25)
\multiput(1.25,50.25)(.05168269,-.03365385){208}{\line(1,0){.05168269}}
%\end
%\emline(9.25,45.25)(30,50.5)
\multiput(9.25,45.25)(.13301282,.03365385){156}{\line(1,0){.13301282}}
%\end
%\emline(1.5,50.25)(18.25,52.75)
\multiput(1.5,50.25)(.2233333,.0333333){75}{\line(1,0){.2233333}}
%\end
%\emline(28,50)(7.5,35.75)
\multiput(28,50)(-.048463357,-.033687943){423}{\line(-1,0){.048463357}}
%\end
%\emline(14.75,46.75)(2,35.75)
\multiput(14.75,46.75)(-.038990826,-.033639144){327}{\line(-1,0){.038990826}}
%\end
%\emline(2,35.75)(10.5,35.5)
\multiput(2,35.75)(1.0625,-.03125){8}{\line(1,0){1.0625}}
%\end
%\emline(25.5,48)(10.75,27.5)
\multiput(25.5,48)(-.033675799,-.046803653){438}{\line(0,-1){.046803653}}
%\end
%\emline(14.25,40.5)(2.75,26.5)
\multiput(14.25,40.5)(-.03372434,-.041055718){341}{\line(0,-1){.041055718}}
%\end
%\emline(2.5,26.5)(14,27.75)
\multiput(2.5,26.5)(.3026316,.0328947){38}{\line(1,0){.3026316}}
%\end
%\emline(24.5,46.5)(18,21.5)
\multiput(24.5,46.5)(-.03367876,-.12953368){193}{\line(0,-1){.12953368}}
%\end
%\emline(16.25,35.25)(10.5,16.5)
\multiput(16.25,35.25)(-.03362573,-.10964912){171}{\line(0,-1){.10964912}}
%\end
%\emline(10.5,16.5)(20,23)
\multiput(10.5,16.5)(.0492228,.03367876){193}{\line(1,0){.0492228}}
%\end
%\emline(24,44.5)(24.5,19.25)
\multiput(24,44.5)(.033333,-1.683333){15}{\line(0,-1){1.683333}}
%\end
%\emline(19.75,29.25)(20,13.75)
\multiput(19.75,29.25)(.03125,-1.9375){8}{\line(0,-1){1.9375}}
%\end
%\emline(20,13.75)(26.75,22)
\multiput(20,13.75)(.03358209,.04104478){201}{\line(0,1){.04104478}}
%\end
%\emline(24,39.75)(42.75,19)
\multiput(24,39.75)(.033723022,-.037320144){556}{\line(0,-1){.037320144}}
%\end
%\emline(24.25,35)(33.75,17.75)
\multiput(24.25,35)(.033687943,-.061170213){282}{\line(0,-1){.061170213}}
%\end
\put(26,65){\line(1,0){3.25}}
%\emline(15,63.5)(18.75,65.5)
\multiput(15,63.5)(.0625,.0333333){60}{\line(1,0){.0625}}
%\end
%\emline(6.5,47)(6.75,51.25)
\multiput(6.5,47)(.03125,.53125){8}{\line(0,1){.53125}}
%\end
%\emline(6,36)(5.5,39)
\multiput(6,36)(-.033333,.2){15}{\line(0,1){.2}}
%\end
%\emline(8.5,27.5)(6.25,30.75)
\multiput(8.5,27.5)(-.0335821,.0485075){67}{\line(0,1){.0485075}}
%\end
%\emline(40,65.5)(43.75,62.75)
\multiput(40,65.5)(.0457317,-.0335366){82}{\line(1,0){.0457317}}
%\end
%\emline(20,19.75)(23,17.75)
\multiput(20,19.75)(.05,-.0333333){60}{\line(1,0){.05}}
%\end
%\emline(38.25,16.25)(42.75,17.25)
\multiput(38.25,16.25)(.15,.033333){30}{\line(1,0){.15}}
%\end
%\emline(56.75,48)(56,52.5)
\multiput(56.75,48)(-.032609,.195652){23}{\line(0,1){.195652}}
%\end
\put(30.75,3.75){\makebox(0,0)[cc]{Figure 5.4}}
\end{picture}
\end{center}

\begin{description}

  \item[Remark:] We have been concerned to give an algorithm; that is, a systematic procedure for folding regular $\left\{ \frac{b}{a} \right\}$-gons. Of course, we do not advocate this algorithm as the best procedure in any particular case. For example, if $a < \left( \frac{1}{2} \right) b'$, it is plainly best to take the tape prepared for folding the $\left\{ \frac{b'}{a} \right\}$-gon, on which there will already be a crease line making an angle of $\left( \frac{a}{b'} \right) \pi$ with the forward direction of the top of the tape, and then to bisect this angle by folding $k$ times to produce the required angle of $\left( \frac{a}{b} \right) \pi$. This remark may seem always applicable if $a = 1$, that is when regular \emph{convex} polygons are in question; however, with $a = 1$, this simplified procedure actually coincides with our algorithm!

\end{description}



\newpage

\section*{Appendix - The Further Generation \\ Of Folding Numbers}
\addcontentsline{toc}{section}{\numberline{}Appendix - The Further Generation Of Folding Numbers}
\markboth{Appendix - Further Generation}{Appendix - Further Generation}


In Section 3, we described the set of \emph{$t$-folding numbers}, $F_{t}$, consisting of integers $s$ of the form
\begin{equation}
s = \frac{t^{xy}-1}{t^{x}-1}, \quad x \geq 1, \quad y \geq 2  \tag{A1}
\end{equation}

We report here the main results of [10], in which we generalize even further, to consider the set $F_{tu}$ of \emph{$(t,u)$-folding numbers} $s$, being integers of the form
\begin{equation}
s = \frac{t^{xy}-u^{xy}}{t^{x}-u^{x}}, \quad x \geq 1, \quad y \geq 2  \tag{A2}
\end{equation}

In (A2), $t,u$ are fixed positive integers with $t > u \geq 1$ and $\gcd(t,u) = 1$. We are content in this appendix simply to report the results; the proofs will be found in [10]. In each case, we indicate the result of Section 3 which is being generalized.

\begin{description}

  \item[Theorem A1:] If $\gcd(a,b) = d$, then $\gcd (t^{a}-u^{a}, t^{b}-u^{b}) = (t^{d}-u^{d})$. \emph{(Theorem 3.1)}

\end{description}

Now let $c$ be a positive integer $\geq 2$, prime to $t$ and $u$; and let $\frac{F}{c}$ be the set defined by the rule
\begin{equation*}
\frac{F}{c} = \{ s \in F_{tu} \text{ such that } c\ |\ s \}
\end{equation*}

Then $\frac{F}{c}$ is non-empty; indeed, we may fix $x$ in (A2) and the set of $s$ of the form (A2) with a given $x$ such that $c\ |\ s$ is non-empty.

We fix $x=1$ in (A2), let $y_{0}$ be minimal such that $c\ \left|\ \frac{(t^{y_{0}}-u^{t_{0}})}{(t-u)} \right.$, and write $y_{0} = h(c)$, the \emph{height} of $c$. Then

\begin{description}

  \item[Lemma A2:] If $c\ \left|\ \frac{(t^{xy}-u^{xy})}{(t^{x}-u^{x})} \right.$, then $h(c)\ |\ xy$. \emph{(Theorem 3.4 (ii) )}

  \item[Theorem A3:] Let $s_{0}$ be the smallest integer in $\frac{F}{c}$, where $s_{0} = \frac{(t^{x_{0}y_{0}}-u^{x_{0}y_{0}})}{(t^{x_{0}}-u^{x_{0}})}$. Then $h(c) = x_{0} y_{0}$. \emph{(Theorem 3.5)}

  \item[Corollary A4:] If there exist $(t,u)$ such that $\frac{(t^{xy}-u^{xy})}{(t^{x}-u^{x})}\ \left|\ \frac{(t^{x'y'}-u^{x'y'})}{(t^{x'}-u^{y'})} \right.$, then $xy\ |\ x'y'$. \emph{(Corollary 3.6)}

\end{description}

The main theorem of [10] is the following. Suppose $s,s'$ are the two $(t,u)$-folding numbers
\begin{equation}
s = \frac{t^{xy}-u^{xy}}{t^{x}-u^{x}}, \quad s' = \frac{t^{x'y'}-u^{x'y'}}{t^{x'}-u^{x'}}  \tag{A4}
\end{equation}
and suppose that $xy\ |\ x' y'$. Let $d = \gcd (xy, x')$, $h = \gcd (x,x')$, and set $xy = ad$, $x' = bd$. Then $ad\ |\ bdy'$, so that $a\ |\ by'$, hence $a\ |\ y'$.

\begin{description}

  \item[Theorem A5:] $s\ \left|\ s'\ \Longleftrightarrow\ \frac{(t^{d}-u^{d})}{(t^{h}-u^{h})} \right|\ \frac{y'}{a}$. \emph{(Theorem 3.7)}

\end{description}

This theorem provides an effective algorithm for deciding for what values of $(t,u)$ we have $s\ |\ s'$. For there are 2 cases to be considered:
\begin{enumerate}[\hspace{1cm}(1)]
  \item $\gcd (xy, x')\ |\ x$, i.e., $d = h$. Then $s\ |\ s'$ for all pairs $(t,u)$.
  \item $\gcd (xy, x') \not|\ x$, i.e., $h\ |\ d$, $h \neq d$. Then $\frac{(t^{d}-u^{d})}{(t^{h}-u^{h})} < \frac{y'}{a}$ for only finitely many pairs $(t,u)$, so that we may test for the required divisibility in a finite number of steps.
\end{enumerate}

An easy computer algorithm answers our question quickly. With regard to Case 1, we have the following important refinement.

\begin{description}

  \item[Theorem A6:] For $s,s'$ as in (A4), the following statements are equivalent:
    \begin{enumerate}[\hspace{1cm}(i)]
      \item $s\ |\ s'$ for all pairs $(t,u)$;
      \item $s\ |\ s'$ for infinitely many pairs $(t,u)$;
      \item $s\ |\ s'$ as polynomials in $Z[t]$ (or $Z[t,u]$);
      \item $xy\ |\ x'y'$ and $\gcd(xy,x')\ |\ x$ 
    \end{enumerate}
    \emph{(Theorem 3.13)}

  \item[Example:] We may now extend an example given in Section 3. First, we find that the relation $7\ |\ 21$ is a special case of the general fact that
    \begin{equation*}
    \frac{t^{3}-u^{3}}{t-u}\ \left|\ \frac{t^{6}-u^{6}}{t^{2}-u^{2}} \right. \text{ for all pairs } (t,u)
    \end{equation*}

    Second, however, we find that the relation $3\ |\ 21$ generalizes to the statement $\frac{(t^{3}-u^{3})}{(t-u)}\ \left|\ \frac{(t^{6}-u^{6})}{(t^{2}-u^{2})} \right.$, which is \emph{false} for \emph{all} pairs $(t,u)$, with the single exception of $t = 2$, $u = 1$. No two generalizations could possibly have more divergent truth values!

\end{description}

More sophisticated, but very striking, is the following divisibility result involving integers in $F_{tu}$. Given any positive integer $n$, let $\psi(n)$ be the exponent of the multiplicative group $\left( \frac{Z}{n} \right) *$ of residues prime to $n$ (or units of the ring $\frac{Z}{n}$). Consider the set $S_{q}$ of all integers $s$ of the form $s = \frac{(t^{xy}-u^{xy})}{(t^{x}-u^{x})}$, such that $xy = q$.

\begin{description}

  \item[Theorem A7:] If $s \in S_{q}$, then $q\ |\ \psi(s)$. \emph{(Theorem 3.15 (ii) and the final note of Section 3)}

\end{description}

We describe in [10] how to calculate $\psi(s)$. For example, let $q=6$, $t=3$, $u=2$. Then
\begin{equation*}
s = \frac{3^{6}-2^{6}}{3^{2}-2^{2}} = 133 \in S_{6}
\end{equation*}

Now $133 = 7 \times 19$, and $\Phi(7) = 6$, $\Phi(19) = 18$. Thus $\psi(133) = \text{lcm}(6,18) = 18$, and $6\ |\ 18$. We may also take $s = \frac{(3^{6}-2^{6})}{(3^{3}-2^{3})} = 35 \in S_{6}$, and $\Phi(5) = 4$, $\Phi(7) = 6$. Thus $\psi(35) = \text{lcm}(4,6) = 12$, and $6\ |\ 12$. These calculations show how fine Theorem A7 is, for $\gcd(18,12) = 6$.




\newpage

\section*{Postscript - Some Questions For Diligent Readers}
\addcontentsline{toc}{section}{\numberline{}Postscript - Some Questions For Diligent Readers}
\markboth{Postscript - Questions}{Postscript - Questions}


\begin{description}

  \item[2.1:] What configurations can be obtained by taking tape folded by our systematic procedure and then applying some rule other than the FAT-algorithm? As special cases, consider the systematic procedures $D^{n} U^{n}$, $D^{m} U^{n} (m \neq n)$, $D^{1} U^{1} D^{3}$.

  \item[2.2:] Suppose that regular convex $(2^{n}+1)$-gons are obtained from $D^{n} U^{n}$-folded tape in various ways. How are these polygons related geometrically? (We know that there are at least $(n+1)$ different types of regular $(2n+1)$-gon obtainable from this tape.)

  \item[2.3:] Generalize Lemma 2.1, especially in ways relevant to our folding procedures. Here is one suggestion:

    Let $\{ x_{ik} \}$, $i = 1,2,\cdots,r$, $k = 0,1,2,\cdots$, be $r$ sequences of real numbers such that
    \begin{align*}
    x_{1k} & + a_{1} x_{2k} = b_{1}, \\
    x_{2k} & + a_{2} x_{3k} = b_{2}, \\
    \vdots & \\
    x_{r-1,k} & + a_{r-1} x_{rk} = b_{r-1}, \\
    x_{rk} & + a_{r} x_{1,k+1} = b_{r}
    \end{align*}
    for all $k$. Find conditions under which all the sequences converge. If the sequences converge, find their limits.

  \item[3.1:] Find a proof of Theorem 3.1 based on the fact that $\gcd(u,v)$ is a linear combination of $u$ and $v$ with integer coefficients.

  \item[3.2:] Let $a$ be prime to $t$, and let $h(a)$ be the $t$-height of $a$. We saw that if $a$ is prime, the smallest $s$ in $F_{t}$ with $a\ |\ s$ is given by
    \begin{equation*}
    s = \frac{t^{h(a)}-1}{t^{x_{1}}-1}
    \end{equation*}
    where $x_{1}$ is the largest proper factor of $h(a)$. Generalize this result by weakening the hypothesis that $a$ is prime.

  \item[3.3:] How many regular star polygons with $N$ vertices are there? If $N$ is a folding number, how many of these may be obtained from tape folded to produce the regular convex $N$-gon?

  \item[4.1:] Prove that (4.3) yields the unique solution to equation (4.2).

  \item[4.2:] Let $S$ be any set. Show that the sequence $(k_{1}, k_{2}, \cdots, k_{r-1})$, where $k_{i} \in S$, may be extended in at most one way to a repeating sequence $(k_{1}, k_{2}, \cdots,$ $k_{r-1}, k_{r})$, with $k_{r} \in S$. Here the sequence $(k_{1}, k_{2}, \cdots, k_{r})$ is said to \emph{repeat} if there exists a proper divisor $s$ of $r$ such that $k_{j} = k_{s+j}$ for all $j$. What is the relevance of this to the content of Section 4?

  \item[4.3:] Give an algorithm for determining what star polygons are available from tape folded according to the sequence of instructions $(k_{1}, k_{2}, \cdots, k_{r})$. Here of course, the $k_{i}$ are positive integers, and you may assume that the sequence is non-repeating.

  \item[5.1:] Describe explicitly how to fold a regular $\left\{ \frac{60}{7} \right\}$-gon by the algorithm of Section 5.

  \item[5.2:] Give a formula for the number of secondary fold-lines to be put into our (already folded) tape to produce a regular star $\left\{ \frac{b}{a} \right\}$-gon, where $b$ is even.

  \item[A.1:] Prove Theorem A1. (Hint: Look at Question 3.1.) What can be said if $\gcd(t,u) \neq 1$?

  \item[A.2:] State and prove the generalization of Lemma 3.10 needed in the proof of Theorem A5.

  \item[A.3:] Let $\gcd(t,u) = 1$. Define the equalizer of $t,u \mod{b}$ (where $\gcd(t,b) = \gcd(u,b) = 1$) to be the smallest positive integer $n$ such that $t^{n} \equiv u^{n} \mod{b}$. Define the \emph{quasi-equalizer} of $t,u \mod{b}$ similarly. Give an algorithm for determining the quasi-equalizer of $t,u \mod{b}$, and for deciding whether $t^{k} \equiv u^{k} \mod{b}$ or $t^{k} \equiv -u^{k} \mod{b}$, where $k$ is the equalizer.

    (Hint: Base this on the algorithm of Section 4 for determining quasi-orders $\mod{b}$.)

\end{description}




\newpage

\section*{References}
\addcontentsline{toc}{section}{\numberline{}References}
\markboth{References}{References}


\begin{enumerate}[1.]
  \item H. S. M. Coxeter. \emph{Regular Polytopes}. Methuen (1948).
  \item Jon Froemke and Jerrold W. Grossman. ``An algebraic approach to some number-theoretic problems arising from paper-folding regular polygons''. $\ $ \emph{American Mathematical Monthly}, 95 (1988), p289-307.
  \item Peter Hilton and Jean Pedersen. ``Approximating any regular polygon by folding paper; An interplay of geometry, analysis and number theory''. \emph{Math. Mag.}, 56 (1983), p141-155.
  \item Peter Hilton and Jean Pedersen. ``Folding regular star polygons and number theory''. \emph{Math Intelligencer}, 7, No. 1 (1985), p15-26.
  \item Peter Hilton and Jean Pedersen. ``Certain algorithms in the practice of geometry and the theory of numbers''. \emph{Publ. Sec. Mat. Univ. Autonoma Barcelona}, 29, No. 1 (1985), p31-64.
  \item Peter Hilton and Jean Pedersen. ``On generalized symbols, orders and quasi-orders''. \emph{Publ. Sec. Mat. Univ. Autonoma Barcelona}, 29, No. 2-3 (1985), p123-144.
  \item Peter Hilton and Jean Pedersen. ``The general quasi-order algorithm in number theory''. \emph{Int. Journ. Math. and Math. Sci.}, 9, No. 2 (1986) p245-252.
  \item Peter Hilton and Jean Pedersen. ``On the complementary factor in a new congruence algorithm''. \emph{Int. Journ. Math. and Math. Sci.}, 10, No. 1 (1987), p113-123.
  \item Peter Hilton and Jean Pedersen. \emph{Build Your Own Polyhedra}. Addison-Wesley (1988).
  \item Peter Hilton and Jean Pedersen. ``On a generalization of folding numbers''. \emph{Southeast Asian Bulletin of Mathematics}, 12, No. 1 (1988), p53-63.
  \item Carroll O. Wilde. \emph{The Contraction Mapping Principle}. UMAP Unit 326 (Lexington, MA: COMAP, Inc).
\end{enumerate}


\end{document}

